% !TEX program = pdflatex
% =====================================================================
%  GUÍA TÉCNICA DE HERRAMIENTAS DE LÍNEA DE COMANDOS
%  Documentos, imágenes y PDF en Ubuntu 26.04 LTS sobre WSL2 (Windows 11)
%
%  Herramientas: Marker · Poppler · img2pdf · ImageMagick · Potrace ·
%                Pandoc · QPDF · OCRmyPDF · Ghostscript
% =====================================================================
\documentclass[11pt,a4paper]{article}

% ---------------------------------------------------------------------
% Codificación, idioma y tipografía
% ---------------------------------------------------------------------
\usepackage[T1]{fontenc}
\usepackage[utf8]{inputenc}
\usepackage{lmodern}
\usepackage[spanish,es-noshorthands,es-tabla,es-nodecimaldot,es-nolists,es-nosectiondot]{babel}
\usepackage[a4paper,top=2.6cm,bottom=2.5cm,left=2.3cm,right=2.3cm,
headheight=15pt,headsep=0.75cm,footskip=1.1cm]{geometry}
\usepackage[table]{xcolor}
\usepackage{amsmath,amssymb}
\usepackage{pifont}
\usepackage{microtype}

% ---------------------------------------------------------------------
% Tablas, listas, figuras, encabezados y títulos
% ---------------------------------------------------------------------
\usepackage{array}
\usepackage{booktabs}
\usepackage{longtable}
\usepackage{tabularx}
\usepackage{enumitem}
\usepackage{graphicx}
\usepackage{fancyhdr}
\usepackage{titlesec}

% ---------------------------------------------------------------------
% Código, cajas y diagramas
% ---------------------------------------------------------------------
\usepackage{listings}
\usepackage{tcolorbox}
\tcbuselibrary{skins,breakable,raster}
\usepackage{tikz}
\usetikzlibrary{arrows.meta,chains,positioning}

% ---------------------------------------------------------------------
% URL e hipervínculos
% ---------------------------------------------------------------------
\usepackage{xurl}
\usepackage{hyperref}
\usepackage{bookmark}

% =====================================================================
% Colores
% =====================================================================
\definecolor{acento}{HTML}{1F4E79}
\definecolor{acentoclaro}{HTML}{E3ECF6}
\definecolor{codigofondo}{HTML}{F5F7FA}
\definecolor{codigoborde}{HTML}{C3CCD7}
\definecolor{codigocoment}{HTML}{5E6A78}
\definecolor{codigoclave}{HTML}{0B4F8A}
\definecolor{codigocadena}{HTML}{8A3B12}
\definecolor{colcmd}{HTML}{0B4F8A}
\definecolor{colopt}{HTML}{7A1F3D}
\definecolor{colarch}{HTML}{1E6B3A}
\definecolor{colnota}{HTML}{2B6CB0}
\definecolor{colconsejo}{HTML}{2F7D4F}
\definecolor{coladvert}{HTML}{B7791F}
\definecolor{colimport}{HTML}{B83232}
\definecolor{collimit}{HTML}{5A6270}
\definecolor{colejemplo}{HTML}{5B4B9A}
\definecolor{tablacab}{HTML}{E6EDF5}

\hypersetup{
	colorlinks=true,
	linkcolor=acento,
	urlcolor=colnota!85!black,
	citecolor=acento,
	bookmarksnumbered=true,
	bookmarksopen=true,
	pdfstartview=FitH,
	pdftitle={Guía técnica de herramientas de línea de comandos (Ubuntu 26.04 + WSL2)},
	pdfsubject={Marker, Poppler, img2pdf, ImageMagick, Potrace, Pandoc, QPDF, OCRmyPDF, Ghostscript},
	pdfkeywords={WSL2, Ubuntu, PDF, OCR, Markdown, LaTeX, Bash}
}

% =====================================================================
% Ajustes generales de composición
% =====================================================================
\setlength{\parindent}{0pt}
\setlength{\parskip}{0.55em plus 0.15em minus 0.1em}
\setlength{\emergencystretch}{3em}
\renewcommand{\arraystretch}{1.25}
\setlength{\LTpre}{6pt}
\setlength{\LTpost}{8pt}
\setlist{itemsep=0.25em, topsep=0.35em, parsep=0.15em, leftmargin=1.6em}
\newcolumntype{L}[1]{>{\raggedright\arraybackslash}p{#1}}
\newcolumntype{C}[1]{>{\centering\arraybackslash}p{#1}}
\newcolumntype{Y}{>{\raggedright\arraybackslash}X}

% Títulos de sección
\setcounter{secnumdepth}{3}
\setcounter{tocdepth}{2}
\titleformat{\section}
{\normalfont\LARGE\bfseries\color{acento}}{\thesection}{0.6em}{}
[\vspace{3pt}{\color{acento}\titlerule[0.8pt]}]
\titleformat{\subsection}
{\normalfont\Large\bfseries\color{acento}}{\thesubsection}{0.6em}{}
\titleformat{\subsubsection}
{\normalfont\large\bfseries\color{acento!85!black}}{\thesubsubsection}{0.6em}{}
\titleformat{\paragraph}[runin]
{\normalfont\normalsize\bfseries\color{acento!85!black}}{}{0pt}{}[.]
\titlespacing*{\section}{0pt}{0pt}{14pt}
\titlespacing*{\subsection}{0pt}{18pt}{8pt}
\titlespacing*{\subsubsection}{0pt}{14pt}{6pt}
\titlespacing*{\paragraph}{0pt}{8pt}{0.6em}


% Índice
\makeatletter
\renewcommand*\l@section[2]{%
	\ifnum \c@tocdepth >\z@
	\addpenalty\@secpenalty
	\addvspace{1.0em \@plus\p@}%
	\setlength\@tempdima{2.0em}%
	\begingroup
	\parindent \z@ \rightskip \@pnumwidth
	\parfillskip -\@pnumwidth
	\leavevmode \bfseries
	\advance\leftskip\@tempdima
	\hskip -\leftskip
	#1\nobreak\hfil
	\nobreak\hb@xt@\@pnumwidth{\hss #2\kern-\p@\kern\p@}\par
	\endgroup
	\fi}
\renewcommand*\l@subsection{\@dottedtocline{2}{2.0em}{3.4em}}
% (4)(5) Reservar espacio: si no cabe, pasa a la página siguiente
\newcommand{\necesitaespacio}[1]{\par\penalty-100\begingroup
	\setlength{\dimen@}{#1}\dimen@ii\pagegoal\advance\dimen@ii-\pagetotal
	\ifdim\dimen@>\dimen@ii\ifdim\dimen@ii>\z@\vfil\fi\penalty-\@M\fi\endgroup}
\makeatother

% Encabezados y pies de página
\pagestyle{fancy}
\fancyhf{}
\renewcommand{\sectionmark}[1]{\markboth{\thesection\quad #1}{}}
\renewcommand{\subsectionmark}[1]{}
\fancyhead[L]{\small\color{acento!85!black}\nouppercase{\leftmark}}
\fancyhead[R]{\small\color{acento!85!black}Herramientas CLI \textperiodcentered{} Ubuntu 26.04 + WSL2}
\fancyfoot[C]{\small\thepage}
\renewcommand{\headrulewidth}{0.4pt}
\fancypagestyle{plain}{%
	\fancyhf{}\fancyfoot[C]{\small\thepage}\renewcommand{\headrulewidth}{0pt}}

% =====================================================================
% Bloques de código (listings)
% =====================================================================
\lstdefinestyle{base}{
	basicstyle=\ttfamily\small,
	columns=fullflexible,
	keepspaces=true,
	showstringspaces=false,
	upquote=true,
	breaklines=true,
	breakatwhitespace=true,
	postbreak=\mbox{\textcolor{codigocoment}{$\hookrightarrow$}\space},
	tabsize=4,
	aboveskip=8pt,
	belowskip=8pt,
	xleftmargin=6pt,
	xrightmargin=6pt,
	framesep=5pt,
	framerule=0.4pt,
	extendedchars=true,
	literate=
	{á}{{\'a}}1 {é}{{\'e}}1 {í}{{\'i}}1 {ó}{{\'o}}1 {ú}{{\'u}}1
	{Á}{{\'A}}1 {É}{{\'E}}1 {Í}{{\'I}}1 {Ó}{{\'O}}1 {Ú}{{\'U}}1
	{ñ}{{\~n}}1 {Ñ}{{\~N}}1 {ü}{{\"u}}1 {Ü}{{\"U}}1
	{¿}{{\textquestiondown}}1 {¡}{{\textexclamdown}}1
	{«}{{\guillemotleft}}1 {»}{{\guillemotright}}1
	{→}{{$\rightarrow$}}2 {←}{{$\leftarrow$}}2 {↓}{{$\downarrow$}}1
	{—}{{\textemdash}}1 {–}{{\textendash}}1 {·}{{\textperiodcentered}}1
	{º}{{\textordmasculine}}1 {ª}{{\textordfeminine}}1
	{\%}{{\char37}}1
}
\lstdefinestyle{bash}{
	style=base,
	language=bash,
	keywordstyle=\color{codigoclave}\bfseries,
	commentstyle=\color{codigocoment}\itshape,
	stringstyle=\color{codigocadena},
	backgroundcolor=\color{codigofondo},
	frame=single,
	rulecolor=\color{codigoborde},
}
\lstdefinestyle{salida}{
	style=base,
	language={},
	basicstyle=\ttfamily\footnotesize,
	backgroundcolor=\color{white},
	frame=single,
	rulecolor=\color{codigoborde!70!white},
}
\lstdefinestyle{texto}{
	style=base,
	language={},
	backgroundcolor=\color{white},
	frame=none,
}
\lstdefinestyle{sintaxis}{
	style=base,
	language={},
	basicstyle=\ttfamily\small\color{acento!90!black},
	backgroundcolor=\color{acentoclaro!45!white},
	frame=leftline,
	framerule=2pt,
	rulecolor=\color{acento},
}
\lstset{style=bash}

% =====================================================================
% Entornos de tabla reutilizables (longtable: se dividen entre páginas)
% =====================================================================
% Tabla de opciones: \begin{opciones}[Encabezado col. 1] ... \end{opciones}
\newenvironment{opciones}[1][Opción]{%
	\par\begingroup\small
	\begin{longtable}{L{4.25cm}L{11.2cm}}
		\toprule\rowcolor{tablacab}\cab{#1} & \cab{Significado, valores y uso} \\ \midrule
		\endfirsthead
		\toprule\rowcolor{tablacab}\cab{#1} & \cab{Significado, valores y uso (continuación)} \\ \midrule
		\endhead
		\bottomrule
		\endfoot}
	{\end{longtable}\endgroup}
% Tabla de errores frecuentes
\newenvironment{errores}{%
	\par\begingroup\small
	\begin{longtable}{L{4.1cm}L{4.9cm}L{6.0cm}}
		\toprule\rowcolor{tablacab}\cab{Problema} & \cab{Causa probable} & \cab{Solución} \\ \midrule
		\endfirsthead
		\toprule\rowcolor{tablacab}\cab{Problema} & \cab{Causa probable} & \cab{Solución} \\ \midrule
		\endhead
		\bottomrule
		\endfoot}
	{\end{longtable}\endgroup}
% Tabla de formatos de una herramienta
\newenvironment{formatos}{%
	\par\begingroup\small
	\begin{longtable}{L{2.9cm}C{1.5cm}C{1.5cm}L{8.6cm}}
		\toprule\rowcolor{tablacab}\cab{Formato} & \cab{Entrada} & \cab{Salida} & \cab{Condiciones y observaciones} \\ \midrule
		\endfirsthead
		\toprule\rowcolor{tablacab}\cab{Formato} & \cab{Entrada} & \cab{Salida} & \cab{Condiciones y observaciones} \\ \midrule
		\endhead
		\bottomrule
		\endfoot}
	{\end{longtable}\endgroup}

% =====================================================================
% Convenciones tipográficas en línea
% =====================================================================
\newcommand{\cmd}[1]{\texttt{\textcolor{colcmd}{#1}}}      % ejecutables
\newcommand{\opt}[1]{\texttt{\textcolor{colopt}{#1}}}      % opciones
\newcommand{\arch}[1]{\texttt{\textcolor{colarch}{#1}}}    % archivos, rutas
\newcommand{\val}[1]{\texttt{#1}}                          % valores literales
\newcommand{\oblig}[1]{\texttt{<\textit{#1}>}}             % <obligatorio>
\newcommand{\opcional}[1]{\texttt{[\textit{#1}]}}          % [opcional]
\newcommand{\tecla}[1]{\tcbox[on line, boxsep=0pt, left=3pt, right=3pt,
	top=1.5pt, bottom=1.5pt, colback=white, colframe=codigoborde,
	boxrule=0.5pt, arc=2pt]{\footnotesize\sffamily #1}}
\newcommand{\compat}{\textcolor{colconsejo}{\ding{51}}}
\newcommand{\nocompat}{\textcolor{colimport}{\ding{55}}}
\newcommand{\condic}{\textcolor{coladvert}{$\boldsymbol{\sim}$}}
\newcommand{\cab}[1]{\textbf{#1}}
\newcommand{\cabmc}[2]{\multicolumn{1}{>{\centering\arraybackslash}b{#1}}{\cab{#2}}}
\newcommand{\cabml}[2]{\multicolumn{1}{>{\raggedright\arraybackslash}b{#1}}{\cab{#2}}}

% =====================================================================
% Cajas de aviso
% =====================================================================
\newcommand{\nuevacaja}[3]{%
	\newtcolorbox{#1}[1][]{%
		enhanced, breakable,
		colback=#3!6!white, colframe=#3,
		boxrule=0pt, leftrule=3pt, arc=0pt, outer arc=0pt,
		left=8pt, right=8pt, top=5pt, bottom=5pt,
		before skip=9pt, after skip=9pt,
		parbox=false,
		title={#2\ifx\relax##1\relax\else\ ---\ ##1\fi},
		fonttitle=\bfseries, coltitle=#3!80!black,
		detach title,
		before upper={\tcbtitle\textcolor{#3!80!black}{\textbf{.}}\enspace},
	}%
}
\nuevacaja{nota}{Nota}{colnota}
\nuevacaja{consejo}{Consejo}{colconsejo}
\nuevacaja{advertencia}{Advertencia}{coladvert}
\nuevacaja{importante}{Importante}{colimport}
\nuevacaja{limitacion}{Limitación}{collimit}
\nuevacaja{ejemplo}{Ejemplo}{colejemplo}

% =====================================================================
% Diagramas de flujo: entrada -> herramienta -> salida
%   Elementos impares = datos/archivos; pares = herramientas/comandos.
%   \flujo{a,b,c}  (vertical)      \flujoh{a,b,c}  (horizontal)
% =====================================================================
\tikzset{
	fdato/.style={draw=codigoborde, fill=codigofondo, rounded corners=2pt,
		font=\small\ttfamily, inner xsep=7pt, inner ysep=4pt, align=center},
	fherr/.style={draw=acento, fill=acentoclaro, rounded corners=2pt,
		font=\small\sffamily\bfseries, text=acento, inner xsep=7pt,
		inner ysep=4pt, align=center},
	every join/.style={-{Stealth[length=2.2mm]}, semithick, draw=acento!75!black},
}
\newcommand{\flujo}[1]{%
	\begin{center}
		\begin{tikzpicture}[start chain=going below, node distance=4.5mm]
			\foreach \x [count=\i] in {#1} {%
				\ifodd\i\relax
				\node[fdato, on chain, join] {\x};
				\else
				\node[fherr, on chain, join] {\x};
				\fi
			}
		\end{tikzpicture}
\end{center}}
\newcommand{\flujoh}[1]{%
	\begin{center}
		\begin{tikzpicture}[start chain=going right, node distance=6mm]
			\foreach \x [count=\i] in {#1} {%
				\ifodd\i\relax
				\node[fdato, on chain, join] {\x};
				\else
				\node[fherr, on chain, join] {\x};
				\fi
			}
		\end{tikzpicture}
\end{center}}

% =====================================================================
% Atajos de redacción
% =====================================================================
\newcommand{\wsl}{WSL2}
\newcommand{\ubu}{Ubuntu 26.04 LTS}

% =====================================================================
\begin{document}
	% =====================================================================
	
	% ---------------------------------------------------------------------
	% PORTADA
	% ---------------------------------------------------------------------
	\begin{titlepage}
		\thispagestyle{empty}
		\begin{tikzpicture}[remember picture, overlay]
			\fill[acento] (current page.north west) rectangle ([yshift=-8.3cm]current page.north east);
			\node[anchor=north west, text=white, font=\sffamily\bfseries\fontsize{30}{36}\selectfont,
			align=left, text width=17cm]
			at ([xshift=2.3cm, yshift=-1.4cm]current page.north west)
			{Guía técnica de herramientas\\de línea de comandos};
			\node[anchor=north west, text=white, font=\sffamily\fontsize{20}{24}\selectfont,
			align=left, text width=17cm]
			at ([xshift=2.3cm, yshift=-4cm]current page.north west)
			{\textit{Angel David Ojeda Peralta · Physicist} | \textit{Database Specialist}};
			\node[anchor=north west, text=acentoclaro, font=\sffamily\Large, align=left, text width=17cm]
			at ([xshift=2.3cm, yshift=-5.5cm]current page.north west)
			{Documentos, imágenes y PDF desde la terminal de\\Ubuntu 26.04 LTS en WSL2 (Windows 11)};
		\end{tikzpicture}
		
		\vspace*{6.4cm}
		
		{\sffamily\large\color{acento}\textbf{Manual de consulta rápida y referencia técnica}}
		
		\vspace{0.6cm}
		
		\begin{tcbraster}[raster columns=3, raster equal height, raster column skip=6pt,
			raster row skip=6pt, colback=acentoclaro, colframe=acentoclaro, boxrule=0pt,
			arc=2pt, left=6pt, right=6pt, top=5pt, bottom=5pt, fontupper=\sffamily\small]
			\begin{tcolorbox}\textbf{\color{acento}Marker}\\ PDF y documentos $\rightarrow$ Markdown, JSON, HTML\end{tcolorbox}
			\begin{tcolorbox}\textbf{\color{acento}Poppler}\\ Inspeccionar, extraer y renderizar PDF\end{tcolorbox}
			\begin{tcolorbox}\textbf{\color{acento}img2pdf}\\ Imágenes $\rightarrow$ PDF sin recompresión\end{tcolorbox}
			\begin{tcolorbox}\textbf{\color{acento}ImageMagick}\\ Procesamiento y conversión de imágenes\end{tcolorbox}
			\begin{tcolorbox}\textbf{\color{acento}Potrace}\\ Vectorización: mapa de bits $\rightarrow$ SVG/PDF/EPS\end{tcolorbox}
			\begin{tcolorbox}\textbf{\color{acento}Pandoc}\\ Conversión entre formatos de documento\end{tcolorbox}
			\begin{tcolorbox}\textbf{\color{acento}QPDF}\\ Estructura del PDF: páginas, cifrado, reparación\end{tcolorbox}
			\begin{tcolorbox}\textbf{\color{acento}OCRmyPDF}\\ Capa de texto OCR para PDF escaneados\end{tcolorbox}
			\begin{tcolorbox}\textbf{\color{acento}Ghostscript}\\ Intérprete PostScript/PDF: rasterizar, reescribir, comprimir\end{tcolorbox}
		\end{tcbraster}
		
		\vspace{0.9cm}
		
		\begin{tcolorbox}[enhanced, colback=white, colframe=codigoborde, boxrule=0.5pt, arc=2pt,
			left=10pt, right=10pt, top=8pt, bottom=8pt, fontupper=\small]
			\textbf{\color{acento}Versiones de referencia (repositorios oficiales de Ubuntu 26.04 «resolute»)}\par\smallskip
			\begin{tabular}{@{}L{3.9cm}L{2.3cm}L{3.4cm}L{4.6cm}@{}}
				\textbf{Herramienta} & \textbf{Versión} & \textbf{Herramienta} & \textbf{Versión}\\[2pt]
				Poppler (\texttt{poppler-utils}) & 26.01 & Ghostscript & 10.06.0\\
				img2pdf & 0.6.2 & QPDF & 12.3.2\\
				ImageMagick (\texttt{magick}) & 7.1.2-18 & OCRmyPDF & 16.13.0\\
				Potrace & 1.16 & Tesseract & 5.5.0\\
				Pandoc & 3.7.0.2 & Marker (\texttt{pip}) & serie 2.x, en \texttt{\textasciitilde/marker-env} (Python 3.12)\\
			\end{tabular}
		\end{tcolorbox}
		
		\vfill
		{\small\color{collimit}Las opciones documentadas se contrastaron con la ayuda integrada de cada programa, sus páginas \texttt{man}, el código o la documentación oficial de cada proyecto. Si su sistema muestra otra versión, compruébela con los comandos de la sección~\ref{sec:instalacion}.}
		
		\vspace{0.4cm}
		{\sffamily\color{acento}Octubre de 2026}
	\end{titlepage}
	
	% ---------------------------------------------------------------------
	% ÍNDICE
	% ---------------------------------------------------------------------
	\pagenumbering{roman}
	\begingroup
	\hypersetup{linkcolor=black}
	\tableofcontents
	\endgroup
	\clearpage
	\pagenumbering{arabic}
	% =====================================================================
	% Atajos para caracteres especiales en texto
	% =====================================================================
	\newcommand{\bs}{\textbackslash{}}
	\newcommand{\tl}{\textasciitilde{}}
	
	% =====================================================================
	\section{Introducción}\label{sec:intro}
	% =====================================================================
	
	\necesitaespacio{8\baselineskip}
	\subsection{Propósito de la guía y cómo usarla}
	
	Esta guía es un manual de consulta para trabajar desde la terminal de Ubuntu con nueve herramientas que procesan documentos, imágenes y archivos PDF. No está pensada para leerse de principio a fin: se abre cuando hay una tarea concreta («quiero pasar este PDF a Markdown», «necesito unir diez PDF», «este escaneo no tiene texto seleccionable») y se busca el comando adecuado.
	
	Cada una de las nueve herramientas tiene un capítulo con la misma estructura, para que siempre sepa dónde mirar:
	
	\begin{enumerate}[label=\arabic*.]
		\item \textbf{Descripción general}: qué es, qué problema resuelve, cuándo conviene usarla y cuándo no.
		\item \textbf{Instalación y comprobación} en Ubuntu 26.04.
		\item \textbf{Comandos}: sintaxis, explicación, tabla de parámetros y ejemplos copiables.
		\item \textbf{Formatos de entrada y salida}, distinguiendo soporte nativo, condicionado o inexistente.
		\item \textbf{Limitaciones} prácticas.
		\item \textbf{Errores frecuentes}, con causa probable y solución, y cómo pedir ayuda a la propia herramienta.
	\end{enumerate}
	
	Después de los capítulos individuales vienen las comparaciones entre herramientas (sección~\ref{sec:comparacion}), los flujos de trabajo que combinan varias de ellas (sección~\ref{sec:flujos}), las tablas de formatos, dependencias y comandos, los errores comunes a todo el entorno, la seguridad, las buenas prácticas y, al final, una referencia rápida en el apéndice~\ref{sec:refrapida}.
	
	\necesitaespacio{14\baselineskip}
	\subsection{Localizador rápido: ¿qué necesito?}\label{sec:localizador}
	
	La siguiente tabla responde a las preguntas más habituales indicando la herramienta recomendada, el comando esencial y la sección donde se explica con detalle. Los comandos completos, con sus variantes, están en cada sección.
	
	\begingroup\small
	\necesitaespacio{10\baselineskip}
	\begin{longtable}{L{4.1cm}L{2.6cm}L{6.26cm}C{1.6cm}}
		\toprule\rowcolor{tablacab}\cab{Quiero\ldots} & \cab{Herramienta} & \cab{Comando esencial} & \cab{Sección}\\\midrule
		\endfirsthead
		\toprule\rowcolor{tablacab}\cab{Quiero\ldots} & \cab{Herramienta} & \cab{Comando esencial} & \cab{Sección}\\\midrule
		\endhead
		\bottomrule
		\endfoot
		Convertir un PDF en imágenes & Poppler & \cmd{pdftoppm -png -r 150 doc.pdf pag} & \ref{sec:pdftoppm}\\*
		Extraer el texto de un PDF digital & Poppler & \cmd{pdftotext -layout doc.pdf doc.txt} & \ref{sec:pdftotext}\\
		Hacer OCR a un PDF escaneado & OCRmyPDF & \cmd{ocrmypdf -l spa entrada.pdf salida.pdf} & \ref{sec:ocrmypdf}\\
		Convertir un PDF a Markdown & Marker & \cmd{marker\_\allowbreak{}single doc.pdf -{}-output\_\allowbreak{}dir md} & \ref{sec:marker}\\
		Unir varios PDF & QPDF / Poppler & \cmd{qpdf -{}-empty -{}-pages a.pdf b.pdf -{}- unido.pdf} & \ref{sec:qpdf-paginas}\\
		Separar un PDF por páginas & QPDF / Poppler & \cmd{qpdf -{}-split-\allowbreak{}pages doc.pdf pag-\%d.pdf} & \ref{sec:qpdf-paginas}\\
		Inspeccionar o reparar un PDF & QPDF & \cmd{qpdf -{}-check doc.pdf} & \ref{sec:qpdf-inspeccion}\\
		Reducir el tamaño de un PDF & Ghostscript & \cmd{gs -sDEVICE=\allowbreak{}pdfwrite -dPDFSETTINGS=\allowbreak{}/\allowbreak{}ebook -o s.pdf e.pdf} & \ref{sec:gs-reduccion}\\
		Rasterizar un PDF (páginas a mapa de bits) & Poppler / Ghostscript & \cmd{pdftoppm -r 300 -png doc.pdf p} & \ref{sec:pdftoppm}\\
		Convertir imágenes en un PDF & img2pdf & \cmd{img2pdf *.jpg -o album.pdf} & \ref{sec:img2pdf}\\
		Convertir una imagen raster a SVG & ImageMagick + Potrace & \cmd{potrace logo.pbm -s -o logo.svg} & \ref{sec:potrace}\\
		Convertir Markdown a DOCX o PDF & Pandoc & \cmd{pandoc doc.md -o doc.docx} & \ref{sec:pandoc}\\
		Redimensionar o convertir imágenes & ImageMagick & \cmd{magick foto.png -resize 50\% foto.jpg} & \ref{sec:imagemagick}\\
		Procesar cientos de archivos & Bash + cualquiera & \cmd{for f in *.pdf; do ...; done} & \ref{sec:bash}\\*
		Decidir qué herramienta usar & --- & Tablas comparativas & \ref{sec:comparacion}\\
	\end{longtable}
	\endgroup
	
	\necesitaespacio{22\baselineskip}
	\subsection{Entorno de referencia y versiones verificadas}\label{sec:versiones}
	
	La guía está escrita para \textbf{Ubuntu 26.04 LTS («resolute») ejecutándose en WSL2 sobre Windows 11}, con las versiones que instala \cmd{apt} desde los repositorios oficiales de esa versión de Ubuntu. Marker es la excepción: no está en los repositorios y se instala con \cmd{pip} dentro de un entorno virtual de Python; en esta guía se supone el entorno \arch{\tl/marker-env} con Python 3.12.
	
	\begingroup\small
	\necesitaespacio{18\baselineskip}
	\begin{longtable}{L{2.5cm}L{4.0cm}L{2.4cm}L{5.16cm}}
		\toprule\rowcolor{tablacab}\cab{Herramienta} & \cab{Paquete} & \cab{Versión} & \cab{Comprobar la versión instalada}\\\midrule
		\endfirsthead
		\toprule\rowcolor{tablacab}\cab{Herramienta} & \cab{Paquete} & \cab{Versión} & \cab{Comprobar la versión instalada}\\\midrule
		\endhead
		\bottomrule
		\endfoot
		Marker & \texttt{marker-pdf} (PyPI) & serie 2.x & \cmd{pip show marker-\allowbreak{}pdf} (con el entorno activo)\\*
		Poppler & \texttt{poppler-utils} & 26.01.0 & \cmd{pdftoppm -v}\\
		img2pdf & \texttt{img2pdf} & 0.6.2 & \cmd{img2pdf -{}-version}\\
		ImageMagick & \texttt{imagemagick} & 7.1.2-18 & \cmd{magick -version}\\
		Potrace & \texttt{potrace} & 1.16 & \cmd{potrace -{}-version}\\
		Pandoc & \texttt{pandoc} & 3.7.0.2 & \cmd{pandoc -{}-version}\\
		QPDF & \texttt{qpdf} & 12.3.2 & \cmd{qpdf -{}-version}\\
		OCRmyPDF & \texttt{ocrmypdf} & 16.13.0 & \cmd{ocrmypdf -{}-version}\\
		Tesseract & \texttt{tesseract-ocr} & 5.5.0 & \cmd{tesseract -{}-version}\\*
		Ghostscript & \texttt{ghostscript} & 10.06.0 & \cmd{gs -{}-version}\\
	\end{longtable}
	\endgroup
	
	\begin{importante}[Versiones]
		Las opciones de esta guía corresponden a esas versiones. Algunas opciones que aparecen en documentación más reciente en Internet \emph{no existen} en las versiones de Ubuntu 26.04: por ejemplo, \opt{-{}-syntax-\allowbreak{}highlighting} (Pandoc 3.8 y posteriores) o \opt{-{}-mode} (OCRmyPDF 17). Si un comando de Internet falla con «opción desconocida», compare primero la versión.
	\end{importante}
	
	\necesitaespacio{8\baselineskip}
	\subsection{Las nueve herramientas en una frase}
	
	\begin{description}[style=nextline, leftmargin=1.2em, font=\normalfont\bfseries\color{acento}]
		\item[Marker] Convierte PDF (y, con extras, imágenes, DOCX, PPTX, XLSX, HTML y EPUB) en Markdown, JSON, HTML o fragmentos, reconstruyendo títulos, tablas y ecuaciones con modelos de aprendizaje automático.
		\item[Poppler] Conjunto de utilidades (\cmd{pdfinfo}, \cmd{pdftotext}, \cmd{pdftoppm}\ldots) para inspeccionar, extraer y renderizar PDF.
		\item[img2pdf] Empaqueta imágenes en un PDF sin recomprimirlas cuando el formato lo permite (JPEG, por ejemplo).
		\item[ImageMagick] Navaja suiza del procesamiento de imágenes: convertir, redimensionar, recortar, ajustar, combinar, anotar.
		\item[Potrace] Vectoriza mapas de bits en blanco y negro: de una imagen de píxeles a curvas SVG, PDF o EPS.
		\item[Pandoc] Convierte entre formatos de documento (Markdown, DOCX, HTML, LaTeX, EPUB, ODT\ldots) y genera PDF mediante un motor externo.
		\item[QPDF] Transforma la \emph{estructura} del PDF: páginas, cifrado, linealización, reparación, inspección.
		\item[OCRmyPDF] Añade una capa de texto reconocido por OCR (con Tesseract) a PDF escaneados.
		\item[Ghostscript] Intérprete de PostScript y PDF: rasteriza, reescribe, convierte y reduce el tamaño de documentos.
	\end{description}
	
	% =====================================================================
	\clearpage
	\section{Entorno Ubuntu + WSL2}\label{sec:wsl}
	% =====================================================================
	
	Este capítulo no es una guía completa de WSL2: recoge solo lo necesario para que los comandos de las nueve herramientas encuentren sus archivos, funcionen con nombres reales (con espacios y acentos) y rindan bien.
	
	\necesitaespacio{8\baselineskip}
	\subsection{Dos sistemas de archivos en la misma máquina}
	
	WSL2 ejecuta Ubuntu dentro de una pequeña máquina virtual. Eso significa que conviven dos sistemas de archivos:
	
	\begin{itemize}
		\item El \textbf{sistema de archivos de Linux}, propio de Ubuntu, donde está su carpeta personal \arch{/home/USUARIO}, abreviada \arch{\tl}. Aquí los programas de Linux trabajan a máxima velocidad.
		\item Las \textbf{unidades de Windows}, que Ubuntu ve «montadas» bajo \arch{/mnt}: la unidad \arch{C:} aparece como \arch{/mnt/c}, la \arch{D:} como \arch{/mnt/d}, y así sucesivamente.
	\end{itemize}
	
	\begingroup\small
	\necesitaespacio{13\baselineskip}
	\begin{longtable}{L{5.16cm}L{5.4cm}L{4.5cm}}
		\toprule\rowcolor{tablacab}\cab{Ruta vista desde Windows} & \cab{La misma ruta vista desde Ubuntu} & \cab{Comentario}\\\midrule
		\endfirsthead
		\toprule\rowcolor{tablacab}\cab{Ruta vista desde Windows} & \cab{La misma ruta vista desde Ubuntu} & \cab{Comentario}\\\midrule
		\endhead
		\bottomrule
		\endfoot
		\arch{C:\bs{}\allowbreak{}Users\bs{}\allowbreak{}NOMBRE\bs{}\allowbreak{}Downloads} & \arch{/\allowbreak{}mnt/\allowbreak{}c/\allowbreak{}Users/\allowbreak{}NOMBRE/\allowbreak{}Downloads} & Descargas de Windows\\*
		\arch{C:\bs{}\allowbreak{}Users\bs{}\allowbreak{}NOMBRE\bs{}\allowbreak{}Documents} & \arch{/\allowbreak{}mnt/\allowbreak{}c/\allowbreak{}Users/\allowbreak{}NOMBRE/\allowbreak{}Documents} & «Documentos» se llama \arch{Documents} en disco\\
		\arch{D:\bs{}Proyectos} & \arch{/\allowbreak{}mnt/\allowbreak{}d/\allowbreak{}Proyectos} & Otras unidades\\*
		\arch{\bs\bs{}\allowbreak{}wsl.localhost\bs{}\allowbreak{}Ubuntu\bs{}\allowbreak{}home\bs{}\allowbreak{}USUARIO} & \arch{/home/USUARIO} o \arch{\tl} & Carpeta personal de Linux\\
	\end{longtable}
	\endgroup
	
	Las reglas de traducción manual son tres: la letra de unidad \arch{C:} pasa a ser \arch{/mnt/c}; las barras invertidas \arch{\bs} se convierten en barras normales \arch{/}; y las comillas siguen siendo necesarias si hay espacios.
	
	\begin{nota}[Nombres de usuario]
		En los ejemplos, \texttt{NOMBRE} es su usuario de \emph{Windows} (la carpeta dentro de \arch{C:\bs{}Users}) y \texttt{USUARIO} su usuario de \emph{Linux}. Pueden coincidir o no; por ejemplo, si su usuario de Windows es \texttt{ana}, la ruta de descargas es \arch{/\allowbreak{}mnt/\allowbreak{}c/\allowbreak{}Users/\allowbreak{}ana/\allowbreak{}Downloads}. Compruebe los nombres reales con \cmd{ls /\allowbreak{}mnt/\allowbreak{}c/\allowbreak{}Users} y \cmd{whoami}.
	\end{nota}
	
	\necesitaespacio{8\baselineskip}
	\subsection{Acceder a archivos de Windows desde Ubuntu}
	
	\begin{lstlisting}
		cd /mnt/c/Users/NOMBRE/Downloads     # ir a Descargas de Windows
		pwd                                  # mostrar dónde está
		ls -l                                # listar con detalles
		ls -l *.pdf                          # listar solo los PDF
\end{lstlisting}
	
	Si usa a menudo la misma carpeta de Windows, cree un \textbf{enlace simbólico} (un acceso directo de Linux) en su carpeta personal:
	\newpage
	
	\begin{lstlisting}
		ln -s /mnt/c/Users/NOMBRE/Downloads ~/descargas
		cd ~/descargas          # ahora es equivalente a la ruta larga
\end{lstlisting}
	
	\necesitaespacio{8\baselineskip}
	\subsection{Acceder a archivos de Ubuntu desde Windows}
	
	Desde Windows, los archivos de Ubuntu están en la ruta de red \arch{\bs\bs{}\allowbreak{}wsl.localhost\bs{}\allowbreak{}Ubuntu\bs{}\allowbreak{}home\bs{}\allowbreak{}USUARIO} (el nombre de la distribución puede variar: compruébelo en PowerShell con \cmd{wsl -l -v}). Desde la propia terminal de Ubuntu hay dos atajos muy útiles:
	
	\begin{lstlisting}
		explorer.exe .  # abre la carpeta actual en el Explorador
		# abre un archivo con su programa de Windows
		explorer.exe "$(wslpath -w ~/trabajo/salida.pdf)"
\end{lstlisting}
	
	\necesitaespacio{18\baselineskip}
	\subsection{Traducir rutas con \texttt{wslpath}}
	
	\cmd{wslpath} convierte rutas entre los dos mundos sin errores de transcripción.
	
	\necesitaespacio{14\baselineskip}
	\begin{opciones}
		\opt{wslpath 'C:\bs{}\allowbreak{}...'} & Sin opciones (equivale a \opt{-u}): de ruta de Windows a ruta de Linux. Use comillas \emph{simples} para que Bash no interprete las barras invertidas.\\*
		\opt{-w} & De ruta de Linux a ruta de Windows con barras invertidas.\\
		\opt{-m} & De ruta de Linux a ruta de Windows con barras normales (\arch{C:/Users/...}).\\*
		\opt{-a} & Fuerza el resultado como ruta absoluta.\\
	\end{opciones}
	
	\begin{lstlisting}
  wslpath 'C:\Users\NOMBRE\Downloads'       # -> /mnt/c/Users/NOMBRE/Downloads
  wslpath -w ~/trabajo  # -> \\wsl.localhost\Ubuntu\home\USUARIO\trabajo
  cd "$(wslpath 'C:\Users\NOMBRE\Mis Documentos')"  # cd con ruta copiada de Windows
\end{lstlisting}
	
	\necesitaespacio{8\baselineskip}
	\subsection{Rendimiento: dónde conviene trabajar}\label{sec:rendimiento}
	
	El acceso de Ubuntu a \arch{/mnt/c} atraviesa la frontera entre la máquina virtual y Windows, y es bastante más lento que el acceso al sistema de archivos propio de Linux. Para un archivo aislado la diferencia no importa; para procesar cientos de páginas con OCR, Marker o ImageMagick, sí. La práctica recomendada es:
	
	\flujoh{/mnt/c/\ldots,cp,\textasciitilde/trabajo/entrada,procesar,\textasciitilde/trabajo/salida}
	\newpage
	
	\begin{lstlisting}
		mkdir -p ~/trabajo/entrada ~/trabajo/salida
		cp /mnt/c/Users/NOMBRE/Downloads/*.pdf ~/trabajo/entrada/
		# ... procesar dentro de ~/trabajo ...
		cp ~/trabajo/salida/* /mnt/c/Users/NOMBRE/Downloads/
\end{lstlisting}
	
	\necesitaespacio{8\baselineskip}
	\subsection{Diferencias prácticas entre rutas de Linux y de Windows}\label{sec:diferencias-wsl}
	
	\paragraph{Mayúsculas y minúsculas} En el sistema de archivos de Linux, \arch{Informe.pdf} e \arch{informe.pdf} son archivos distintos. Las carpetas de Windows, en cambio, normalmente no distinguen entre ambas. Un script que funciona en \arch{/mnt/c} puede fallar en \arch{\tl} (o al revés) solo por una mayúscula.
	
	\paragraph{Caracteres no válidos en Windows} Linux admite en un nombre de archivo casi cualquier carácter, pero Windows no admite \texttt{\bs{} / : * ? " < > |}. Si una herramienta genera nombres con dos puntos (por ejemplo, una hora como \texttt{12:30}) y la salida está en \arch{/mnt/c}, la escritura puede fallar o el nombre puede verse alterado desde Windows. Para las salidas en carpetas de Windows, use nombres sin esos caracteres.
	
	\paragraph{Finales de línea} Los editores de Windows pueden guardar los scripts con finales de línea CRLF; Bash los interpreta mal y aparecen errores como \texttt{\$'\bs{}r': command not found}. Guarde los scripts con finales LF (en VS Code, indicador «CRLF/LF» de la barra inferior) o corríjalos con \cmd{sed -i 's/\bs{}r\$//' script.sh}.
	
	\paragraph{Archivos \texttt{Zone.Identifier}} Al copiar archivos con el Explorador de Windows hacia carpetas de Ubuntu pueden aparecer archivos acompañantes con nombres como \arch{informe.pdf:Zone.Identifier}. No son documentos: son metadatos de origen de Windows. Molestan en los procesos por lotes que recorren \emph{todos} los archivos de una carpeta (por ejemplo, el comando \cmd{marker}). Para localizarlos y, tras revisarlos, borrarlos:
	
	\begin{lstlisting}
	find ~/trabajo -name '*:Zone.Identifier'           # 1) revisar la lista
	find ~/trabajo -name '*:Zone.Identifier' -delete   # 2) borrar (irreversible)
\end{lstlisting}
	
	\necesitaespacio{8\baselineskip}
	\subsection{Memoria y procesadores disponibles en WSL2}\label{sec:wsl-memoria}
	
	WSL2 asigna a Ubuntu solo una parte de la memoria del equipo (por defecto, una fracción de la RAM total de Windows). Las herramientas pesadas (Marker, OCRmyPDF con muchos procesos, ImageMagick con imágenes enormes) pueden quedarse sin memoria aunque Windows tenga RAM libre. Compruebe los recursos reales desde Ubuntu:
	
	\begin{lstlisting}
		free -h      # memoria total, usada y disponible dentro de WSL2
		nproc        # número de procesadores visibles
		df -h ~      # espacio libre en el disco de Linux
\end{lstlisting}
	
	Para cambiar los límites, cree o edite en Windows el archivo \arch{\%UserProfile\%\bs{}\allowbreak{}.wslconfig} (es decir, \arch{C:\bs{}\allowbreak{}Users\bs{}\allowbreak{}NOMBRE\bs{}\allowbreak{}.wslconfig}) y reinicie WSL con \cmd{wsl -{}-shutdown} desde PowerShell. También puede hacerse desde la aplicación «Configuración de WSL» de Windows.
	\newpage
	
	\begin{lstlisting}[style=salida]
										[wsl2]
										memory=8GB
										processors=4
										swap=4GB
\end{lstlisting}
	
	\begin{advertencia}
		No asigne a WSL2 toda la memoria del equipo: Windows también la necesita. Ajuste los valores a su equipo y verifique el resultado con \cmd{free -h} tras reiniciar WSL.
	\end{advertencia}
	
	\necesitaespacio{17\baselineskip}
	\subsection{Rutas absolutas y relativas}
	
	Una \textbf{ruta absoluta} empieza en la raíz \arch{/} y funciona desde cualquier carpeta: \arch{/\allowbreak{}mnt/\allowbreak{}c/\allowbreak{}Users/\allowbreak{}NOMBRE/\allowbreak{}Downloads/\allowbreak{}doc.pdf}. Una \textbf{ruta relativa} se interpreta desde la carpeta actual (la que muestra \cmd{pwd}).
	
	\begingroup\small
	\necesitaespacio{13\baselineskip}
	\begin{longtable}{L{2.4cm}L{12.66cm}}
		\toprule\rowcolor{tablacab}\cab{Símbolo} & \cab{Significado}\\\midrule
		\endfirsthead
		\bottomrule
		\endfoot
		\arch{/} & Raíz del sistema de archivos de Linux.\\*
		\arch{\tl} & Su carpeta personal (\arch{/home/USUARIO}). No se expande dentro de comillas: escriba \arch{\tl/\allowbreak{}"Mis archivos"} o \arch{"\$HOME/\allowbreak{}Mis archivos"}, no \arch{"\tl/\allowbreak{}Mis archivos"}.\\
		\arch{.} & La carpeta actual. \cmd{explorer.exe .} abre «aquí».\\
		\arch{..} & La carpeta superior. \cmd{cd ..} sube un nivel.\\*
		\cmd{cd -} & Vuelve a la carpeta anterior.\\
	\end{longtable}
	\endgroup
	
	\necesitaespacio{8\baselineskip}
	\subsection{Espacios, comillas y caracteres especiales}\label{sec:comillas}
	
	Bash separa los argumentos por espacios. Un nombre como \arch{Mis Documentos} se interpretaría como \emph{dos} argumentos (\arch{Mis} y \arch{Documentos}) si no se protege. Hay tres formas de protegerlo:
	
	\begin{lstlisting}
		# 1) Comillas dobles: la forma más habitual
		pdfinfo "/mnt/c/Users/NOMBRE/Mis Documentos/informe final.pdf"
		
		# 2) Barra invertida delante de cada espacio
		pdfinfo /mnt/c/Users/NOMBRE/Mis\ Documentos/informe\ final.pdf
		
		# 3) Comillas simples: protegen TODO, también $ y \
		pdfinfo '/mnt/c/Users/NOMBRE/Mis Documentos/informe final.pdf'
\end{lstlisting}
	
	\begingroup\small
	\necesitaespacio{21\baselineskip}
	\begin{longtable}{L{3.6cm}L{11.46cm}}
		\toprule\rowcolor{tablacab}\cab{Mecanismo} & \cab{Qué protege y cuándo usarlo}\\\midrule
		\endfirsthead
		\toprule\rowcolor{tablacab}\cab{Mecanismo} & \cab{Qué protege y cuándo usarlo}\\\midrule
		\endhead
		\bottomrule
		\endfoot
		\texttt{"comillas dobles"} & Espacios y la mayoría de caracteres especiales; \emph{sí} expanden variables (\texttt{\$f}, \texttt{\$HOME}) y sustituciones \texttt{\$(...)}. Es la opción por defecto para rutas y para variables: escriba siempre \texttt{"\$f"}, nunca \texttt{\$f} a secas.\\*
		\texttt{'comillas simples'} & Todo de forma literal: no expanden nada. Ideales para rutas de Windows con \texttt{\bs} y para textos con \texttt{\$} o \texttt{!}. No pueden contener una comilla simple.\\
		\texttt{\bs} (barra invertida) & Protege el carácter siguiente: \texttt{Mis\bs{} Documentos}. Al final de una línea, continúa el comando en la línea siguiente.\\
		Tecla \tecla{Tab} & Autocompleta nombres y añade los escapes necesarios. Es la forma más segura de escribir rutas complicadas.\\*
		\cmd{printf '\%q\bs{}n'} & Muestra cómo escribir un nombre «raro» de forma segura: \cmd{printf '\%q\bs{}\allowbreak{}n' "año (v2).pdf"}.\\
	\end{longtable}
	\endgroup
	
	Caracteres que exigen protección porque Bash les da un significado especial: espacio, \texttt{\& ; | ( ) < > \$ ` \bs{}\allowbreak{} " ' * ? [ ] \# \tl{} !} y el salto de línea. Los acentos y la eñe no necesitan protección en Ubuntu 26.04 (que usa UTF-8), pero sí las comillas si el nombre contiene además espacios.
	
	\paragraph{Nombres que empiezan por guion} Un archivo llamado \arch{-borrador.pdf} se confundiría con una opción. Use \arch{./-borrador.pdf} o el separador \opt{-{}-}, que en la mayoría de programas indica «fin de las opciones»: \cmd{ls -l -{}- -borrador.pdf}.
	
	\paragraph{Rutas largas} Guarde las rutas en variables, entre comillas, y úselas siempre entre comillas dobles:
	
	\begin{lstlisting}
	ORIGEN="/mnt/c/Users/NOMBRE/OneDrive/Documentos/Congreso 2026/Presentación"
	ls "$ORIGEN"
	pdfinfo "$ORIGEN/Presentación final.pdf"
\end{lstlisting}
	
	\begin{consejo}
		Cuando un comando con una ruta larga falle, ejecute primero \cmd{ls -l "\$ORIGEN"}: si \cmd{ls} no encuentra el archivo, el problema es la ruta, no la herramienta.
	\end{consejo}
	% =====================================================================
	\clearpage
	\section{Instalación y comprobación de las herramientas}\label{sec:instalacion}
	% =====================================================================
	
	\necesitaespacio{8\baselineskip}
	\subsection{Cómo se instala cada herramienta en Ubuntu 26.04}
	
	En Ubuntu hay tres vías de instalación, y conviene saber cuál se usa en cada caso:
	
	\begin{itemize}
		\item \textbf{\cmd{apt}}: el gestor de paquetes de Ubuntu. Instala versiones probadas e integradas con el sistema y las actualiza con el resto del sistema. Es la vía recomendada para ocho de las nueve herramientas.
		\item \textbf{Entorno virtual de Python (\cmd{venv}) + \cmd{pip}}: carpeta aislada con su propio Python y sus propios paquetes. Es la vía para Marker, que no está en los repositorios de Ubuntu y depende de PyTorch.
		\item \textbf{\cmd{pipx}}: instala programas de línea de comandos escritos en Python, cada uno en su propio entorno aislado. Solo es útil si necesita una versión más nueva que la de \cmd{apt} (por ejemplo, de OCRmyPDF).
	\end{itemize}
	
	\necesitaespacio{8\baselineskip}
	\subsection{Instalación con \texttt{apt}}
	
	\begin{lstlisting}
	sudo apt update
	sudo apt install poppler-utils img2pdf imagemagick potrace pandoc qpdf \
	ocrmypdf tesseract-ocr-spa ghostscript unpaper pngquant jbig2
\end{lstlisting}
	
	\cmd{sudo} ejecuta la orden con permisos de administrador (pide su contraseña de Linux, que no se ve al escribirla). La barra invertida al final de la primera línea de \cmd{apt install} solo continúa el comando en la línea siguiente.
	
	\begingroup\small
	\necesitaespacio{10\baselineskip}
	\begin{longtable}{L{3.6cm}L{2.6cm}L{8.86cm}}
		\toprule\rowcolor{tablacab}\cab{Paquete} & \cab{Herramienta} & \cab{Qué aporta}\\\midrule
		\endfirsthead
		\toprule\rowcolor{tablacab}\cab{Paquete} & \cab{Herramienta} & \cab{Qué aporta}\\\midrule
		\endhead
		\bottomrule
		\endfoot
		\texttt{poppler-utils} & Poppler & \cmd{pdfinfo}, \cmd{pdffonts}, \cmd{pdftotext}, \cmd{pdftoppm}, \cmd{pdftocairo}, \cmd{pdfimages}, \cmd{pdftohtml}, \cmd{pdftops}, \cmd{pdfseparate}, \cmd{pdfunite}, \cmd{pdfdetach}, \cmd{pdfattach}, \cmd{pdfsig}.\\*
		\texttt{img2pdf} & img2pdf & El programa \cmd{img2pdf}.\\
		\texttt{imagemagick} & ImageMagick & ImageMagick 7: \cmd{magick} y los nombres heredados (\cmd{convert}, \cmd{mogrify}, \cmd{identify}\ldots).\\
		\texttt{potrace} & Potrace & \cmd{potrace} y el preprocesador \cmd{mkbitmap}.\\
		\texttt{pandoc} & Pandoc & El conversor \cmd{pandoc}. Para generar PDF necesita además LaTeX (abajo).\\
		\texttt{qpdf} & QPDF & \cmd{qpdf}, \cmd{fix-qdf} y \cmd{zlib-flate}.\\
		\texttt{ocrmypdf} & OCRmyPDF & \cmd{ocrmypdf}; instala Tesseract, Ghostscript y los datos de inglés como dependencias.\\
		\texttt{tesseract-\allowbreak{}ocr-\allowbreak{}spa} & Tesseract & Datos de reconocimiento del español (cada idioma es un paquete \texttt{tesseract-\allowbreak{}ocr-\allowbreak{}XXX}).\\
		\texttt{ghostscript} & Ghostscript & \cmd{gs} y los guiones \cmd{ps2pdf}, \cmd{pdf2ps}, \cmd{eps2eps}, \cmd{ps2ascii}\ldots\\
		\texttt{unpaper} & (OCRmyPDF) & Limpieza de escaneos para \opt{-{}-clean} y \opt{-{}-clean-final}.\\
		\texttt{pngquant} & (OCRmyPDF) & Cuantización de PNG para la optimización \opt{-O 2} y \opt{-O 3}.\\*
		\texttt{jbig2} & (OCRmyPDF) & Codificador JBIG2 (proyecto \emph{jbig2enc}): mejor compresión de páginas en blanco y negro.\\
	\end{longtable}
	\endgroup
	
	Para que Pandoc genere PDF a través de LaTeX (y para compilar esta guía):
	
	\begin{lstlisting}
		sudo apt install texlive-latex-recommended texlive-latex-extra \
		texlive-fonts-recommended texlive-lang-spanish lmodern latexmk
		# opcional: motor xelatex (Unicode y fuentes del sistema)
		sudo apt install texlive-xetex
\end{lstlisting}
	
	\begin{nota}[Paquetes recomendados]
		\cmd{apt} instala también los paquetes «recomendados» de cada programa. Si alguna función se queja de una dependencia ausente (por ejemplo, un idioma de Tesseract o el motor LaTeX), instálela explícitamente; la sección~\ref{sec:dependencias} resume qué depende de qué.
	\end{nota}
	
	\necesitaespacio{8\baselineskip}
	\subsection{Herramientas de Python: por qué no usar \texttt{sudo pip}}\label{sec:pep668}
	
	Ubuntu 26.04 protege el Python del sistema (versión 3.14): \cmd{pip install} fuera de un entorno virtual termina con \texttt{error: externally-\allowbreak{}managed-\allowbreak{}environment}. Esta protección es deliberada: las herramientas que instala \cmd{apt} (OCRmyPDF, img2pdf\ldots) usan bibliotecas de Python concretas, y si \cmd{pip} instala otras versiones encima, esas herramientas se rompen.
	
	\begin{importante}[No fuerce la instalación]
		Evite \cmd{sudo pip install} y la opción \opt{-{}-break-\allowbreak{}system-\allowbreak{}packages}. Un fallo real causado así: un \cmd{ocrmypdf} instalado con \cmd{apt} deja de arrancar con \texttt{ImportError: cannot import name 'PdfMatrix' from 'pikepdf'}, porque otro \texttt{pikepdf} instalado con \cmd{pip} en \arch{/usr/local} tiene prioridad sobre el del sistema.
	\end{importante}
	
	Las alternativas correctas son un entorno virtual (lo que usa Marker) o \cmd{pipx}:
	
	\begin{lstlisting}
	sudo apt install pipx
	pipx ensurepath  # añade ~/.local/bin al PATH (abra una terminal nueva después)
	pipx install ocrmypdf    # ejemplo: versión de PyPI, aislada de la de apt
	pipx list                # programas instalados con pipx
\end{lstlisting}
	\vspace{-0.2cm}
	
	\begin{advertencia}
		Si instala con \cmd{pipx} un programa que también tiene instalado con \cmd{apt}, ejecutará el que aparezca primero en el \texttt{PATH} (normalmente \arch{\tl/.local/bin}, es decir, el de \cmd{pipx}). Compruébelo con \cmd{type -a ocrmypdf}.
	\end{advertencia}
	
	\necesitaespacio{8\baselineskip}
	\subsection{Marker en el entorno \texttt{\textasciitilde/marker-env}}\label{sec:marker-instalacion}
	
	Marker se ejecuta dentro de su entorno virtual. Cada vez que abra una terminal nueva, actívelo antes de usarlo:
	
	\begin{lstlisting}
	source ~/marker-env/bin/activate  # el indicador pasa a empezar por (marker-env)
	python --version                     # debe mostrar Python 3.12.x
	pip show marker-pdf                  # versión de Marker instalada
	deactivate                           # salir del entorno cuando termine
\end{lstlisting}
	
	Si creó el entorno en otra ruta, sustituya \arch{\tl/marker-env} por la suya. El Python del sistema (3.14) y el del entorno (3.12) son independientes: no instale Marker en el primero.
	
	\paragraph{Actualizar Marker} Con el entorno activo:
	
	\begin{lstlisting}
	pip install --upgrade marker-pdf          # actualiza Marker y sus dependencias
	pip install --upgrade "marker-pdf[full]"  # además: DOCX, PPTX, XLSX, HTML y EPUB
\end{lstlisting}
	
	\paragraph{El servidor \texttt{llama-server}} En equipos sin GPU NVIDIA, Marker~2 necesita el programa \cmd{llama-server} (del proyecto llama.cpp) para ejecutar su modelo de visión. Marker lo busca en este orden: primero la ruta indicada en la variable \texttt{LLAMA\_CPP\_BINARY} y, si no existe, un \cmd{llama-server} en el \texttt{PATH}. Hay dos formas de tenerlo:
	
	\begin{lstlisting}
	# Opción 1: binario compilado por usted (por ejemplo en ~/llama.cpp)
	echo "$LLAMA_CPP_BINARY"                 # ¿está definida la variable?
	"$LLAMA_CPP_BINARY" --version            # ¿responde el programa?
	
	# Opción 2: paquete de Ubuntu 26.04 (instala /usr/bin/llama-server)
	sudo apt install llama.cpp-tools
	llama-server --version
\end{lstlisting}
	
	\begin{nota}
		Si \texttt{LLAMA\_\allowbreak{}CPP\_\allowbreak{}BINARY} apunta a un archivo existente, Marker usará ese binario aunque haya otro en el \texttt{PATH}. Si cambia de opción, elimine o corrija la línea \texttt{export LLAMA\_\allowbreak{}CPP\_\allowbreak{}BINARY=\allowbreak{}...} de \arch{\tl/.bashrc}. La sección~\ref{sec:marker-arquitectura} explica para qué sirve este servidor.
	\end{nota}
	
	\necesitaespacio{10\baselineskip}
	\subsection{Comprobar que todo está instalado}
	
	\necesitaespacio{6\baselineskip}
	\subsubsection{Ver la versión de cada herramienta}
	
	\begin{lstlisting}
  pdftoppm -v            # Poppler (todas sus utilidades aceptan -v)
  img2pdf --version
  magick -version        # ImageMagick usa un solo guion: -version
  potrace --version
  pandoc --version
  qpdf --version
  ocrmypdf --version
  tesseract --version
  gs --version  # Ghostscript: solo el número; gs -h muestra además los dispositivos
\end{lstlisting}
	
	\necesitaespacio{26\baselineskip}
	\subsubsection{Localizar un ejecutable: \texttt{which}, \texttt{command -v} y \texttt{type -a}}
	
	\necesitaespacio{22\baselineskip}
	\begin{opciones}[Comando]
		\cmd{which programa} & Busca el programa en las carpetas del \texttt{PATH} y muestra la primera coincidencia. No conoce alias ni funciones de Bash.\\*
		\cmd{command -v programa} & Orden interna de Bash (estándar POSIX): muestra lo que se ejecutaría de verdad, incluidos alias y funciones. Útil en scripts: \cmd{command -v gs >/\allowbreak{}dev/\allowbreak{}null |\allowbreak{}| echo "Falta Ghostscript"}.\\
		\cmd{type -a programa} & Muestra \emph{todas} las coincidencias en orden de prioridad. Ideal para detectar duplicados (\cmd{apt} frente a \cmd{pipx} o entornos virtuales).\\
		\cmd{dpkg -L paquete} & Lista los archivos que instaló un paquete de \cmd{apt}: \cmd{dpkg -L poppler-\allowbreak{}utils | grep /\allowbreak{}usr/\allowbreak{}bin/}.\\*
		\cmd{apt policy paquete} & Versión instalada y versión disponible en los repositorios.\\
	\end{opciones}
	
	\necesitaespacio{6\baselineskip}
	\subsubsection{Script de comprobación general}
	
	\begin{lstlisting}
	for c in pdfinfo pdftoppm pdftotext img2pdf magick potrace mkbitmap pandoc \
	qpdf ocrmypdf tesseract gs; do
	if command -v "$c" >/dev/null; then
	printf '%-10s OK  -> %s\n' "$c" "$(command -v "$c")"
	else
	printf '%-10s FALTA\n' "$c"
	fi
	done
	tesseract --list-langs        # idiomas de OCR disponibles (debe aparecer spa)
\end{lstlisting}
	
	\necesitaespacio{6\baselineskip}
	\subsubsection{Cómo pedir ayuda a cada herramienta}
	
	\begin{lstlisting}
	pdftoppm -h            # Poppler: -h o -help
	qpdf --help            # QPDF: ayuda general; qpdf --help=all muestra todo
	man ocrmypdf           # manual completo (salir con la tecla q)
	magick -help           # ImageMagick: lista de opciones
	# resumen; el manual completo está en https://pandoc.org/MANUAL.html
	pandoc --help
	marker_single --help   # con el entorno marker-env activo
\end{lstlisting}
	
	% =====================================================================
	\clearpage
	\section{Convenciones utilizadas en la guía}\label{sec:convenciones}
	% =====================================================================
	
	\necesitaespacio{24\baselineskip}
	\subsection{Notación de sintaxis}
	
	\necesitaespacio{20\baselineskip}
	\begin{opciones}[Notación]
		\oblig{entrada} & Elemento obligatorio que debe sustituir por su valor (sin escribir los signos \texttt{< >}).\\*
		\opcional{opciones} & Elemento opcional (no escriba los corchetes).\\
		\texttt{a | b} & Elija una de las alternativas.\\
		\texttt{...} & El elemento anterior puede repetirse (por ejemplo, varios archivos).\\
		\texttt{NOMBRE} & Su usuario de Windows en rutas \arch{/\allowbreak{}mnt/\allowbreak{}c/\allowbreak{}Users/\allowbreak{}NOMBRE/\allowbreak{}...}\\
		\texttt{USUARIO} & Su usuario de Linux (\cmd{whoami}).\\
		\texttt{\bs} al final de línea & El comando continúa en la línea siguiente; puede escribirse todo en una sola línea omitiendo la barra.\\*
		\texttt{\# comentario} & Texto explicativo: Bash ignora lo que va detrás de \texttt{\#}. Puede copiarse sin problema.\\
	\end{opciones}
	
	Ejemplo de sintaxis general y su lectura:
	
	\begin{lstlisting}[style=sintaxis]
  pdftoppm [opciones] <archivo.pdf> <prefijo_de_salida>
\end{lstlisting}
	
	Se lee así: el programa \cmd{pdftoppm}, cero o más opciones, un PDF obligatorio y un prefijo obligatorio para los nombres de las imágenes resultantes.
	
	\necesitaespacio{20\baselineskip}
	\subsection{Convenciones tipográficas}
	
	\begingroup\small
	\necesitaespacio{16\baselineskip}
	\begin{longtable}{L{4.3cm}L{10.76cm}}
		\toprule\rowcolor{tablacab}\cab{Aspecto} & \cab{Significado}\\\midrule
		\endfirsthead
		\bottomrule
		\endfoot
		\cmd{magick} & Programas y órdenes que se escriben en la terminal.\\*
		\opt{-{}-output\_dir} & Opciones (parámetros) de un programa.\\
		\arch{\tl/\allowbreak{}trabajo/\allowbreak{}doc.pdf} & Archivos, carpetas y rutas.\\
		\val{300}, \val{A4} & Valores literales que se pasan a una opción.\\
		\textcolor{codigofondo!50!black}{\rule{0.5cm}{0.25cm}} Bloque gris & Comandos listos para copiar (sin el indicador \texttt{\$}).\\
		Bloque con barra azul & Sintaxis general: no se copia tal cual, se adapta.\\*
		Bloque blanco con marco & Salida que muestra la terminal o contenido de un archivo.\\
	\end{longtable}
	\endgroup
	
	\begin{nota}[Copiar comandos]
		Los bloques de código no incluyen el indicador \texttt{\$} del sistema para que puedan copiarse directamente. Si una línea larga se corta en el PDF, aparece una flecha \textcolor{codigocoment}{$\hookrightarrow$} al comienzo de la continuación: no forma parte del comando.
	\end{nota}
	
	\necesitaespacio{22\baselineskip}
	\subsection{Cajas de aviso}
	
	\begingroup\small
	\necesitaespacio{18\baselineskip}
	\begin{longtable}{L{3.2cm}L{11.86cm}}
		\toprule\rowcolor{tablacab}\cab{Caja} & \cab{Uso}\\\midrule
		\endfirsthead
		\bottomrule
		\endfoot
		\textbf{\textcolor{colnota}{Nota}} & Información complementaria o aclaraciones.\\*
		\textbf{\textcolor{colconsejo}{Consejo}} & Forma más cómoda, rápida o segura de hacer algo.\\
		\textbf{\textcolor{coladvert}{Advertencia}} & Riesgo de resultados incorrectos, pérdida de calidad o sobrescritura.\\
		\textbf{\textcolor{colimport}{Importante}} & Punto crítico que, si se ignora, provoca errores o pérdida de datos.\\
		\textbf{\textcolor{collimit}{Limitación}} & Algo que la herramienta no puede hacer o hace con restricciones.\\*
		\textbf{\textcolor{colejemplo}{Ejemplo}} & Caso práctico completo.\\
	\end{longtable}
	\endgroup
	
	\necesitaespacio{12\baselineskip}
	\subsection{Archivos y carpetas de ejemplo}
	
	Los ejemplos usan nombres genéricos: \arch{doc.pdf} (PDF digital), \arch{escaneado.pdf} (PDF formado por imágenes), \arch{foto.jpg}, \arch{logo.png}, \arch{doc.md}. La carpeta de trabajo de Linux es \arch{\tl/trabajo}, y la de Windows, \arch{/\allowbreak{}mnt/\allowbreak{}c/\allowbreak{}Users/\allowbreak{}NOMBRE/\allowbreak{}Downloads}.
	
	\necesitaespacio{8\baselineskip}
	\subsection{Numeración de páginas}
	
	\begin{importante}[Páginas desde 0 o desde 1]
		Poppler, QPDF, Ghostscript, OCRmyPDF e ImageMagick numeran de forma distinta. Poppler, QPDF, Ghostscript y OCRmyPDF cuentan las páginas desde~1. \textbf{Marker} (\opt{-{}-page\_range}) e \textbf{ImageMagick} (\arch{doc.pdf[0]}) cuentan desde~0. Así, la «página 1» del documento es \val{1} en \cmd{pdftoppm -f 1} pero \val{0} en \cmd{marker\_\allowbreak{}single -{}-page\_\allowbreak{}range 0}.
	\end{importante}
	
	% =====================================================================
	\clearpage
	\section{Bash: comodines y procesamiento por lotes}\label{sec:bash}
	% =====================================================================
	
	\necesitaespacio{24\baselineskip}
	\subsection{Comodines}\label{sec:comodines}
	
	Un comodín es un patrón que Bash sustituye por la lista de archivos que coinciden con él \emph{antes} de ejecutar el programa. El programa nunca ve el asterisco: recibe la lista ya expandida.
	
	\necesitaespacio{20\baselineskip}
	\begin{opciones}[Patrón]
		\texttt{*} & Cualquier secuencia de caracteres (incluida la vacía): \arch{*.png}, \arch{archivo\_*.pdf}, \arch{cap*}.\\*
		\texttt{?} & Exactamente un carácter: \arch{pag-?.png} coincide con \arch{pag-1.png} pero no con \arch{pag-10.png}.\\
		\texttt{[...]} & Un carácter del conjunto o rango: \arch{pag-[0-9].png}, \arch{*.[jJ][pP][gG]} (jpg en mayúsculas o minúsculas).\\
		\texttt{[!...]} & Un carácter que \emph{no} esté en el conjunto: \arch{[!\_]*.pdf}.\\
		\texttt{\{a,b\}} & No es un comodín sino una expansión de llaves: genera los textos aunque los archivos no existan. \arch{*.\{jpg,png\}} equivale a \arch{*.jpg *.png}.\\*
		\texttt{**} & Cualquier nivel de subcarpetas, solo tras \cmd{shopt -s globstar}: \arch{**/*.pdf}.\\
	\end{opciones}
	
	\begin{lstlisting}
	ls *.png                  # ver qué archivos coinciden ANTES de usar el patrón
	echo archivo_*.pdf        # echo muestra la expansión sin ejecutar nada
	ls -1 *.{jpg,png}         # varias extensiones
\end{lstlisting}
	
	\necesitaespacio{6\baselineskip}
	\subsubsection{Orden de los archivos}\label{sec:orden}
	
	Bash ordena la expansión alfabéticamente, carácter a carácter, según la configuración regional. Por eso \arch{pag10.png} queda \emph{antes} que \arch{pag2.png}: se compara el carácter \texttt{1} con el \texttt{2}. Para procesos donde el orden importa (unir páginas, crear un PDF con img2pdf), hay dos soluciones:
	
	\begin{enumerate}
		\item \textbf{Nombres con ceros a la izquierda} (\arch{pag-01.png}, \arch{pag-02.png}\ldots, \arch{pag-10.png}). Poppler los genera así automáticamente cuando el documento tiene diez o más páginas.
		\item \textbf{Orden natural} con \cmd{sort -V} (orden de «versiones», que entiende los números):
	\end{enumerate}
	
	\begin{lstlisting}
	ls -v *.png                                  # ver el orden natural
	mapfile -t imgs < <(printf '%s\n' *.png | sort -V)   # guardar la lista ordenada
	printf '%s\n' "${imgs[@]}"                   # revisar
	img2pdf "${imgs[@]}" -o documento.pdf        # usar la lista
\end{lstlisting}
	
	\necesitaespacio{6\baselineskip}
	\subsubsection{Comportamientos que conviene conocer}
	
	\begin{itemize}
		\item \textbf{Si nada coincide}, Bash deja el patrón tal cual y el programa recibe literalmente \arch{*.png}, lo que provoca errores del tipo «No such file». Con \cmd{shopt -s nullglob} el patrón sin coincidencias se convierte en una lista vacía (útil en scripts).
		\item \textbf{Mayúsculas}: \arch{*.jpg} no coincide con \arch{FOTO.JPG}. Use \arch{*.[jJ][pP][gG]} o active \cmd{shopt -s nocaseglob}.
		\item \textbf{Archivos ocultos} (los que empiezan por punto) no coinciden con \arch{*}.
		\item \textbf{Límite de longitud}: con decenas de miles de archivos, la lista expandida puede superar la longitud máxima de una orden (\texttt{Argument list too long}). Use bucles, \cmd{find} o \cmd{xargs}.
	\end{itemize}
	
	\begin{advertencia}[Comodines y órdenes destructivas]
		Un comodín mal escrito con \cmd{rm}, \cmd{mv} o con opciones que sobrescriben (\opt{-{}-replace-\allowbreak{}input} de QPDF, \cmd{magick mogrify}) puede afectar a muchos archivos a la vez y sin papelera. Ejecute antes el mismo patrón con \cmd{ls} o \cmd{echo}, y trabaje sobre copias.
	\end{advertencia}
	
	\necesitaespacio{8\baselineskip}
	\subsection{Bucles \texttt{for}}
	
	\begin{lstlisting}[style=sintaxis]
	for VARIABLE in LISTA; do
	ORDEN "$VARIABLE"
	done
\end{lstlisting}
	
	\begin{lstlisting}
	# Extraer el texto de cada PDF de la carpeta actual
	for f in *.pdf; do
	pdftotext -layout "$f" "${f%.pdf}.txt"
	done
\end{lstlisting}
	
	Las comillas en \texttt{"\$f"} son imprescindibles: sin ellas, un nombre con espacios se partiría en varios argumentos.
	
	\necesitaespacio{18\baselineskip}
	\subsubsection{Manipular nombres de archivo}
	
	\necesitaespacio{14\baselineskip}
	\begin{opciones}[Expresión]
		\texttt{\$\{f\%.pdf\}} & Quita el sufijo \texttt{.pdf}: \arch{informe.pdf} $\rightarrow$ \arch{informe}.\\*
		\texttt{\$\{f\%.*\}} & Quita la última extensión, sea cual sea: \arch{foto.jpeg} $\rightarrow$ \arch{foto}.\\
		\texttt{\$\{f\#\#*/\}} & Quita la ruta y deja el nombre: \arch{/a/b/doc.pdf} $\rightarrow$ \arch{doc.pdf} (como \cmd{basename}).\\
		\texttt{\$\{f\%/*\}} & Deja solo la carpeta: \arch{/a/b/doc.pdf} $\rightarrow$ \arch{/a/b} (como \cmd{dirname}).\\*
		\texttt{\$\{f\#\#*.\}} & Deja solo la extensión: \arch{doc.pdf} $\rightarrow$ \arch{pdf}.\\
	\end{opciones}
	
	\necesitaespacio{6\baselineskip}
	\subsubsection{Plantilla robusta para procesos por lotes}\label{sec:plantilla-lote}
	
	\begin{lstlisting}
	#!/usr/bin/env bash
	# lote.sh - procesa todos los PDF de ENTRADA y deja los resultados en SALIDA
	set -u                           # error si se usa una variable no definida
	shopt -s nullglob                # si no hay PDF, el bucle no se ejecuta
	ENTRADA="$HOME/trabajo/entrada"
	SALIDA="$HOME/trabajo/salida"
	mkdir -p "$SALIDA"
	for f in "$ENTRADA"/*.pdf; do
	nombre="${f##*/}"                      # doc.pdf
	base="${nombre%.pdf}"                  # doc
	destino="$SALIDA/$base.txt"
	[ -e "$destino" ] && { echo "Ya existe, se omite: $destino"; continue; }
	echo "Procesando: $nombre"
	if ! pdftotext -layout "$f" "$destino"; then
	echo "ERROR: $nombre" >> "$SALIDA/errores.log"
	fi
	done
\end{lstlisting}
	
	\begin{lstlisting}
	nano ~/trabajo/lote.sh        # crear el archivo y pegar el contenido
	chmod +x ~/trabajo/lote.sh    # darle permiso de ejecución
	~/trabajo/lote.sh             # ejecutarlo
\end{lstlisting}
	
	\begin{consejo}[Ensayo en seco]
		Antes de ejecutar un bucle que modifica archivos, anteponga \cmd{echo} a la orden: el bucle imprimirá los comandos que \emph{ejecutaría}, sin hacer nada. Cuando la lista sea correcta, quite el \cmd{echo}.
	\end{consejo}
	
	\necesitaespacio{8\baselineskip}
	\subsection{Archivos en subcarpetas: \texttt{find}}
	
	\begin{lstlisting}
	# Listar todos los PDF bajo ~/trabajo, en cualquier subcarpeta
	find ~/trabajo -type f -name '*.pdf'
	
	# Procesarlos uno a uno, con nombres que contengan espacios o saltos de línea
	find ~/trabajo -type f -name '*.pdf' -print0 |
	while IFS= read -r -d '' f; do
	pdfinfo "$f" | grep -E '^Pages:'
	done
\end{lstlisting}
	
	Las comillas simples en \texttt{'*.pdf'} impiden que Bash expanda el patrón: así es \cmd{find} quien lo interpreta en cada subcarpeta. \opt{-print0} separa los nombres con un carácter nulo, el único que no puede aparecer en un nombre de archivo.
	
	\necesitaespacio{8\baselineskip}
	\subsection{Procesamiento en paralelo con \texttt{xargs}}
	
	\begin{lstlisting}
	# Hasta 4 procesos simultáneos; pdftotext crea doc.txt junto a cada doc.pdf
	find ~/trabajo -name '*.pdf' -print0 | xargs -0 -n 1 -P 4 pdftotext -layout
\end{lstlisting}
	
	\begin{advertencia}
		No combine paralelismo externo (\opt{-P} de \cmd{xargs}) con herramientas que ya usan varios núcleos por sí mismas (OCRmyPDF con \opt{-{}-jobs}, Marker con \opt{-{}-workers}): multiplicaría los procesos y podría agotar la memoria de WSL2.
	\end{advertencia}
	
	\necesitaespacio{8\baselineskip}
	\subsection{Ejemplos de lotes habituales}
	
	\begin{lstlisting}
	# Varias imágenes: reducir todas las JPG a 50 % en otra carpeta (originales
	# intactos)
	mkdir -p reducidas && magick mogrify -path reducidas -resize 50% *.jpg
	
	# Varios PDF: rasterizar cada uno en su propia carpeta
	for f in *.pdf; do
	d="${f%.pdf}_png"; mkdir -p "$d"
	pdftoppm -png -r 150 "$f" "$d/pag"
	done
	
	# Varios Markdown: convertir todos a DOCX
	for f in *.md; do pandoc "$f" -o "${f%.md}.docx"; done
	
	# Varias conversiones encadenadas: OCR y después texto plano
	mkdir -p ocr texto
	for f in escaneados/*.pdf; do
	b="${f##*/}"; b="${b%.pdf}"
	ocrmypdf -l spa "$f" "ocr/$b.pdf" && pdftotext "ocr/$b.pdf" "texto/$b.txt"
	done
\end{lstlisting}
	
	El operador \texttt{\&\&} ejecuta la segunda orden solo si la primera terminó bien (código de salida 0). Para encadenar sin importar el resultado se usaría \texttt{;}.
	% =====================================================================
	\clearpage
	\section{Marker}\label{sec:marker}
	% =====================================================================
	
	\necesitaespacio{8\baselineskip}
	\subsection{Descripción general}
	
	\textbf{Qué es.} Marker es un conversor de documentos desarrollado por Datalab que transforma PDF (y otros formatos) en \textbf{Markdown}, \textbf{JSON}, \textbf{HTML} o \textbf{fragmentos} («chunks») para sistemas de búsqueda. A diferencia de \cmd{pdftotext}, que devuelve texto plano, Marker reconstruye la \emph{estructura}: jerarquía de títulos, párrafos en orden de lectura, listas, tablas, bloques de código, ecuaciones en LaTeX, notas al pie e imágenes. Su código se publica con licencia Apache 2.0; los pesos de sus modelos usan una licencia Open Rail-M modificada, gratuita para investigación, uso personal y empresas con menos de 5 millones de dólares de financiación o ingresos.
	
	\textbf{Qué problema resuelve.} Convertir documentos visualmente complejos (artículos científicos a dos columnas, libros, informes con tablas) en texto estructurado y editable, apto para Pandoc, para notas, para procesadores de texto o para sistemas de inteligencia artificial.
	
	\textbf{Cuándo conviene usarlo:}
	\begin{itemize}
		\item Para pasar a Markdown un PDF conservando títulos, tablas y fórmulas.
		\item Para documentos con varias columnas, encabezados y pies de página que deben eliminarse.
		\item Para extraer las tablas de un documento (\opt{-{}-converter\_\allowbreak{}cls} con el conversor de tablas).
	\end{itemize}
	
	\textbf{Cuándo no conviene:}
	\begin{itemize}
		\item Si solo necesita el texto sin estructura: \cmd{pdftotext} es instantáneo.
		\item Si necesita un PDF con texto seleccionable (no un Markdown): use OCRmyPDF.
		\item Si el PDF debe conservar su aspecto exacto: Marker reconstruye el contenido, no la maquetación.
		\item En equipos con poca memoria y documentos escaneados largos, el reconocimiento en CPU puede ser muy lento.
	\end{itemize}
	
	\textbf{Complementos naturales:} Poppler (\cmd{pdffonts}, \cmd{pdfinfo}) para saber antes si el PDF tiene texto; QPDF para dividir documentos enormes; OCRmyPDF para preparar escaneados en equipos modestos; Pandoc para convertir el Markdown resultante en DOCX o PDF.
	
	\necesitaespacio{8\baselineskip}
	\subsection{Cómo funciona Marker 2}\label{sec:marker-arquitectura}
	
	Marker~2 combina modelos pequeños que corren en la CPU con un modelo de visión y lenguaje (VLM) llamado \textbf{surya}, de unos 650 millones de parámetros, que se ejecuta en un \textbf{servidor de inferencia local}. Para cada documento:
	
	\flujo{PDF o imagen,pdftext,texto incrustado,detector de maquetación,bloques (títulos{,} tablas{,} figuras\ldots),surya (solo donde hace falta),OCR{,} ecuaciones{,} tablas difíciles,renderizador,Markdown / JSON / HTML / chunks}
	
	\begin{enumerate}
		\item Extrae el texto incrustado del PDF con \texttt{pdftext}, en orden de lectura.
		\item Detecta la maquetación de cada página (en modo \val{fast}, con un detector ligero; en modo \val{balanced}, con el VLM).
		\item Decide página a página si el texto incrustado es utilizable. Las páginas escaneadas o con texto corrupto se reconocen con el VLM.
		\item Reconoce ecuaciones con el VLM (el texto incrustado no puede representar matemáticas).
		\item Reconstruye las tablas a partir del texto, y recurre al VLM cuando la reconstrucción es poco fiable.
		\item Opcionalmente, mejora el resultado con un LLM externo (\opt{-{}-use\_llm}).
	\end{enumerate}
	
	\necesitaespacio{18\baselineskip}
	\subsubsection{Backend del servidor de inferencia}
	
	El servidor se arranca \emph{automáticamente} la primera vez que se necesita, y se detiene al terminar:
	
	\begingroup\small
	\necesitaespacio{14\baselineskip}
	\begin{longtable}{L{3.3cm}L{5.4cm}L{6.36cm}}
		\toprule\rowcolor{tablacab}\cab{Equipo} & \cab{Backend} & \cab{Requisitos}\\\midrule
		\endfirsthead
		\bottomrule
		\endfoot
		Sin GPU NVIDIA (CPU) & llama.cpp (\cmd{llama-server}) & \cmd{llama-server} en el \texttt{PATH} o indicado en \texttt{LLAMA\_\allowbreak{}CPP\_\allowbreak{}BINARY} (sección~\ref{sec:marker-instalacion}).\\
		Con GPU NVIDIA & vLLM dentro de Docker & Docker y NVIDIA Container Toolkit funcionando dentro de WSL2. En WSL2 el controlador NVIDIA se instala en Windows, no en Ubuntu.\\
	\end{longtable}
	\endgroup
	
	La detección es automática: si Marker encuentra una GPU NVIDIA, elige vLLM; si no, llama.cpp. Para forzar uno u otro: \texttt{export SURYA\_\allowbreak{}INFERENCE\_\allowbreak{}BACKEND=\allowbreak{}llamacpp} (o \val{vllm}).
	
	\necesitaespacio{25\baselineskip}
	\subsubsection{Modos de conversión}\label{sec:marker-modos}
	
	\necesitaespacio{21\baselineskip}
	\begin{opciones}[Modo]
		\opt{-{}-mode balanced} & Predeterminado con GPU. Usa el VLM para la maquetación, reconoce las fórmulas en línea y vuelve a reconocer la página entera si parte de su texto es defectuoso. Máxima calidad, mayor coste.\\*
		\opt{-{}-mode fast} & Predeterminado en CPU. Detector de maquetación ligero, texto incrustado y uso mínimo del VLM: ecuaciones, reparación de bloques concretos defectuosos y reconocimiento de página completa solo en páginas escaneadas. Un documento digital limpio y sin ecuaciones \emph{nunca} arranca el servidor.\\*
		\opt{-{}-disable\_ocr} & No llama nunca al VLM (válido en ambos modos): solo texto incrustado. Muy rápido y sin servidor, pero se pierden las ecuaciones y las páginas escaneadas.\\
	\end{opciones}
	
	Como orientación, el propio proyecto publica estas puntuaciones en la prueba independiente \emph{olmOCR-bench} (porcentaje de pruebas superadas; más alto es mejor):
	
	\begingroup\small
	\necesitaespacio{10\baselineskip}
	\begin{longtable}{L{5.06cm}C{2.4cm}C{2.4cm}C{2.4cm}}
		\toprule\rowcolor{tablacab}\cab{Categoría} & \cab{balanced} & \cab{fast} & \cab{sin OCR}\\\midrule
		\endfirsthead
		\bottomrule
		\endfoot
		Matemáticas (arXiv) & 83,9 & 23,4 & 0,0\\*
		Tablas & 73,4 & 69,0 & 46,1\\
		Varias columnas & 76,6 & 76,0 & 67,0\\
		Escaneos antiguos & 43,2 & 43,2 & 14,3\\*
		\textbf{Total} & \textbf{76,0} & \textbf{66,6} & \textbf{43,6}\\
	\end{longtable}
	\endgroup
	
	La lectura práctica: para documentos con muchas fórmulas, el modo \val{fast} solo las recupera bien si se fuerza el reconocimiento (\opt{-{}-force\_ocr} o \opt{-{}-ocr\_\allowbreak{}inline\_\allowbreak{}math true}); \opt{-{}-disable\_ocr} es adecuado para documentos digitales de texto corrido.
	
	\necesitaespacio{20\baselineskip}
	\subsubsection{Componentes que se descargan y dónde se guardan}
	
	\begingroup\small
	\necesitaespacio{16\baselineskip}
	\begin{longtable}{L{4.5cm}L{5.2cm}L{5.36cm}}
		\toprule\rowcolor{tablacab}\cab{Componente} & \cab{Para qué sirve} & \cab{Ubicación predeterminada}\\\midrule
		\endfirsthead
		\bottomrule
		\endfoot
		Modelos pequeños (maquetación, detección de texto, detección de errores de OCR) & Análisis en CPU de cada página & \arch{\tl/\allowbreak{}.cache/\allowbreak{}datalab/\allowbreak{}models} (variable \texttt{MODEL\_\allowbreak{}CACHE\_\allowbreak{}DIR})\\*
		Modelo surya en formato GGUF (\texttt{surya-2.gguf} y \texttt{surya-\allowbreak{}2-\allowbreak{}mmproj.gguf}) & VLM para llama.cpp: OCR, ecuaciones, tablas difíciles & Caché de Hugging Face: \arch{\tl/\allowbreak{}.cache/\allowbreak{}huggingface} (variable \texttt{HF\_HOME})\\*
		Imagen Docker de vLLM y modelo completo & Solo con GPU NVIDIA & Almacenamiento de Docker y caché de Hugging Face\\
	\end{longtable}
	\endgroup
	
	La primera conversión tarda más porque descarga los modelos. Para ver cuánto ocupan en su equipo:
	
	\begin{lstlisting}
			   du -sh ~/.cache/datalab ~/.cache/huggingface 2>/dev/null
\end{lstlisting}
	
	\necesitaespacio{6\baselineskip}
	\subsubsection{CPU, GPU y memoria en WSL2}
	
	\begin{itemize}
		\item \textbf{Sin GPU}: Marker funciona en CPU con el modo \val{fast}. El reconocimiento con el VLM es la parte más lenta; los documentos digitales sin ecuaciones son rápidos.
		\item \textbf{Con GPU NVIDIA}: compruebe que WSL2 la ve con \cmd{nvidia-smi}. Para el backend vLLM necesita Docker y NVIDIA Container Toolkit dentro de Ubuntu; si no los tiene, fuerce \texttt{SURYA\_\allowbreak{}INFERENCE\_\allowbreak{}BACKEND=\allowbreak{}llamacpp}.
		\item \textbf{Memoria}: el proyecto no fija un requisito mínimo. El consumo depende del número de procesos de conversión (\opt{-{}-workers}), del servidor de inferencia y del tamaño de las páginas. Si falta memoria: reduzca \opt{-{}-workers}, divida los PDF largos (sección~\ref{sec:qpdf-paginas}) y revise el límite de memoria de WSL2 (sección~\ref{sec:wsl-memoria}).
	\end{itemize}
	
	\necesitaespacio{8\baselineskip}
	\subsection{Instalación y comprobación}
	
	La instalación en \arch{\tl/marker-env} y la configuración de \cmd{llama-server} se describen en la sección~\ref{sec:marker-instalacion}. Comprobación rápida:
	
	\begin{lstlisting}
  source ~/marker-env/bin/activate
  pip show marker-pdf | head -2      # nombre y versión
  marker_single --help               # opciones comunes
  marker_single config --help  # TODAS las opciones de configuración (lista larga)
\end{lstlisting}
	
	\begin{nota}[Ejecutables de Marker 2]
		El paquete instala cuatro órdenes: \cmd{marker\_single} (un archivo), \cmd{marker} (una carpeta), \cmd{marker\_server} (API web local) y \cmd{marker\_gui} (interfaz web interactiva). La orden \cmd{marker\_\allowbreak{}chunk\_\allowbreak{}convert} de versiones anteriores ya no existe.
	\end{nota}
	
	\necesitaespacio{8\baselineskip}
	\subsection{\texttt{marker\_single}: convertir un archivo}\label{sec:marker-single}
	
	\begin{lstlisting}[style=sintaxis]
  marker_single [opciones] <archivo>
\end{lstlisting}
	
	Convierte un PDF o una imagen (y, con \texttt{marker-\allowbreak{}pdf[full]}, DOCX, PPTX, XLSX, HTML o EPUB). El resultado se guarda en una carpeta con el nombre del archivo dentro de \opt{-{}-output\_dir}.
	
	\flujoh{informe.pdf,marker\_single,salida/informe/informe.md}
	
	\necesitaespacio{14\baselineskip}
	\subsubsection{Opciones principales}
	
	\necesitaespacio{10\baselineskip}
	\begin{opciones}
		\opt{-{}-output\_\allowbreak{}dir RUTA} & Carpeta de salida; se crea si no existe. \textbf{Indíquela siempre}: el valor predeterminado es una carpeta \arch{conversion\_\allowbreak{}results} dentro de la instalación de Python del entorno, difícil de encontrar.\\*
		\opt{-{}-output\_\allowbreak{}format F} & \val{markdown} (predeterminado), \val{json} (árbol de bloques con coordenadas), \val{html} o \val{chunks} (lista plana de bloques para sistemas de búsqueda).\\
		\opt{-{}-page\_\allowbreak{}range "R"} & Páginas a convertir, \textbf{numeradas desde 0}: \val{"0,5-10,20"} convierte la primera página, de la sexta a la undécima y la vigésimo primera.\\
		\opt{-{}-mode balanced|\allowbreak{}fast} & Modo de conversión (sección~\ref{sec:marker-modos}). Sin la opción, depende del equipo: \val{balanced} con GPU, \val{fast} en CPU.\\
		\opt{-{}-disable\_ocr} & Sin VLM: solo texto incrustado. No arranca el servidor.\\
		\opt{-{}-force\_ocr} & Reconoce todo el documento con el VLM, aunque tenga texto. Útil con texto incrustado corrupto (caracteres extraños) y para máxima calidad matemática. Mucho más lento en CPU.\\
		\opt{-{}-strip\_\allowbreak{}existing\_\allowbreak{}ocr} & Elimina la capa OCR que ya tuviera el PDF y vuelve a reconocer con surya.\\
		\opt{-{}-ocr\_\allowbreak{}inline\_\allowbreak{}math true|\allowbreak{}false} & Reconoce las fórmulas en línea de los párrafos. Por defecto es «automático» (activado solo en \val{balanced}). \textbf{Requiere el valor} \val{true} o \val{false}: escribir solo \opt{-{}-ocr\_\allowbreak{}inline\_\allowbreak{}math} provoca un error.\\
		\opt{-{}-disable\_\allowbreak{}image\_\allowbreak{}extraction} & No guarda las imágenes. Con \opt{-{}-use\_llm}, las sustituye por una descripción textual.\\
		\opt{-{}-paginate\_\allowbreak{}output} & Separa las páginas en la salida con su número seguido de una línea de 48 guiones.\\
		\opt{-{}-keep\_\allowbreak{}pageheader\_\allowbreak{}in\_\allowbreak{}output} & Conserva los encabezados de página (por defecto se eliminan). Su pareja es \opt{-{}-keep\_\allowbreak{}pagefooter\_\allowbreak{}in\_\allowbreak{}output}.\\
		\opt{-{}-use\_llm} & Modo híbrido con un LLM externo (sección~\ref{sec:marker-llm}).\\
		\opt{-{}-converter\_\allowbreak{}cls CLASE} & \val{marker.converters.pdf.PdfConverter} (predeterminado) o \val{marker.converters.table.TableConverter} (solo tablas).\\
		\opt{-{}-config\_\allowbreak{}json ARCHIVO} & Archivo JSON con opciones de configuración (mismos nombres, sin guiones).\\
		\opt{-{}-processors LISTA} & Sustituye la cadena de procesadores (rutas completas de módulos, separadas por comas). Uso avanzado.\\
		\opt{-{}-disable\_\allowbreak{}multiprocessing} & Extrae el texto en un solo proceso (menos memoria, más lento).\\*
		\opt{-{}-debug}, \opt{-d} & Modo de depuración: guarda imágenes de cada página con la maquetación detectada e información adicional.\\
	\end{opciones}
	
	\necesitaespacio{6\baselineskip}
	\subsubsection{Ejemplos}
	
	\begin{lstlisting}
	source ~/marker-env/bin/activate
	cd /mnt/c/Users/NOMBRE/Downloads
	
	# 1) Prueba inicial con las dos primeras páginas (descarga modelos la primera vez)
	marker_single "informe.pdf" --output_dir ./md_salida --page_range "0-1"
	
	# 2) Documento completo, modo rápido con fórmulas en línea
	marker_single "informe.pdf" --output_dir ./md_salida --mode fast \
	--ocr_inline_math true
	
	# 3) PDF digital de texto corrido: sin VLM, muy rápido
	marker_single "manual.pdf" --output_dir ./md_salida --disable_ocr
	
	# 4) PDF con texto incrustado corrupto o escaneado: reconocimiento completo
	marker_single "escaneado.pdf" --output_dir ./md_salida --force_ocr
	
	# 5) Otros formatos de salida
	marker_single "informe.pdf" --output_dir ./json --output_format json
	marker_single "informe.pdf" --output_dir ./html --output_format html
	
	# 6) Solo texto, sin imágenes, con separadores de página
	marker_single "informe.pdf" --output_dir ./md_salida --disable_image_extraction \
	--paginate_output
	
	# 7) Extraer solo las tablas
	marker_single "datos.pdf" --output_dir ./tablas --converter_cls \
	marker.converters.table.TableConverter
\end{lstlisting}
	
	\begin{lstlisting}
	# 8) Opciones en un archivo JSON (útil para repetir siempre la misma
	# configuración)
	cat > ~/marker-config.json <<'EOF'
	{
		"paginate_output": true,
		"disable_image_extraction": false,
		"keep_pageheader_in_output": false
	}
	EOF
	marker_single "informe.pdf" --output_dir ./md_salida --config_json \
	~/marker-config.json
\end{lstlisting}
	
	\necesitaespacio{6\baselineskip}
	\subsubsection{Estructura de la salida}
	
	\begin{lstlisting}[style=salida]
	md_salida/
	informe/
	informe.md            # documento convertido (o .json / .html)
	informe_meta.json     # metadatos: índice detectado y estadísticas por página
	*.jpeg                # imágenes extraídas, enlazadas desde el Markdown
\end{lstlisting}
	
	El archivo \arch{\_meta.json} incluye el índice de títulos detectado (con su nivel y página) y, por cada página, el método con el que se obtuvo el texto. Es útil para comprobar qué páginas necesitaron reconocimiento.
	
	\necesitaespacio{8\baselineskip}
	\subsection{\texttt{marker}: convertir una carpeta}\label{sec:marker-lotes}
	
	\begin{lstlisting}[style=sintaxis]
	marker [opciones] <carpeta_de_entrada>
\end{lstlisting}
	
	Acepta las mismas opciones que \cmd{marker\_single} y, además, las de la tabla siguiente. Procesa \textbf{todos los archivos} situados directamente en la carpeta (no entra en subcarpetas), con varios procesos que comparten un único servidor de inferencia.
	
	\necesitaespacio{18\baselineskip}
	\begin{opciones}
		\opt{-{}-workers N} & Número de procesos de conversión. Se ajusta automáticamente; aumentarlo acelera a costa de CPU y memoria, reducirlo evita errores de memoria.\\*
		\opt{-{}-skip\_\allowbreak{}existing} & Omite los archivos que ya tienen resultado en \opt{-{}-output\_dir}. Permite reanudar un lote interrumpido.\\
		\opt{-{}-max\_files N} & Convierte como máximo \val{N} archivos (útil para una prueba).\\
		\opt{-{}-num\_\allowbreak{}chunks N}, \opt{-{}-chunk\_idx I} & Reparte la lista en \val{N} partes y procesa la parte \val{I} (para varias máquinas).\\
		\opt{-{}-max\_\allowbreak{}tasks\_\allowbreak{}per\_\allowbreak{}worker N} & Recicla cada proceso tras \val{N} archivos (predeterminado 200); limita la acumulación de memoria.\\*
		\opt{-{}-debug\_print} & Muestra información de depuración del lote.\\
	\end{opciones}
	
	\begin{lstlisting}
	# Copiar los PDF al disco de Linux (más rápido) y convertir la carpeta
	mkdir -p ~/trabajo/pdf ~/trabajo/md
	cp /mnt/c/Users/NOMBRE/Downloads/*.pdf ~/trabajo/pdf/
	marker ~/trabajo/pdf --output_dir ~/trabajo/md --mode fast
	
	# Reanudar el mismo lote tras una interrupción
	marker ~/trabajo/pdf --output_dir ~/trabajo/md --mode fast --skip_existing
	
	# Lote de documentos digitales sin ecuaciones: sin servidor y en paralelo
	marker ~/trabajo/pdf --output_dir ~/trabajo/md --disable_ocr --workers 4
\end{lstlisting}
	
	\begin{advertencia}[Carpeta de entrada limpia]
		Como \cmd{marker} intenta convertir cada archivo de la carpeta, deje en ella solo los documentos. Archivos como \arch{*:Zone.Identifier} (sección~\ref{sec:diferencias-wsl}), textos o restos de otros procesos generarán errores.
	\end{advertencia}
	
	\necesitaespacio{29\baselineskip}
	\subsection{Variables de entorno útiles}
	
	\necesitaespacio{25\baselineskip}
	\begin{opciones}[Variable]
		\texttt{LLAMA\_\allowbreak{}CPP\_\allowbreak{}BINARY} & Ruta del \cmd{llama-server} que debe usarse.\\*
		\texttt{SURYA\_\allowbreak{}INFERENCE\_\allowbreak{}BACKEND} & \val{llamacpp} o \val{vllm}; sin definir, se elige automáticamente.\\
		\texttt{SURYA\_\allowbreak{}INFERENCE\_\allowbreak{}KEEP\_\allowbreak{}ALIVE} & Con \val{true}, el servidor sigue funcionando al terminar, y las conversiones siguientes no esperan a que arranque. Deténgalo después con \cmd{pkill llama-\allowbreak{}server}.\\
		\texttt{SURYA\_\allowbreak{}INFERENCE\_\allowbreak{}URL} & Usa un servidor ya en marcha (\val{http:/\allowbreak{}/\allowbreak{}host:puerto/\allowbreak{}v1}) en lugar de arrancar uno.\\
		\texttt{SURYA\_\allowbreak{}INFERENCE\_\allowbreak{}PARALLEL} & Número de peticiones simultáneas al servidor; por defecto se calcula solo.\\
		\texttt{TORCH\_DEVICE} & Dispositivo de los modelos pequeños (\val{cpu} o \val{cuda}); no afecta al VLM.\\
		\texttt{MODEL\_\allowbreak{}CACHE\_\allowbreak{}DIR}, \texttt{HF\_HOME} & Cambian las carpetas de caché de modelos (útil si el disco de Linux tiene poco espacio).\\*
		\texttt{VLLM\_GPUS}, \texttt{VLLM\_\allowbreak{}GPU\_\allowbreak{}TYPE} & Solo backend vLLM: GPU que se usan y modelo de GPU (ajusta la memoria reservada).\\
	\end{opciones}
	
	\begin{lstlisting}
		# Mantener el servidor arrancado durante una sesión de trabajo
		export SURYA_INFERENCE_KEEP_ALIVE=true
		marker_single a.pdf --output_dir md
		marker_single b.pdf --output_dir md
		pkill llama-server          # detenerlo al terminar
\end{lstlisting}
	
	\necesitaespacio{20\baselineskip}
	\subsection{Modo híbrido con un LLM}\label{sec:marker-llm}
	
	Con \opt{-{}-use\_llm}, Marker envía fragmentos del documento a un modelo de lenguaje para fusionar tablas que continúan en otra página, mejorar las fórmulas en línea, dar formato a tablas difíciles y extraer valores de formularios. Por defecto usa Gemini de Google; también admite Claude, OpenAI y servicios compatibles, Azure, Vertex, OpenRouter y Ollama (modelos locales).
	
	\necesitaespacio{16\baselineskip}
	\begin{opciones}
		\opt{-{}-use\_llm} & Activa el modo híbrido.\\*
		\opt{-{}-llm\_\allowbreak{}service CLASE} & Servicio, por ejemplo \val{marker.services.ollama.OllamaService} o \val{marker.services.claude.ClaudeService}.\\
		\opt{-{}-gemini\_\allowbreak{}api\_\allowbreak{}key CLAVE} & Clave de Gemini (alternativa: variable \texttt{GOOGLE\_\allowbreak{}API\_\allowbreak{}KEY}). Cada servicio tiene sus propias opciones de clave y modelo: consúltelas con \cmd{marker\_\allowbreak{}single config -{}-help}.\\
		\opt{-{}-redo\_\allowbreak{}inline\_\allowbreak{}math} & Con \opt{-{}-use\_llm}: máxima calidad en matemáticas en línea.\\*
		\opt{-{}-block\_\allowbreak{}correction\_\allowbreak{}prompt "..."} & Con \opt{-{}-use\_llm}: instrucción adicional para corregir o dar formato al resultado.\\
	\end{opciones}
	
	\begin{advertencia}[Privacidad y coste]
		Salvo con Ollama (local), el modo híbrido envía contenido del documento a un servicio externo, que puede tener coste. No lo use con documentos confidenciales sin revisar las condiciones del proveedor.
	\end{advertencia}
	
	\necesitaespacio{8\baselineskip}
	\subsection{Servidor web e interfaz gráfica}
	
	\begin{lstlisting}
		pip install -U streamlit streamlit-ace    # dependencias de la interfaz
		marker_gui                                # abre una interfaz web local
		
		pip install -U uvicorn fastapi python-multipart
		# API local; documentación en http://localhost:8001/docs
		marker_server --port 8001
\end{lstlisting}
	
	La API local solo acepta los parámetros \texttt{page\_range}, \texttt{mode}, \texttt{force\_ocr}, \texttt{paginate\_\allowbreak{}output} y \texttt{output\_format}; el modo híbrido y \opt{-{}-disable\_ocr} solo están disponibles desde la línea de comandos.
	
	\necesitaespacio{22\baselineskip}
	\subsection{PDF digital frente a PDF escaneado en Marker}
	
	\begingroup\small
	\necesitaespacio{18\baselineskip}
	\begin{longtable}{L{3.1cm}L{6.0cm}L{5.96cm}}
		\toprule\rowcolor{tablacab}\cab{Tipo de PDF} & \cab{Cómo lo trata Marker} & \cab{Recomendación}\\\midrule
		\endfirsthead
		\bottomrule
		\endfoot
		Digital (texto seleccionable) & Usa el texto incrustado; el VLM solo para ecuaciones y bloques defectuosos. & \val{fast}; \opt{-{}-disable\_ocr} si no hay fórmulas; \opt{-{}-ocr\_\allowbreak{}inline\_\allowbreak{}math true} si las hay.\\*
		Escaneado (solo imágenes) & Reconoce la página completa con el VLM. & Funciona directamente, pero en CPU es lento. Alternativa ligera: OCRmyPDF y después Marker con \opt{-{}-disable\_ocr} (sección~\ref{sec:flujo2}).\\*
		Escaneado con OCR previo & Usa la capa OCR existente si es de buena calidad; si no, reconoce de nuevo. & Si la capa OCR es mala: \opt{-{}-strip\_\allowbreak{}existing\_\allowbreak{}ocr} o \opt{-{}-force\_ocr}.\\
	\end{longtable}
	\endgroup
	
	Para saber de qué tipo es un PDF: \cmd{pdffonts doc.pdf} (sin fuentes $\Rightarrow$ escaneado) y \cmd{pdftotext -l 1 doc.pdf -} (sin texto $\Rightarrow$ escaneado).
	
	\necesitaespacio{29\baselineskip}
	\subsection{Formatos de entrada y salida}
	
	\necesitaespacio{25\baselineskip}
	\begin{formatos}
		PDF & \compat & \nocompat & Formato principal. Digital, escaneado o mixto.\\*
		Imágenes (PNG, JPEG\ldots) & \compat & \nocompat & Una imagen por conversión; requiere reconocimiento con el VLM.\\
		DOCX, PPTX, XLSX & \condic & \nocompat & Requieren \texttt{marker-\allowbreak{}pdf[full]}.\\
		HTML, EPUB & \condic & \nocompat & Requieren \texttt{marker-\allowbreak{}pdf[full]}.\\
		Markdown & \nocompat & \compat & Salida predeterminada: tablas, ecuaciones \texttt{\$\$...\$\$}, código, enlaces a imágenes.\\
		JSON & \nocompat & \compat & Árbol de bloques con tipo, coordenadas y HTML de cada bloque.\\
		HTML & \nocompat & \compat & Imágenes con \texttt{<img>} y ecuaciones con \texttt{<math>}.\\
		Chunks (JSON plano) & \nocompat & \compat & Bloques de primer nivel de cada página, pensados para sistemas de recuperación (RAG).\\*
		JPEG (imágenes extraídas) & --- & \compat & Imágenes y figuras del documento, salvo con \opt{-{}-disable\_\allowbreak{}image\_\allowbreak{}extraction}.\\
	\end{formatos}
	
	\necesitaespacio{8\baselineskip}
	\subsection{Limitaciones}
	
	\begin{limitacion}
		\begin{itemize}
			\item \textbf{Fidelidad}: reconstruye contenido y estructura, no la maquetación. Columnas, colores, tipografías y posiciones exactas se pierden.
			\item \textbf{Maquetas complejas}: el proyecto reconoce dificultades con tablas anidadas y formularios; \opt{-{}-use\_llm} y \opt{-{}-force\_ocr} los mejoran.
			\item \textbf{Velocidad en CPU}: el reconocimiento con el VLM es la etapa lenta; documentos escaneados largos pueden tardar mucho en WSL2 sin GPU.
			\item \textbf{Memoria y disco}: modelos, caché y servidor ocupan recursos. Varios procesos (\opt{-{}-workers}) multiplican el consumo.
			\item \textbf{Dependencia de programas externos}: sin \cmd{llama-server} (CPU) o sin Docker (GPU), solo funciona \opt{-{}-disable\_ocr}.
			\item \textbf{Matemáticas en modo \val{fast}}: sin \opt{-{}-ocr\_\allowbreak{}inline\_\allowbreak{}math true} o \opt{-{}-force\_ocr}, las fórmulas pueden quedar como texto ilegible.
			\item \textbf{PDF protegidos}: un PDF cifrado con contraseña de apertura debe descifrarse antes (QPDF, sección~\ref{sec:qpdf-cifrado}).
			\item \textbf{Licencia}: los pesos de los modelos tienen restricciones de uso comercial por encima de cierto tamaño de empresa.
		\end{itemize}
	\end{limitacion}
	
	\necesitaespacio{14\baselineskip}
	\subsection{Errores frecuentes}
	
	\necesitaespacio{10\baselineskip}
	\begin{errores}
		\texttt{marker\_\allowbreak{}single: command not found} & El entorno virtual no está activo. & \cmd{source \tl/\allowbreak{}marker-\allowbreak{}env/\allowbreak{}bin/\allowbreak{}activate}.\\*
		\texttt{llama-\allowbreak{}server binary not found} & Marker necesita el VLM y no encuentra \cmd{llama-server}. & Defina \texttt{LLAMA\_\allowbreak{}CPP\_\allowbreak{}BINARY} o instale \texttt{llama.cpp-\allowbreak{}tools}; o use \opt{-{}-disable\_ocr}.\\
		\texttt{docker binary not found} & Marker detectó una GPU NVIDIA y eligió vLLM, pero no hay Docker. & \texttt{export SURYA\_\allowbreak{}INFERENCE\_\allowbreak{}BACKEND=\allowbreak{}llamacpp}, o instale Docker y NVIDIA Container Toolkit.\\
		\texttt{Option '-{}-ocr\_\allowbreak{}inline\_\allowbreak{}math' requires an argument} & La opción no es un interruptor. & Escriba \opt{-{}-ocr\_\allowbreak{}inline\_\allowbreak{}math true}.\\
		\texttt{No such option} & Opción mal escrita o de otra versión. & \cmd{marker\_\allowbreak{}single config -{}-help} lista las opciones de su versión.\\
		No encuentra los resultados & No se indicó \opt{-{}-output\_dir}. & Repita con \opt{-{}-output\_dir}; los anteriores están en \arch{conversion\_\allowbreak{}results} dentro de \arch{\tl/\allowbreak{}marker-\allowbreak{}env/\allowbreak{}lib/\allowbreak{}python3.12/\allowbreak{}site-\allowbreak{}packages}.\\
		Fórmulas o texto ilegibles & Texto incrustado defectuoso o modo \val{fast} sin reconocimiento matemático. & \opt{-{}-force\_ocr} o \opt{-{}-ocr\_\allowbreak{}inline\_\allowbreak{}math true}.\\
		\texttt{Killed} o \texttt{MemoryError} & Memoria insuficiente en WSL2. & Menos \opt{-{}-workers}, dividir el PDF, aumentar la memoria en \arch{.wslconfig}.\\
		Primera ejecución muy lenta & Descarga de modelos y arranque del servidor. & Normal; pruebe primero con \opt{-{}-page\_\allowbreak{}range "0-\allowbreak{}1"}.\\
		Errores en un lote con \cmd{marker} & Archivos que no son documentos en la carpeta. & Deje solo PDF; elimine \arch{*:Zone.Identifier}.\\*
		Error al abrir un PDF cifrado & Contraseña de apertura. & Descifre con \cmd{qpdf -{}-password=\allowbreak{}... -{}-decrypt}.\\
	\end{errores}
	
	Ayuda desde la terminal: \cmd{marker\_\allowbreak{}single -{}-help}, \cmd{marker -{}-help} y la lista completa con \cmd{marker\_\allowbreak{}single config -{}-help}. Versión: \cmd{pip show marker-\allowbreak{}pdf}.
	% =====================================================================
	\clearpage
	\section{Poppler}\label{sec:poppler}
	% =====================================================================
	
	\necesitaespacio{8\baselineskip}
	\subsection{Descripción general}
	
	\textbf{Qué es.} Poppler es una biblioteca de lectura y renderizado de PDF. En Ubuntu, el paquete \texttt{poppler-utils} la acompaña de un \textbf{conjunto de pequeñas utilidades independientes}, cada una con su propio nombre y sus propias opciones. No existe un comando llamado «poppler».
	
	\textbf{Qué problema resuelve.} Consultar y extraer información de un PDF existente: metadatos, fuentes, texto, imágenes incrustadas, páginas convertidas en imágenes o en otros formatos vectoriales, y operaciones sencillas de separar y unir.
	
	\textbf{Cuándo usarlo:} para diagnosticar un PDF antes de procesarlo, extraer texto de PDF digitales, rasterizar páginas con buena calidad, extraer imágenes sin pérdida y para divisiones o uniones rápidas.
	
	\textbf{Cuándo no usarlo:} no hace OCR (use OCRmyPDF), no reconstruye estructura en Markdown (use Marker), no comprime ni optimiza (use Ghostscript o QPDF), no cifra ni repara estructuras dañadas (use QPDF).
	
	\begingroup\small
	\necesitaespacio{10\baselineskip}
	\begin{longtable}{L{3.0cm}L{8.16cm}C{2.0cm}}
		\toprule\rowcolor{tablacab}\cab{Utilidad} & \cab{Función} & \cab{Sección}\\\midrule
		\endfirsthead
		\bottomrule
		\endfoot
		\cmd{pdfinfo} & Información general y metadatos. & \ref{sec:pdfinfo}\\*
		\cmd{pdffonts} & Fuentes usadas y si están incrustadas. & \ref{sec:pdffonts}\\
		\cmd{pdftotext} & Extraer el texto. & \ref{sec:pdftotext}\\
		\cmd{pdftoppm} & Páginas $\rightarrow$ imágenes PNG, JPEG, TIFF o PPM. & \ref{sec:pdftoppm}\\
		\cmd{pdftocairo} & Páginas $\rightarrow$ PNG, JPEG, TIFF, SVG, PDF, PS o EPS. & \ref{sec:pdftocairo}\\
		\cmd{pdfimages} & Extraer las imágenes incrustadas. & \ref{sec:pdfimages}\\
		\cmd{pdftohtml} & PDF $\rightarrow$ HTML o XML. & \ref{sec:pdftohtml}\\
		\cmd{pdftops} & PDF $\rightarrow$ PostScript o EPS. & \ref{sec:pdftops}\\
		\cmd{pdfseparate} & Una página por archivo. & \ref{sec:pdfseparate}\\
		\cmd{pdfunite} & Unir varios PDF. & \ref{sec:pdfunite}\\*
		\cmd{pdfdetach}, \cmd{pdfattach}, \cmd{pdfsig} & Adjuntos y firmas digitales. & \ref{sec:pdf-otras}\\
	\end{longtable}
	\endgroup
	
	\necesitaespacio{8\baselineskip}
	\subsection{Instalación y comprobación}
	
	\begin{lstlisting}
	sudo apt install poppler-utils
	pdftoppm -v                           # versión (todas las utilidades aceptan -v)
	dpkg -L poppler-utils | grep /usr/bin/  # utilidades realmente instaladas
\end{lstlisting}
	
	\necesitaespacio{18\baselineskip}
	\subsection{Opciones comunes a casi todas las utilidades}
	
	\necesitaespacio{14\baselineskip}
	\begin{opciones}
		\opt{-f N} & Primera página a procesar (desde 1).\\*
		\opt{-l N} & Última página a procesar. Con \opt{-f 3 -l 3} se procesa solo la página 3.\\
		\opt{-opw CLAVE} & Contraseña de propietario de un PDF cifrado.\\
		\opt{-upw CLAVE} & Contraseña de usuario (la de apertura).\\
		\opt{-q} & Silencia avisos y errores (útil con PDF imperfectos que generan muchos avisos).\\*
		\opt{-v}, \opt{-h} & Versión y ayuda.\\
	\end{opciones}
	
	\necesitaespacio{8\baselineskip}
	\subsection{\texttt{pdfinfo}: información del documento}\label{sec:pdfinfo}
	
	\begin{lstlisting}[style=sintaxis]
  pdfinfo [opciones] <archivo.pdf>
\end{lstlisting}
	
	\necesitaespacio{21\baselineskip}
	\begin{opciones}
		\opt{-box} & Muestra las cajas de página (MediaBox, CropBox, BleedBox, TrimBox, ArtBox).\\*
		\opt{-meta} & Muestra el flujo de metadatos XMP.\\
		\opt{-isodates} & Fechas en formato ISO-8601.\\
		\opt{-rawdates} & Fechas tal como están en el PDF.\\
		\opt{-js} & Muestra el código JavaScript del documento.\\
		\opt{-struct}, \opt{-struct-text} & Árbol de estructura lógica de un PDF etiquetado (con o sin texto).\\
		\opt{-dests} & Destinos con nombre (anclas de enlaces internos).\\
		\opt{-url} & Enlaces (URL) del documento.\\
		\opt{-custom} & Metadatos personalizados del diccionario de información.\\*
		\opt{-f}, \opt{-l} & Con varias páginas, muestra el tamaño y la rotación de cada una.\\
	\end{opciones}
	
	\begin{lstlisting}
		pdfinfo doc.pdf
		pdfinfo -f 1 -l 3 -box doc.pdf      # cajas de las páginas 1 a 3
		pdfinfo -meta doc.pdf | less        # metadatos XMP (salir con q)
\end{lstlisting}
	
	Campos clave de la salida: \texttt{Pages} (número de páginas), \texttt{Page size} (en puntos: 1 pt = 1/72 de pulgada; A4 = 595 × 842 pt), \texttt{Encrypted}, \texttt{Optimized} (si está linealizado para la web), \texttt{PDF version}, \texttt{Producer} y \texttt{Creator} (programas que lo generaron) y \texttt{Tagged} (si tiene estructura lógica).
	
	\necesitaespacio{8\baselineskip}
	\subsection{\texttt{pdffonts}: fuentes del documento}\label{sec:pdffonts}
	
	\begin{lstlisting}[style=sintaxis]
  pdffonts [opciones] <archivo.pdf>
\end{lstlisting}
	
	\begin{lstlisting}
  pdffonts doc.pdf
  pdffonts -subst doc.pdf        # fuentes no incrustadas y por cuál se sustituyen
\end{lstlisting}
	
	Columnas: \texttt{name}, \texttt{type} (Type 1, TrueType, CID\ldots), \texttt{encoding}, \texttt{emb} (incrustada), \texttt{sub} (subconjunto), \texttt{uni} (tiene tabla de correspondencia con Unicode). Interpretación práctica:
	
	\begin{itemize}
		\item \textbf{Tabla vacía}: el PDF no tiene texto; casi seguro es un escaneo. Necesita OCR.
		\item \textbf{\texttt{uni = no}}: el texto puede extraerse con caracteres incorrectos (pdftotext y Marker devolverán basura en esas fuentes).
		\item \textbf{\texttt{emb = no}}: la fuente no viaja dentro del PDF; otro equipo la sustituirá y el aspecto puede cambiar.
		\item \textbf{Fuente \texttt{GlyphLessFont}}: capa de texto invisible añadida por OCR (Tesseract/OCRmyPDF).
	\end{itemize}
	
	\necesitaespacio{8\baselineskip}
	\subsection{\texttt{pdftotext}: extraer texto}\label{sec:pdftotext}
	
	\begin{lstlisting}[style=sintaxis]
  pdftotext [opciones] <archivo.pdf> [<salida.txt> | -]
\end{lstlisting}
	
	Si se omite la salida, crea un \arch{.txt} con el mismo nombre; con \val{-} escribe en la pantalla (salida estándar), lo que permite encadenar con otras órdenes.
	
	\necesitaespacio{10\baselineskip}
	\begin{opciones}
		\opt{-layout} & Mantiene la disposición física (columnas y alineaciones) con espacios. Ideal para tablas simples y documentos a dos columnas que deben leerse «como se ven».\\*
		\opt{-raw} & Texto en el orden en que aparece en el flujo del PDF. Útil con algunas maquetas donde \opt{-layout} mezcla columnas.\\
		\opt{-fixed N} & Como \opt{-layout}, pero con un ancho de carácter fijo de \val{N} puntos.\\
		\opt{-nopgbrk} & No inserta saltos de página (carácter \val{\bs{}f}) entre páginas.\\
		\opt{-enc CODIF} & Codificación de salida; por defecto UTF-8. \opt{-listenc} muestra las disponibles.\\
		\opt{-eol unix|\allowbreak{}dos|\allowbreak{}mac} & Tipo de fin de línea. \val{dos} produce archivos con CRLF para programas de Windows antiguos.\\
		\opt{-nodiag} & Descarta el texto no horizontal (por ejemplo, marcas de agua en diagonal).\\
		\opt{-x -y -W -H} & Extrae solo el área rectangular indicada (en píxeles a la resolución \opt{-r}, 72 ppp por defecto).\\
		\opt{-cropbox} & Usa la CropBox (área visible) en lugar de la MediaBox.\\
		\opt{-htmlmeta} & Genera HTML sencillo con los metadatos en la cabecera.\\
		\opt{-bbox}, \opt{-bbox-layout} & HTML con las coordenadas de cada palabra (y de bloques, líneas y flujos con \opt{-bbox-layout}).\\*
		\opt{-tsv} & Tabla separada por tabuladores con coordenadas y nivel de cada palabra (útil para análisis posterior).\\
	\end{opciones}
	
	\begin{lstlisting}
  pdftotext doc.pdf                          # crea doc.txt
  pdftotext -layout doc.pdf doc.txt          # conservando la disposición
  pdftotext -f 2 -l 5 -layout doc.pdf - | less   # páginas 2-5 en pantalla
  pdftotext -layout doc.pdf - | grep -i "presupuesto"   # buscar sin crear archivos
  pdftotext -nopgbrk -eol dos doc.pdf doc_windows.txt
\end{lstlisting}
	
	\begin{nota}[Texto vacío]
		Si \cmd{pdftotext} produce un archivo vacío o solo saltos de página, el PDF no contiene texto (escaneo). Aplique OCR con OCRmyPDF (sección~\ref{sec:ocrmypdf}) y repita.
	\end{nota}
	
	\necesitaespacio{8\baselineskip}
	\subsection{\texttt{pdftoppm}: páginas a imágenes}\label{sec:pdftoppm}
	
	\begin{lstlisting}[style=sintaxis]
  pdftoppm [opciones] <archivo.pdf> <prefijo>
\end{lstlisting}
	
	Rasteriza cada página (la convierte en una imagen de píxeles). Sin opciones de formato produce PPM; lo habitual es pedir PNG (sin pérdida) o JPEG (con pérdida, archivos menores).
	
	\flujoh{doc.pdf (3 páginas),pdftoppm -png,pag-1.png  pag-2.png  pag-3.png}
	
	\necesitaespacio{10\baselineskip}
	\begin{opciones}
		\opt{-png}, \opt{-jpeg}, \opt{-tiff} & Formato de salida (por defecto PPM, o PGM/PBM con \opt{-gray}/\opt{-mono}).\\*
		\opt{-r PPP} & Resolución en puntos por pulgada (predeterminada 150). \opt{-rx} y \opt{-ry} la fijan por eje.\\
		\opt{-scale-to N} & Ajusta el lado mayor a \val{N} píxeles (ignora \opt{-r}). \opt{-scale-to-x} y \opt{-scale-to-y} fijan ancho o alto (\val{-1} conserva la proporción).\\
		\opt{-f N}, \opt{-l N} & Primera y última página.\\
		\opt{-o}, \opt{-e} & Solo páginas impares o solo pares.\\
		\opt{-singlefile} & Escribe solo la primera página convertida, sin número en el nombre (\arch{prefijo.png}).\\
		\opt{-gray}, \opt{-mono} & Escala de grises o blanco y negro (1 bit).\\
		\opt{-jpegopt OPC} & Opciones JPEG separadas por comas: \val{quality=N} (0--100), \val{progressive=\allowbreak{}y|\allowbreak{}n}, \val{optimize=y|n}.\\
		\opt{-tiffcompression C} & Compresión TIFF: \val{none}, \val{packbits}, \val{jpeg}, \val{lzw} o \val{deflate}.\\
		\opt{-cropbox} & Usa la CropBox (área visible) en lugar de la MediaBox.\\
		\opt{-x -y -W -H} & Recorta un área (en píxeles de la imagen resultante).\\
		\opt{-hide-\allowbreak{}annotations} & No dibuja anotaciones (comentarios, sellos, campos de formulario).\\
		\opt{-aa yes|no}, \opt{-aaVector yes|\allowbreak{}no} & Suavizado de fuentes y de vectores; desactivarlo da bordes duros (a veces útil para OCR o para imágenes en 1 bit).\\
		\opt{-sep C} & Carácter entre prefijo y número (predeterminado \val{-}).\\
		\opt{-forcenum} & Añade el número de página aunque se use \opt{-singlefile}.\\*
		\opt{-progress} & Muestra el avance en la terminal.\\
	\end{opciones}
	
	\paragraph{Nombres de salida} \cmd{pdftoppm} añade al prefijo el número de página, con tantos dígitos como tenga el número de la última página: \arch{pag-1.png} en un PDF de 3 páginas, \arch{pag-01.png} en uno de 12 y \arch{pag-001.png} en uno de 150. Esto mantiene el orden correcto con los comodines.
	
	\paragraph{Resolución y tamaño en píxeles} El ancho en píxeles es aproximadamente \emph{ancho en puntos} $\times$ \emph{ppp} / 72. Una página A4 (595 × 842 pt) a 150 ppp produce unos 1240 × 1754 píxeles; a 300 ppp, el doble en cada lado y cuatro veces más píxeles. Referencias: 72--100 ppp para miniaturas, 150 ppp para pantalla, 300 ppp para impresión y OCR, 600 ppp solo para detalles finos (archivos muy grandes).
	\newpage
	
	\begin{lstlisting}
	pdftoppm -png -r 150 doc.pdf pag                       # todas las páginas, PNG
	pdftoppm -jpeg -jpegopt quality=85 -r 150 doc.pdf pag  # JPEG con calidad 85
	pdftoppm -png -r 300 -f 3 -l 3 doc.pdf p3              # solo la página 3
	pdftoppm -png -r 200 -f 1 -singlefile doc.pdf portada  # crea portada.png
	pdftoppm -png -scale-to 1024 doc.pdf mini              # lado mayor = 1024 px
	# para OCR externo
	pdftoppm -gray -r 300 -tiff -tiffcompression lzw escaneado.pdf ocr
	mkdir -p png && pdftoppm -png -r 150 "/mnt/c/Users/NOMBRE/Downloads/doc.pdf" \
	png/pag
\end{lstlisting}
	
	\begin{advertencia}[Prefijo frente a carpeta]
		El segundo argumento es un \emph{prefijo de nombre}, no una carpeta. \cmd{pdftoppm -png doc.pdf salida} crea \arch{salida-1.png}\ldots{} en la carpeta actual. Para escribir en una carpeta, créela y úsela en el prefijo: \arch{png/pag}.
	\end{advertencia}
	
	\necesitaespacio{8\baselineskip}
	\subsection{\texttt{pdftocairo}: imágenes y formatos vectoriales}\label{sec:pdftocairo}
	
	\begin{lstlisting}[style=sintaxis]
 pdftocairo <-png|-jpeg|-tiff|-svg|-pdf|-ps|-eps> [opciones] <archivo.pdf> [<salida>]
\end{lstlisting}
	
	Usa el motor gráfico Cairo. Para imágenes acepta casi las mismas opciones que \cmd{pdftoppm} (con un prefijo como salida); además produce formatos vectoriales, donde la salida es un nombre de archivo.
	
	\necesitaespacio{24\baselineskip}
	\begin{opciones}
		\opt{-png}, \opt{-jpeg}, \opt{-tiff} & Imágenes, como en \cmd{pdftoppm} (\opt{-r}, \opt{-scale-to}, \opt{-singlefile}, \opt{-gray}, \opt{-mono}, \opt{-jpegopt}\ldots).\\*
		\opt{-transp} & PNG o TIFF con fondo transparente en lugar de blanco.\\
		\opt{-antialias M} & Suavizado: \val{default}, \val{none}, \val{gray}, \val{subpixel}, \val{fast}, \val{good}, \val{best}.\\
		\opt{-svg} & SVG vectorial. Use una página por archivo (\opt{-f N -l N}).\\
		\opt{-pdf} & Vuelve a escribir el PDF a través de Cairo (normaliza algunos PDF problemáticos).\\
		\opt{-ps}, \opt{-eps} & PostScript (todas las páginas) o EPS (una sola página). \opt{-level2}/\opt{-level3} fijan el nivel de PostScript.\\
		\opt{-paper T} & Tamaño de papel para PS/PDF: \val{letter}, \val{legal}, \val{A4}, \val{A3} o \val{match} (el de cada página).\\*
		\opt{-origpagesizes}, \opt{-nocrop}, \opt{-expand}, \opt{-noshrink}, \opt{-nocenter} & Control del ajuste de cada página al papel en salidas PS/PDF.\\
	\end{opciones}
	
	\begin{lstlisting}
 pdftocairo -svg -f 1 -l 1 doc.pdf pagina1.svg          # página 1 como SVG
 pdftocairo -png -transp -r 300 -f 2 -l 2 grafico.pdf g  # PNG con transparencia
 pdftocairo -eps -f 1 -l 1 figura.pdf figura.eps        # para LaTeX antiguo
 pdftocairo -pdf problematico.pdf normalizado.pdf
 for n in $(seq 1 "$(pdfinfo doc.pdf | awk '/^Pages:/{print $2}')"); do
 pdftocairo -svg -f "$n" -l "$n" doc.pdf "pag-$n.svg"   # todas las páginas a SVG
 done
\end{lstlisting}
	
	\begin{nota}[¿pdftoppm o pdftocairo?]
		Para imágenes ambos dan resultados de calidad; \cmd{pdftoppm} (motor Splash) es la referencia habitual, y \cmd{pdftocairo} aporta transparencia (\opt{-transp}) y salidas vectoriales. Si una página se ve mal con uno, pruebe el otro.
	\end{nota}
	
	\necesitaespacio{8\baselineskip}
	\subsection{\texttt{pdfimages}: extraer imágenes incrustadas}\label{sec:pdfimages}
	
	\begin{lstlisting}[style=sintaxis]
  pdfimages [opciones] <archivo.pdf> <prefijo>
\end{lstlisting}
	
	Extrae las imágenes tal como están guardadas dentro del PDF, con su resolución original (a diferencia de \cmd{pdftoppm}, que «fotografía» la página entera).
	
	\necesitaespacio{18\baselineskip}
	\begin{opciones}
		\opt{-list} & No extrae: lista cada imagen con página, tamaño en píxeles, color, compresión y resolución efectiva (\texttt{x-ppi}, \texttt{y-ppi}).\\*
		\opt{-png}, \opt{-tiff} & Guarda en PNG o TIFF (por defecto PPM/PBM).\\
		\opt{-j} & Guarda las imágenes JPEG como \arch{.jpg} sin recomprimir.\\
		\opt{-jp2}, \opt{-jbig2}, \opt{-ccitt} & Guarda en su formato original JPEG2000, JBIG2 o CCITT.\\
		\opt{-all} & Equivale a \opt{-png -tiff -j -jp2 -jbig2 -ccitt}: cada imagen en su formato nativo o sin pérdida.\\
		\opt{-p} & Incluye el número de página en el nombre: \arch{img-\allowbreak{}001-\allowbreak{}000.png}.\\*
		\opt{-print-\allowbreak{}filenames} & Muestra el nombre de cada archivo creado.\\
	\end{opciones}
	
	\begin{lstlisting}
		pdfimages -list doc.pdf               # inventario antes de extraer
		mkdir -p imagenes && pdfimages -all -p doc.pdf imagenes/img
		pdfimages -png -f 5 -l 5 doc.pdf p5   # solo las de la página 5
\end{lstlisting}
	
	\begin{limitacion}
		Solo extrae imágenes de mapa de bits. Los gráficos vectoriales (diagramas de Excel, matplotlib, dibujos) no son imágenes y no aparecen: rasterice la página con \cmd{pdftoppm} o conviértala a SVG con \cmd{pdftocairo}. Las máscaras de transparencia pueden aparecer como imágenes separadas (tipo \texttt{smask} en \opt{-list}).
	\end{limitacion}
	
	\necesitaespacio{8\baselineskip}
	\subsection{\texttt{pdftohtml}: PDF a HTML o XML}\label{sec:pdftohtml}
	
	\begin{lstlisting}[style=sintaxis]
  pdftohtml [opciones] <archivo.pdf> [<salida>]
\end{lstlisting}
	
	\necesitaespacio{24\baselineskip}
	\begin{opciones}
		\opt{-s} & Un único documento HTML con todas las páginas.\\*
		\opt{-noframes} & Sin marcos (frames); combinado con \opt{-s} da un solo archivo.\\
		\opt{-c} & Modo «complejo»: posiciona cada texto en su lugar y genera una imagen de fondo por página. Más fiel visualmente, menos editable.\\
		\opt{-i} & Ignora las imágenes.\\
		\opt{-dataurls} & Incrusta las imágenes dentro del HTML como \texttt{data:} URL.\\
		\opt{-fmt png|jpg} & Formato de las imágenes de fondo (con \opt{-c}).\\
		\opt{-xml} & Salida XML con coordenadas, útil para procesar con programas.\\
		\opt{-zoom Z} & Factor de ampliación (predeterminado 1,5).\\
		\opt{-stdout} & Escribe en la salida estándar.\\
		\opt{-enc C} & Codificación de salida.\\*
		\opt{-hidden} & Incluye texto oculto (por ejemplo, la capa de OCR).\\
	\end{opciones}
	
	\begin{lstlisting}
	pdftohtml -s -noframes -i doc.pdf doc.html     # un único HTML, sin imágenes
	pdftohtml -c -s doc.pdf doc_fiel.html          # aspecto similar al PDF
	pdftohtml -xml -i doc.pdf doc                  # crea doc.xml
\end{lstlisting}
	
	\necesitaespacio{8\baselineskip}
	\subsection{\texttt{pdftops}: PDF a PostScript o EPS}\label{sec:pdftops}
	
	\begin{lstlisting}[style=sintaxis]
  pdftops [opciones] <archivo.pdf> [<salida.ps>]
\end{lstlisting}
	
	\necesitaespacio{18\baselineskip}
	\begin{opciones}
		\opt{-level1}, \opt{-level2}, \opt{-level3} & Nivel de PostScript (predeterminado: nivel 2). Las variantes \opt{-level2sep}/\opt{-level3sep} separan colores.\\*
		\opt{-eps} & Genera EPS (requiere una sola página: \opt{-f N -l N}).\\
		\opt{-paper T} & \val{letter}, \val{legal}, \val{A4}, \val{A3} o \val{match}. \opt{-paperw}/\opt{-paperh} en puntos.\\
		\opt{-origpagesizes} & Conserva el tamaño original de cada página.\\
		\opt{-duplex} & Marca el trabajo para impresión a doble cara.\\
		\opt{-r PPP} & Resolución para las partes que deban rasterizarse (transparencias).\\*
		\opt{-binary} & Escribe datos binarios (archivo menor; algunas impresoras no lo admiten).\\
	\end{opciones}
	
	\begin{lstlisting}
				pdftops doc.pdf doc.ps
				pdftops -level3 -paper A4 doc.pdf doc_a4.ps
				pdftops -eps -f 2 -l 2 doc.pdf figura.eps
\end{lstlisting}
	
	\necesitaespacio{8\baselineskip}
	\subsection{\texttt{pdfseparate}: una página por archivo}\label{sec:pdfseparate}
	
	\begin{lstlisting}[style=sintaxis]
  pdfseparate [-f N] [-l N] <archivo.pdf> <patrón_con_%d.pdf>
\end{lstlisting}
	
	\begin{lstlisting}
  pdfseparate doc.pdf pagina-%d.pdf           # pagina-1.pdf, pagina-2.pdf, ...
  pdfseparate -f 3 -l 5 doc.pdf p-%03d.pdf    # p-003.pdf ... p-005.pdf (3 cifras)
\end{lstlisting}
	
	\verb|%d| se sustituye por el número de página; \verb|%03d| lo rellena con ceros hasta tres cifras.
	
	\necesitaespacio{8\baselineskip}
	\subsection{\texttt{pdfunite}: unir PDF}\label{sec:pdfunite}
	
	\begin{lstlisting}[style=sintaxis]
  pdfunite <entrada1.pdf> <entrada2.pdf> [...] <salida.pdf>
\end{lstlisting}
	
	\begin{lstlisting}
	pdfunite portada.pdf capitulo1.pdf capitulo2.pdf libro.pdf
	pdfunite pagina-*.pdf unido.pdf      # ojo al orden: ver sección de comodines
\end{lstlisting}
	
	\begin{importante}[El último nombre es la salida]
		\cmd{pdfunite} no tiene opción \opt{-o}: el \emph{último} archivo de la lista es siempre el destino. Si olvida escribirlo, el último PDF de entrada será tratado como salida y sobrescrito. Para uniones con selección de páginas, rotaciones o mejor conservación de la estructura, use QPDF (sección~\ref{sec:qpdf-paginas}).
	\end{importante}
	
	\necesitaespacio{8\baselineskip}
	\subsection{Otras utilidades: adjuntos y firmas}\label{sec:pdf-otras}
	
	\begin{lstlisting}
	pdfdetach -list doc.pdf                         # listar archivos adjuntos
	pdfdetach -saveall -o adjuntos/ doc.pdf          # extraerlos todos
	pdfdetach -save 1 -o datos.xlsx doc.pdf          # extraer el adjunto número 1
	pdfattach doc.pdf datos.csv doc_con_adjunto.pdf  # añadir un adjunto
	pdfsig doc.pdf                                   # verificar firmas digitales
\end{lstlisting}
	
	\necesitaespacio{29\baselineskip}
	\subsection{Formatos de entrada y salida}
	
	\necesitaespacio{25\baselineskip}
	\begin{formatos}
		PDF & \compat & \condic & Única entrada de todas las utilidades. Salida PDF solo con \cmd{pdfseparate}, \cmd{pdfunite}, \cmd{pdfattach} y \cmd{pdftocairo -pdf}.\\*
		TXT, TSV & \nocompat & \compat & \cmd{pdftotext} (texto UTF-8; TSV con coordenadas).\\
		PNG, JPEG, TIFF & \nocompat & \compat & \cmd{pdftoppm}, \cmd{pdftocairo}; \cmd{pdfimages} extrae las incrustadas.\\
		PPM, PGM, PBM & \nocompat & \compat & Salida predeterminada de \cmd{pdftoppm} y \cmd{pdfimages}.\\
		JPEG2000, JBIG2, CCITT & \nocompat & \condic & Solo \cmd{pdfimages}, y solo si la imagen ya estaba así dentro del PDF.\\
		SVG & \nocompat & \compat & \cmd{pdftocairo -svg} (una página por archivo).\\
		PS, EPS & \nocompat & \compat & \cmd{pdftops}, \cmd{pdftocairo}. EPS: una página.\\
		HTML, XML & \nocompat & \compat & \cmd{pdftohtml}; HTML con coordenadas mediante \cmd{pdftotext -bbox}.\\*
		Imágenes, DOCX\ldots & \nocompat & \nocompat & Poppler no lee nada que no sea PDF.\\
	\end{formatos}
	
	\necesitaespacio{8\baselineskip}
	\subsection{Limitaciones}
	
	\begin{limitacion}
		\begin{itemize}
			\item \textbf{Sin OCR}: en un PDF escaneado, \cmd{pdftotext} no devuelve nada.
			\item \textbf{Orden de lectura}: en maquetas complejas (varias columnas, cuadros laterales), el texto puede salir desordenado; pruebe \opt{-layout}, \opt{-raw} o Marker.
			\item \textbf{Tablas}: \cmd{pdftotext} no reconoce tablas, solo alinea texto con espacios.
			\item \textbf{Fuentes sin Unicode}: si \cmd{pdffonts} muestra \texttt{uni = no}, el texto extraído puede ser ilegible; la única solución es reconocerlo como imagen (OCR).
			\item \textbf{Memoria y disco}: rasterizar a 600 ppp documentos largos produce archivos enormes; calcule el tamaño antes.
			\item \textbf{Fusiones simples}: \cmd{pdfunite} no combina marcadores ni formularios de forma completa.
			\item \textbf{PDF dañados}: Poppler tolera muchos errores, pero no repara la estructura (use QPDF).
		\end{itemize}
	\end{limitacion}
	
	\necesitaespacio{14\baselineskip}
	\subsection{Errores frecuentes}
	
	\necesitaespacio{10\baselineskip}
	\begin{errores}
		\texttt{I/\allowbreak{}O Error: Couldn't open file} & Ruta incorrecta o espacios sin comillas. & Ruta entre comillas; compruebe con \cmd{ls -l "ruta"}.\\*
		\texttt{Command Line Error: Incorrect password} & PDF cifrado. & \opt{-upw} (o \opt{-opw}) con la contraseña.\\
		Archivo de texto vacío & PDF escaneado. & OCR con OCRmyPDF; compruebe con \cmd{pdffonts}.\\
		Texto con símbolos extraños & Fuentes sin tabla Unicode (\texttt{uni no}). & OCR (\cmd{ocrmypdf -{}-force-\allowbreak{}ocr}) o Marker con \opt{-{}-force\_ocr}.\\
		Imágenes enormes o lentas & Resolución excesiva. & Baje \opt{-r} o use \opt{-scale-to}.\\
		Nombres \arch{-01}, \arch{-001} inesperados & Relleno con ceros según el número de páginas. & Es intencionado; use \opt{-singlefile} o \opt{-sep} si necesita otro formato.\\
		Muchos \texttt{Syntax Warning} & PDF imperfecto pero legible. & Revise el resultado; \opt{-q} los oculta. Si hay errores reales, repare con \cmd{qpdf}.\\
		\texttt{Permission denied} al escribir en \arch{/mnt/c} & El archivo de destino está abierto en un visor de Windows que lo bloquea. & Cierre el visor o escriba con otro nombre.\\*
		Se sobrescribió un PDF con \cmd{pdfunite} & Faltaba el nombre de salida. & Use siempre un nombre nuevo como último argumento; trabaje sobre copias.\\
	\end{errores}
	
	Ayuda: \cmd{pdftoppm -h} (igual en todas las utilidades) y \cmd{man pdftoppm}. Versión: \cmd{pdftoppm -v}.
	% =====================================================================
	\clearpage
	\section{img2pdf}\label{sec:img2pdf}
	% =====================================================================
	
	\necesitaespacio{8\baselineskip}
	\subsection{Descripción general}
	
	\textbf{Qué es.} Un conversor de imágenes a PDF que \emph{no recomprime} las imágenes cuando el formato lo permite. Los JPEG y JPEG~2000 se copian dentro del PDF byte a byte; el resto (PNG, TIFF\ldots) se guarda sin pérdida. El resultado conserva exactamente los píxeles originales y suele ocupar casi lo mismo que las imágenes.
	
	\textbf{Qué problema resuelve.} Unir fotografías, escaneos o capturas en un PDF sin degradar la calidad ni inflar el tamaño, cosa que sí ocurre con conversores que decodifican y vuelven a codificar cada imagen.
	
	\textbf{Cuándo usarlo:} escaneos de un móvil o un escáner que deben convertirse en un PDF; preparar la entrada de OCRmyPDF; reconstruir un PDF tras procesar sus páginas como imágenes.
	
	\textbf{Cuándo no usarlo:} si hay que modificar las imágenes (recortar, girar con precisión, ajustar contraste), hágalo antes con ImageMagick; img2pdf solo coloca. Tampoco añade texto reconocible (para eso, OCRmyPDF).
	
	\necesitaespacio{8\baselineskip}
	\subsection{Instalación y comprobación}
	
	\begin{lstlisting}
	sudo apt install img2pdf
	img2pdf --version
	img2pdf --help | less
\end{lstlisting}
	
	\necesitaespacio{8\baselineskip}
	\subsection{Sintaxis}
	
	\begin{lstlisting}[style=sintaxis]
  img2pdf [opciones] <imagen1> [<imagen2> ...] -o <salida.pdf>
\end{lstlisting}
	
	Cada imagen se convierte en una página, en el orden en que aparecen en la orden. Sin opciones de página, cada página mide lo mismo que su imagen según su resolución (si la imagen no indica resolución, se asumen 96 ppp).
	
	\flujoh{foto1.jpg  foto2.jpg  foto3.jpg,img2pdf,album.pdf (3 páginas)}
	
	\necesitaespacio{14\baselineskip}
	\subsection{Opciones principales}
	
	\necesitaespacio{10\baselineskip}
	\begin{opciones}
		\opt{-o ARCHIVO} & Archivo de salida. Sin \opt{-o}, el PDF se escribe en la salida estándar (para encadenar con \texttt{|}).\\*
		\opt{-S T}, \opt{-{}-pagesize T} & Tamaño de página: \val{A4}, \val{A5}, \val{Letter}, \val{Legal}\ldots{} o medidas \val{21cmx29.7cm}, \val{210mmx297mm}, \val{8.5inx11in}. \val{A4\^{}T} gira el papel a horizontal.\\
		\opt{-s T}, \opt{-{}-imgsize T} & Tamaño de la \emph{imagen} en la página; admite resolución: \val{300dpi}.\\
		\opt{-b M}, \opt{-{}-border M} & Margen alrededor de la imagen: \val{1cm} o \val{vertical:horizontal}, por ejemplo \val{2cm:1.5cm}.\\
		\opt{-f M}, \opt{-{}-fit M} & Cómo encajar la imagen en la página: \val{into} (predeterminado: cabe entera, mantiene proporción), \val{fill} (cubre la página, mantiene proporción; puede sobresalir), \val{exact} (deforma para ajustar exactamente), \val{shrink} (como \val{into} pero solo reduce), \val{enlarge} (solo amplía).\\
		\opt{-a}, \opt{-{}-auto-orient} & Con \opt{-{}-pagesize}, gira el papel para que su orientación coincida con la de cada imagen (vertical u horizontal).\\
		\opt{-r R}, \opt{-{}-rotation R} & Rotación: \val{auto} (predeterminada: usa la etiqueta EXIF de orientación), \val{none}, \val{ifvalid} (como \val{auto}, pero ignora etiquetas EXIF no válidas, frecuentes en móviles y escáneres), \val{0}, \val{90}, \val{180}, \val{270}.\\
		\opt{-{}-title}, \opt{-{}-author}, \opt{-{}-subject}, \opt{-{}-keywords}, \opt{-{}-creator}, \opt{-{}-producer} & Metadatos del PDF.\\
		\opt{-{}-creationdate}, \opt{-{}-moddate} & Fechas de creación y modificación.\\
		\opt{-D}, \opt{-{}-nodate} & No escribe fechas (dos ejecuciones idénticas dan el mismo archivo).\\
		\opt{-C ESPACIO} & Fuerza el espacio de color (\val{L}, \val{RGB}, \val{CMYK}, \val{CMYK;I}\ldots). \val{CMYK;I} corrige JPEG CMYK de Adobe que salen en negativo.\\
		\opt{-{}-first-\allowbreak{}frame-\allowbreak{}only} & En GIF o TIFF de varios fotogramas, solo el primero (por defecto, uno por página).\\
		\opt{-{}-pillow-\allowbreak{}limit-\allowbreak{}break} & Desactiva la protección contra imágenes gigantes («bomba de descompresión»). Solo con imágenes de confianza.\\
		\opt{-{}-pdfa} & Intenta generar un PDF/A-1b.\\
		\opt{-{}-from-file F} & Lee la lista de imágenes de un archivo (o de la entrada estándar con \val{-}) con los nombres \emph{separados por caracteres nulos}. Evita el límite de longitud de la orden.\\
		\opt{-{}-crop-border}, \opt{-{}-bleed-border}, \opt{-{}-trim-border}, \opt{-{}-art-border} & Cajas de página para imprenta.\\
		\opt{-{}-viewer-*} & Preferencias de visualización (página inicial, diseño, ampliación, pantalla completa).\\*
		\opt{-{}-engine M} & Motor de escritura del PDF (interno o pikepdf). Uso avanzado.\\
	\end{opciones}
	
	\necesitaespacio{8\baselineskip}
	\subsection{Ejemplos}
	
	\begin{lstlisting}
  img2pdf imagen1.png imagen2.png -o documento.pdf
  img2pdf *.jpg -o album.pdf  # todas las JPG (orden alfabético)
  # páginas A4, imagen encajada
  img2pdf --pagesize A4 escaneo_*.jpg -o escaneos_a4.pdf
  img2pdf --pagesize A4 --border 1.5cm *.jpg -o con_margen.pdf
  # A4 vertical u horizontal
  img2pdf --pagesize A4 --auto-orient fotos/*.jpg -o fotos.pdf
  img2pdf --pagesize A4^T panoramica.jpg -o horizontal.pdf
  img2pdf --imgsize 300dpi escaneo.png -o tamano_real.pdf  # tamaño físico a 300 ppp
  img2pdf --rotation=ifvalid IMG_*.jpg -o movil.pdf  # EXIF defectuoso del móvil
  img2pdf --title "Actas 2026" --author "Secretaría" --nodate actas/*.png -o \
  actas.pdf
\end{lstlisting}
	
	\begin{lstlisting}
  # Orden natural (pag2 antes que pag10) y nombres con cualquier carácter
  find . -maxdepth 1 -name '*.png' -print0 | sort -zV | img2pdf --from-file - -o \
  natural.pdf
  
  # Tubería directa hacia OCRmyPDF (sin archivo intermedio)
  img2pdf --pagesize A4 escaneo_*.png | ocrmypdf -l spa - escaneo_ocr.pdf
\end{lstlisting}
	
	\necesitaespacio{8\baselineskip}
	\subsection{Comodines, orden y precauciones}
	
	El asterisco lo expande Bash, no img2pdf (sección~\ref{sec:comodines}). Consecuencias prácticas:
	
	\begin{itemize}
		\item \textbf{El orden de las páginas es el orden alfabético de los nombres}. \arch{pag10.jpg} va antes que \arch{pag2.jpg}. Use nombres con ceros (\arch{pag02.jpg}) u orden natural (\cmd{sort -V}).
		\item \textbf{Mayúsculas}: \arch{*.jpg} no incluye \arch{FOTO.JPG}; use \arch{*.[jJ][pP][gG]}.
		\item \textbf{Revise antes}: \cmd{ls -v *.jpg} muestra qué entrará y en qué orden.
		\item \textbf{No mezcle salidas con entradas}: si \arch{documento.pdf} está en la carpeta, un patrón \arch{*} lo incluiría.
	\end{itemize}
	
	\paragraph{Imágenes de distintos tamaños} Sin \opt{-{}-pagesize}, cada página toma el tamaño de su imagen: el PDF tendrá páginas de tamaños diferentes. Con \opt{-{}-pagesize A4}, todas las páginas son A4 y cada imagen se encaja según \opt{-{}-fit}.
	
	\paragraph{Orientaciones distintas} Con \opt{-{}-pagesize A4} y sin \opt{-a}, una foto horizontal queda pequeña en una página vertical. Con \opt{-{}-auto-orient}, la página se gira para seguir a la imagen. La rotación del \emph{contenido} se controla con \opt{-{}-rotation}.
	
	\paragraph{Transparencia} Las imágenes con canal alfa (PNG transparentes) se admiten: img2pdf guarda la transparencia como una máscara aparte e informa de ello con un mensaje.
	
	\necesitaespacio{22\baselineskip}
	\subsection{img2pdf frente a ImageMagick}
	
	\begingroup\small
	\necesitaespacio{18\baselineskip}
	\begin{longtable}{L{3.6cm}L{5.9cm}L{5.56cm}}
		\toprule\rowcolor{tablacab}\cab{Aspecto} & \cab{img2pdf} & \cab{ImageMagick (\cmd{magick *.jpg out.pdf})}\\\midrule
		\endfirsthead
		\bottomrule
		\endfoot
		Calidad & Sin pérdida: JPEG copiado tal cual. & Decodifica y vuelve a codificar: puede perder calidad.\\*
		Tamaño del PDF & Similar a la suma de las imágenes. & Variable; a menudo mayor o con artefactos.\\
		Velocidad y memoria & Muy rápido, poca memoria. & Carga todas las imágenes en memoria.\\
		Edición de imágenes & Ninguna. & Completa (recortar, ajustar, girar\ldots).\\*
		Cuándo elegirlo & Siempre que solo haya que \emph{empaquetar}. & Cuando hay que \emph{transformar}; luego empaquete con img2pdf.\\
	\end{longtable}
	\endgroup
	
	\necesitaespacio{22\baselineskip}
	\subsection{Formatos de entrada y salida}
	
	\necesitaespacio{18\baselineskip}
	\begin{formatos}
		JPEG & \compat & \nocompat & Incrustado sin recomprimir.\\*
		JPEG 2000 & \compat & \nocompat & Incrustado sin recomprimir.\\
		PNG & \compat & \nocompat & Sin pérdida; la transparencia se guarda como máscara.\\
		TIFF & \compat & \nocompat & Sin pérdida; un TIFF de varias páginas da varias páginas.\\
		GIF & \compat & \nocompat & Sin pérdida; un GIF animado da una página por fotograma (salvo \opt{-{}-first-\allowbreak{}frame-\allowbreak{}only}).\\
		Otros legibles por Pillow & \condic & \nocompat & Se convierten sin pérdida; compruebe el resultado.\\*
		PDF & \nocompat & \compat & Única salida.\\
	\end{formatos}
	
	\necesitaespacio{8\baselineskip}
	\subsection{Limitaciones}
	
	\begin{limitacion}
		\begin{itemize}
			\item No edita, no corrige la perspectiva ni recorta: solo coloca imágenes en páginas.
			\item No reduce el tamaño: si las imágenes pesan mucho, el PDF también. Reduzca antes con ImageMagick o después con Ghostscript.
			\item No añade texto buscable (use OCRmyPDF después).
			\item Imágenes enormes pueden rechazarse por la protección de Pillow (\opt{-{}-pillow-\allowbreak{}limit-\allowbreak{}break} con cuidado).
		\end{itemize}
	\end{limitacion}
	
	\necesitaespacio{32\baselineskip}
	\subsection{Errores frecuentes}
	
	\necesitaespacio{28\baselineskip}
	\begin{errores}
		Páginas en orden incorrecto & Orden alfabético del comodín. & Ceros a la izquierda o \cmd{sort -V} con \opt{-{}-from-file}.\\*
		Foto del móvil girada & Etiqueta EXIF ausente o incorrecta. & \opt{-{}-rotation=\allowbreak{}ifvalid}, o un ángulo fijo (\opt{-{}-rotation=90}).\\
		Error sobre la rotación EXIF & Etiqueta de orientación con un valor no válido (frecuente en móviles y escáneres). & \opt{-{}-rotation=\allowbreak{}ifvalid}.\\
		Colores invertidos & JPEG CMYK de Adobe. & \opt{-C CMYK;I}.\\
		\texttt{DecompressionBombError} & Imagen con demasiados píxeles. & Si es de confianza, \opt{-{}-pillow-\allowbreak{}limit-\allowbreak{}break}.\\
		\texttt{Argument list too long} & Demasiados archivos en la orden. & \opt{-{}-from-file} con \cmd{find -print0}.\\*
		La orden se queda esperando & Sin imágenes en la orden (por ejemplo, un comodín sin coincidencias con \cmd{nullglob}), img2pdf lee de la entrada estándar. & \tecla{Ctrl}+\tecla{C}; revise el patrón con \cmd{ls}.\\
	\end{errores}
	
	Ayuda: \cmd{img2pdf -{}-help} (incluye la tabla de tamaños de papel). Versión: \cmd{img2pdf -{}-version}.
	
	% =====================================================================
	\clearpage
	\section{ImageMagick}\label{sec:imagemagick}
	% =====================================================================
	
	\necesitaespacio{8\baselineskip}
	\subsection{Descripción general}
	
	\textbf{Qué es.} Un conjunto de herramientas para crear, convertir y transformar imágenes en cientos de formatos. Ubuntu 26.04 incluye \textbf{ImageMagick 7}, cuya orden principal es \cmd{magick}.
	
	\textbf{Qué problema resuelve.} Cualquier transformación de imágenes desde la terminal: cambiar formato, tamaño, recorte, orientación, color, contraste, combinar varias, crear rótulos o animaciones, y hacerlo en lotes.
	
	\textbf{Cuándo usarlo:} para preparar imágenes (antes de img2pdf, OCR o Potrace), convertir formatos, generar miniaturas y automatizar retoques.
	
	\textbf{Cuándo no usarlo:} para rasterizar PDF (Poppler da mejor resultado y consume menos memoria), para empaquetar imágenes en PDF sin pérdida (img2pdf) o para comprimir PDF (Ghostscript). ImageMagick no es un sustituto universal de esas herramientas.
	
	\necesitaespacio{20\baselineskip}
	\subsection{ImageMagick 7: \texttt{magick} y los nombres heredados}
	
	En la versión 7 todas las funciones se invocan con \cmd{magick}. Las órdenes de la versión 6 siguen existiendo por compatibilidad:
	
	\begingroup\small
	\necesitaespacio{16\baselineskip}
	\begin{longtable}{L{4.2cm}L{5.0cm}L{5.86cm}}
		\toprule\rowcolor{tablacab}\cab{ImageMagick 6 (heredado)} & \cab{ImageMagick 7} & \cab{Uso}\\\midrule
		\endfirsthead
		\bottomrule
		\endfoot
		\cmd{convert} & \cmd{magick} & Convertir y transformar (muestra un aviso de obsolescencia).\\*
		\cmd{mogrify} & \cmd{magick mogrify} & Modificar muchos archivos (por defecto, \emph{sobrescribe}).\\
		\cmd{identify} & \cmd{magick identify} & Información de una imagen.\\
		\cmd{montage} & \cmd{magick montage} & Mosaicos y hojas de contacto.\\
		\cmd{composite} & \cmd{magick composite} & Superponer imágenes.\\*
		\cmd{compare} & \cmd{magick compare} & Comparar dos imágenes.\\
	\end{longtable}
	\endgroup
	
	\begin{nota}[Aviso de \texttt{convert}]
		\cmd{convert} sigue funcionando pero imprime \texttt{WARNING: The convert command is deprecated in IMv7, use "magick" instead}. Use \cmd{magick} en todos los comandos nuevos: además de evitar el aviso, aplica las operaciones estrictamente en el orden escrito.
	\end{nota}
	
	\necesitaespacio{8\baselineskip}
	\subsection{Instalación y comprobación}
	
	\begin{lstlisting}
	# Ghostscript es necesario para leer PDF/PS/EPS
	sudo apt install imagemagick ghostscript
	magick -version                            # versión, delegados y características
	magick -list format | less  # formatos admitidos (r = leer, w = escribir)
	magick -list policy                        # política de seguridad en vigor
	magick -list resource                      # límites de memoria, disco y tamaño
	magick -list font | grep -i dejavu         # nombres de fuentes para -font
\end{lstlisting}
	
	\necesitaespacio{8\baselineskip}
	\subsection{Sintaxis y orden de las operaciones}
	
	\begin{lstlisting}[style=sintaxis]
  magick [ajustes_de_lectura] <entrada...> [operaciones] <salida>
\end{lstlisting}
	
	ImageMagick~7 procesa los argumentos de izquierda a derecha: lee las entradas, aplica cada operación a las imágenes ya leídas y escribe la salida (el formato lo determina la extensión). Hay dos tipos de opciones:
	
	\begin{itemize}
		\item \textbf{Ajustes} que afectan a la lectura o escritura y deben ir \emph{antes} de lo que afectan: \opt{-density} antes de un PDF o SVG, \opt{-size} antes de \val{xc:}, \opt{-quality} antes del nombre de salida.
		\item \textbf{Operadores} que transforman lo ya leído: \opt{-resize}, \opt{-crop}, \opt{-rotate}\ldots
	\end{itemize}
	
	Los paréntesis escapados \verb|\( ... \)| aplican operaciones solo a un subconjunto de imágenes.
	
	\necesitaespacio{26\baselineskip}
	\subsubsection{Geometría: cómo se expresan tamaños}\label{sec:geometria}
	
	\necesitaespacio{22\baselineskip}
	\begin{opciones}[Geometría]
		\val{50\%} & Escala al 50\,\% en ambas dimensiones.\\*
		\val{800x600} & Cabe dentro de 800×600 \emph{manteniendo la proporción} (puede quedar 800×450).\\
		\val{800x} / \val{x600} & Ancho 800 (alto proporcional) / alto 600 (ancho proporcional).\\
		\val{800x600!} & Exactamente 800×600, deformando si hace falta.\\
		\val{800x600>} & Solo reduce si la imagen es mayor; nunca amplía.\\
		\val{800x600<} & Solo amplía si es menor.\\
		\val{800x600\^{}} & Cubre 800×600 como mínimo (para recortar después con \opt{-extent}).\\
		\val{@250000} & Ajusta el número total de píxeles (área).\\*
		\val{WxH+X+Y} & Rectángulo de \val{W}×\val{H} desplazado \val{X} a la derecha e \val{Y} hacia abajo desde la esquina superior izquierda (o desde la referencia de \opt{-gravity}).\\
	\end{opciones}
	
	\begin{advertencia}[Caracteres especiales]
		\val{!}, \val{>}, \val{<} y \val{\^{}} tienen significado para Bash: escriba la geometría entre comillas, por ejemplo \val{"800x600>"}.
	\end{advertencia}
	
	\necesitaespacio{25\baselineskip}
	\subsection{Operaciones principales}
	
	\necesitaespacio{22\baselineskip}
	\subsubsection{Convertir formatos y comprimir}
	
	\necesitaespacio{18\baselineskip}
	\begin{opciones}
		\opt{-quality N} & JPEG/WebP: calidad 1--100 (predeterminada: la estimada de la imagen de entrada o 92). PNG: no es «calidad» sino nivel y filtro de compresión sin pérdida.\\*
		\opt{-strip} & Elimina metadatos y perfiles (EXIF, miniaturas, comentarios). Reduce tamaño y protege la privacidad (borra la ubicación GPS de las fotos).\\
		\opt{-sampling-\allowbreak{}factor 4:2:0} & Submuestreo de color de JPEG (archivos menores).\\
		\opt{-interlace Plane} & JPEG progresivo (o PNG entrelazado con \val{PNG}).\\
		\opt{-colors N} & Reduce la paleta a \val{N} colores.\\*
		\opt{-depth N} & Bits por canal (8 o 16).\\
	\end{opciones}
	
	\begin{lstlisting}
  magick imagen.png imagen.jpg                         # PNG -> JPEG
  magick foto.jpg -quality 82 -strip foto_web.jpg  # comprimir y quitar metadatos
  magick foto.jpg -quality 80 foto.webp                # JPEG -> WebP
  magick foto.jpg -strip -quality 82 -sampling-factor 4:2:0 -interlace Plane web.jpg
\end{lstlisting}
	
	\necesitaespacio{18\baselineskip}
	\subsubsection{Redimensionar}
	
	\necesitaespacio{14\baselineskip}
	\begin{opciones}
		\opt{-resize GEOM} & Redimensiona con filtros de calidad (geometría en la sección~\ref{sec:geometria}).\\*
		\opt{-thumbnail GEOM} & Como \opt{-resize} y además elimina perfiles: ideal para miniaturas.\\
		\opt{-scale GEOM}, \opt{-sample GEOM} & Rápidos y sin suavizado (píxeles duros); útiles para arte de píxeles.\\*
		\opt{-extent GEOM} & Ajusta el lienzo al tamaño indicado (recorta o rellena con \opt{-background}).\\
	\end{opciones}
	
	\begin{lstlisting}
  magick imagen.png -resize 50% salida.png
  magick foto.jpg -resize 1920x1080 grande.jpg                  # cabe en 1920x1080
  magick foto.jpg -resize "1920x1080>" sin_ampliar.jpg          # solo reduce
  magick foto.jpg -resize "300x300^" -gravity center -extent 300x300 cuadrada.jpg
  magick foto.jpg -thumbnail 200x200 miniatura.jpg
\end{lstlisting}
	
	\necesitaespacio{22\baselineskip}
	\subsubsection{Recortar}
	
	\necesitaespacio{18\baselineskip}
	\begin{opciones}
		\opt{-crop WxH+X+Y} & Recorta un rectángulo. Añada siempre \opt{+repage} después para eliminar el desplazamiento residual del lienzo.\\*
		\opt{-crop 2x2@} & Divide en una cuadrícula de 2×2 piezas iguales (genera varias imágenes).\\
		\opt{-gravity P} & Punto de referencia: \val{center}, \val{north}, \val{southeast}\ldots\\
		\opt{-trim} & Recorta los bordes del color del fondo; con \opt{-fuzz 5\%} tolera pequeñas variaciones.\\
		\opt{-shave WxH} & Quita \val{W} píxeles a izquierda y derecha y \val{H} arriba y abajo.\\*
		\opt{-border N}, \opt{-bordercolor C} & Añade un borde.\\
	\end{opciones}
	
	\begin{lstlisting}
	magick foto.jpg -crop 800x600+100+50 +repage recorte.jpg
	magick foto.jpg -gravity center -crop 1000x1000+0+0 +repage centro.jpg
	magick escaneo.png -fuzz 5% -trim +repage sin_bordes.png
	magick plano.png -crop 2x2@ +repage pieza_%02d.png      # pieza_00 ... pieza_03
\end{lstlisting}
	
	\necesitaespacio{6\baselineskip}
	\subsubsection{Rotar y voltear}
	
	\begin{lstlisting}
	magick foto.jpg -rotate 90 girada.jpg
	# rellena esquinas en blanco
	magick escaneo.png -background white -rotate -2.5 enderezada.png
	# aplica la etiqueta EXIF de orientación
	magick foto.jpg -auto-orient orientada.jpg
	magick foto.jpg -flip invertida_vertical.jpg     # -flop: espejo horizontal
	# corrige inclinación pequeña automáticamente
	magick escaneo.png -deskew 40% enderezado.png
\end{lstlisting}
	
	\necesitaespacio{6\baselineskip}
	\subsubsection{Resolución: \texttt{-density} frente a \texttt{-resample}}
	
	\begin{itemize}
		\item \opt{-density N} con \opt{-units PixelsPerInch} \textbf{solo cambia el metadato} de resolución (los píxeles no cambian), salvo al leer formatos vectoriales (PDF, SVG, EPS), donde fija con cuántos píxeles por pulgada se rasterizan.
		\item \opt{-resample N} cambia la resolución \textbf{y} recalcula los píxeles para mantener el tamaño físico.
	\end{itemize}
	
	\begin{lstlisting}
	# solo metadato
	magick escaneo.png -units PixelsPerInch -density 300 escaneo_300.png
	magick escaneo_300.png -resample 150 escaneo_150.png  # mitad de píxeles
\end{lstlisting}
	
	\necesitaespacio{34\baselineskip}
	\subsubsection{Color, escala de grises, brillo y contraste}
	
	\necesitaespacio{30\baselineskip}
	\begin{opciones}
		\opt{-colorspace Gray} & Convierte a escala de grises (la forma más habitual).\\*
		\opt{-grayscale MÉTODO} & Escala de grises con un método explícito (\val{Rec709Luma}, \val{Average}\ldots).\\
		\opt{-type Grayscale}, \opt{-type Bilevel} & Fuerza el tipo de imagen al guardar (grises o blanco y negro).\\
		\opt{-threshold N\%} & Umbral: píxeles más claros que \val{N} pasan a blanco y el resto a negro.\\
		\opt{-monochrome} & Blanco y negro con tramado (difuminado de puntos).\\
		\opt{-negate} & Negativo.\\
		\opt{-brightness-\allowbreak{}contrast BxC} & Brillo y contraste en porcentaje, de \val{-100} a \val{100}: \val{10x20}.\\
		\opt{-level N\%,M\%} & Estira los tonos: lo más oscuro que \val{N} pasa a negro, lo más claro que \val{M} a blanco. Muy útil para blanquear el fondo de un escaneo.\\
		\opt{-auto-level}, \opt{-normalize}, \opt{-contrast-\allowbreak{}stretch BxW\%} & Ajustes automáticos del rango tonal.\\
		\opt{-gamma G} & Corrección gamma (\val{>1} aclara, \val{<1} oscurece).\\*
		\opt{-sharpen 0xS}, \opt{-blur 0xS} & Enfocar o desenfocar con radio automático y desviación \val{S}.\\
	\end{opciones}
	
	\begin{lstlisting}
  magick foto.jpg -colorspace Gray gris.jpg
  magick escaneo.jpg -colorspace Gray -level 15%,85% limpio.png   # fondo más blanco
  magick foto.jpg -brightness-contrast 10x15 mejorada.jpg
  magick firma.png -colorspace Gray -threshold 55% firma_bn.png
\end{lstlisting}
	
	\necesitaespacio{6\baselineskip}
	\subsubsection{Combinar imágenes}
	
	\begin{lstlisting}
  magick a.jpg b.jpg +append horizontal.jpg        # una al lado de otra
  magick a.jpg b.jpg -append vertical.jpg          # una debajo de otra
  magick montage *.jpg -tile 4x -geometry 300x300+6+6 hoja_contactos.jpg
  magick foto.jpg logo.png -gravity southeast -geometry +20+20 -composite marca.jpg
\end{lstlisting}
	
	\necesitaespacio{6\baselineskip}
	\subsubsection{Crear imágenes y añadir texto}
	
	\begin{lstlisting}
  magick -size 1200x630 xc:white lienzo.png                       # lienzo blanco
  magick -size 1200x300 gradient:white-"#1f4e79" degradado.png    # degradado
  magick -background white -fill black -font DejaVu-Sans -pointsize 48 \
  label:"Título del póster" rotulo.png  # rótulo ajustado al texto
  magick -size 800x -background white -font DejaVu-Sans -pointsize 28 \
  caption:"Un texto largo que se reparte en varias líneas" parrafo.png
  magick foto.jpg -gravity south -fill white -font DejaVu-Sans -pointsize 36 \
  -annotate +0+20 "Borrador" anotada.jpg  # texto sobre una imagen
\end{lstlisting}
	
	\opt{-font} acepta los nombres de \cmd{magick -list font} o la ruta a un archivo \arch{.ttf}/\arch{.otf}. Si una fuente no existe, ImageMagick usa otra sin avisar claramente.
	
	\necesitaespacio{20\baselineskip}
	\subsubsection{Transparencias}
	
	\necesitaespacio{16\baselineskip}
	\begin{opciones}
		\opt{-alpha remove} & Mezcla la transparencia con el color de \opt{-background}: la imagen queda opaca.\\*
		\opt{-alpha off} & Desactiva el canal alfa (sin mezclar: las zonas transparentes pueden verse negras).\\
		\opt{-alpha set}, \opt{-alpha on} & Añade o activa el canal alfa.\\
		\opt{-background C -flatten} & Aplana todas las capas sobre un fondo del color \val{C}.\\*
		\opt{-transparent C} & Vuelve transparente un color; con \opt{-fuzz 10\%} incluye colores parecidos.\\
	\end{opciones}
	
	\begin{lstlisting}
  # PNG transparente -> JPG
  magick logo.png -background white -alpha remove -alpha off logo.jpg
  magick logo_fondo_blanco.png -fuzz 8% -transparent white logo_transparente.png
\end{lstlisting}
	
	\necesitaespacio{6\baselineskip}
	\subsubsection{GIF animados}
	
	\begin{lstlisting}
  # 50 = 0,5 s por fotograma; 0 = bucle infinito
  magick -delay 50 -loop 0 frame_*.png animacion.gif
  magick animacion.gif -layers Optimize optimizada.gif # reduce tamaño
  magick animacion.gif -coalesce fotograma_%03d.png  # separar fotogramas completos
\end{lstlisting}
	
	\necesitaespacio{6\baselineskip}
	\subsubsection{Información de una imagen}
	
	\begin{lstlisting}
  magick identify foto.jpg                         # formato, tamaño, profundidad
  magick identify -verbose foto.jpg | less  # todo: EXIF, perfil, estadísticas
  magick identify -format "%f %wx%h %x\n" *.jpg    # nombre, tamaño y resolución
\end{lstlisting}
	
	\necesitaespacio{13\baselineskip}
	\subsection{Procesamiento por lotes: \texttt{magick mogrify}}
	
	\cmd{magick mogrify} aplica las mismas operaciones a muchos archivos. \textbf{Sin \opt{-path} sobrescribe los originales.}
	
	\necesitaespacio{9\baselineskip}
	\begin{opciones}
		\opt{-path CARPETA} & Escribe los resultados en otra carpeta (que debe existir) y deja intactos los originales.\\
		\opt{-format EXT} & Cambia el formato de salida (crea archivos nuevos con otra extensión).\\
	\end{opciones}
	
	\begin{lstlisting}
	mkdir -p web && magick mogrify -path web -resize "1600x1600>" -quality 82 \
	-strip *.jpg
	# TIFF -> PNG en otra carpeta
	mkdir -p png && magick mogrify -path png -format png *.tif
\end{lstlisting}
	
	\begin{advertencia}[mogrify sobrescribe]
		\cmd{magick mogrify -resize 50\% *.jpg} reduce \emph{los propios archivos} y no hay vuelta atrás. Use siempre \opt{-path} o trabaje sobre copias.
	\end{advertencia}
	
	Para lotes con nombres de salida personalizados, use un bucle:
	
	\begin{lstlisting}
  mkdir -p gris
  for f in *.jpg; do magick "$f" -colorspace Gray "gris/${f%.jpg}_gris.jpg"; done
\end{lstlisting}
	
	\necesitaespacio{8\baselineskip}
	\subsection{ImageMagick y PDF}\label{sec:im-pdf}
	
	ImageMagick no interpreta PDF por sí mismo: delega en Ghostscript para rasterizarlo y, al escribir un PDF, coloca las imágenes como mapas de bits. Rasterizar (convertir en píxeles) no es lo mismo que convertir: el resultado pierde texto seleccionable y vectores.
	
	\begin{lstlisting}
	# Página 1 (índice 0) a PNG a 300 ppp con fondo blanco
	magick -density 300 "doc.pdf[0]" -background white -alpha remove -alpha off \
	pag1.png
	
	# Páginas 1 a 3 (índices 0-2) a JPEG
	magick -density 150 "doc.pdf[0-2]" -background white -alpha remove -quality 85 \
	pag-%02d.jpg
	
	# Imágenes a PDF (recomprime; para conservar calidad, mejor img2pdf)
	magick a.png b.jpg resultado.pdf
\end{lstlisting}
	
	\begin{itemize}
		\item \textbf{\opt{-density} va antes del PDF}: fija la resolución de rasterizado. Sin ella se usan 72 ppp, con resultado pixelado.
		\item \textbf{Índices desde 0}: \arch{doc.pdf[0]} es la primera página; las comillas evitan que Bash interprete los corchetes.
		\item \textbf{Fondo}: las páginas pueden tener fondo transparente; sin \opt{-alpha remove}, un JPEG puede salir con fondo negro.
		\item \textbf{Documentos largos}: ImageMagick carga todas las páginas rasterizadas en memoria y puede agotar los límites de recursos. Para PDF de muchas páginas, use \cmd{pdftoppm} y procese después las imágenes.
	\end{itemize}
	
	\necesitaespacio{22\baselineskip}
	\subsubsection{Política de seguridad (\texttt{policy.xml})}\label{sec:im-policy}
	
	ImageMagick limita qué puede hacer mediante \arch{/\allowbreak{}etc/\allowbreak{}ImageMagick-\allowbreak{}7/\allowbreak{}policy.xml}. En Ubuntu 26.04 la política activa \textbf{permite leer y escribir PDF} (versiones antiguas de Ubuntu lo bloqueaban) y establece:
	
	\begingroup\small
	\necesitaespacio{18\baselineskip}
	\begin{longtable}{L{4.6cm}L{10.46cm}}
		\toprule\rowcolor{tablacab}\cab{Política} & \cab{Efecto}\\\midrule
		\endfirsthead
		\bottomrule
		\endfoot
		\texttt{memory} 1024MiB, \texttt{map} 2048MiB & Memoria máxima para la caché de píxeles antes de recurrir al disco.\\*
		\texttt{disk} 2GiB & Disco máximo para la caché de píxeles. Superarlo produce \texttt{cache resources exhausted}.\\
		\texttt{area} 256MP & Área máxima (megapíxeles) que se mantiene en memoria.\\
		\texttt{width}, \texttt{height} 32KP & Ancho y alto máximos: 32\,000 píxeles.\\
		\texttt{path} \texttt{@*} sin permisos & Prohíbe leer listas de archivos con la sintaxis \arch{@lista.txt}.\\*
		Delegados \texttt{URL}, \texttt{HTTP}, \texttt{HTTPS} sin permisos & No descarga imágenes de Internet.\\
	\end{longtable}
	\endgroup
	
	\begin{advertencia}[Modificar la política]
		Si un trabajo legítimo supera los límites (por ejemplo, rasterizar páginas enormes a alta resolución), es preferible cambiar de herramienta (Poppler) o reducir la resolución antes que relajar la política. Si decide editarla, haga una copia (\cmd{sudo cp /\allowbreak{}etc/\allowbreak{}ImageMagick-\allowbreak{}7/\allowbreak{}policy.xml\{,\allowbreak{}.orig\}}), cambie solo el valor necesario y nunca procese así archivos de origen desconocido.
	\end{advertencia}
	
	\necesitaespacio{26\baselineskip}
	\subsection{Formatos de entrada y salida}
	
	\necesitaespacio{22\baselineskip}
	\begin{formatos}
		PNG, JPEG, GIF, TIFF, BMP, WebP & \compat & \compat & Soporte nativo de la versión de Ubuntu.\\*
		PNM (PBM, PGM, PPM) & \compat & \compat & Formatos sencillos; PBM/PGM son la entrada de Potrace.\\
		PDF, PS, EPS & \condic & \condic & Lectura mediante Ghostscript (rasteriza). Escritura: imágenes de mapa de bits dentro del PDF.\\
		SVG & \condic & \condic & Lectura rasterizando con el intérprete disponible en la compilación; escritura muy limitada (no vectoriza: use Potrace).\\
		HEIC, AVIF, JPEG XL & \condic & \condic & Dependen de los delegados de la compilación: compruébelo con \cmd{magick -list format}.\\*
		GIF animado & \compat & \compat & \opt{-delay}, \opt{-loop}, \opt{-layers Optimize}.\\
	\end{formatos}
	
	\necesitaespacio{8\baselineskip}
	\subsection{Limitaciones}
	
	\begin{limitacion}
		\begin{itemize}
			\item \textbf{Pérdida por recompresión}: cada guardado en JPEG pierde calidad. Trabaje en PNG o TIFF y exporte a JPEG al final.
			\item \textbf{Memoria}: las imágenes se manejan descomprimidas (alta profundidad de color): una imagen de 10\,000 × 10\,000 píxeles ocupa mucha más memoria que su archivo.
			\item \textbf{PDF}: rasteriza; pierde texto y vectores. Para PDF largos, Poppler es más eficiente.
			\item \textbf{Fuentes}: si la fuente pedida no existe, se sustituye. Compruebe con \cmd{magick -list font}.
			\item \textbf{Perfiles de color}: \opt{-strip} elimina también el perfil ICC; en fotografías con perfil distinto de sRGB, los colores pueden cambiar.
			\item \textbf{Política de seguridad}: límites de recursos y prohibiciones descritas en la sección~\ref{sec:im-policy}.
		\end{itemize}
	\end{limitacion}
	
	\necesitaespacio{14\baselineskip}
	\subsection{Errores frecuentes}
	
	\necesitaespacio{10\baselineskip}
	\begin{errores}
		\texttt{WARNING: The convert command is deprecated} & Uso del nombre heredado. & Use \cmd{magick}.\\*
		\texttt{cache resources exhausted} & Se superaron los límites de memoria o disco de la política. & Reduzca la resolución o el tamaño; procese por páginas; use Poppler para PDF.\\
		\texttt{width or height exceeds limit} & Imagen mayor de 32\,000 píxeles de lado. & Reduzca \opt{-density} o divida la imagen.\\
		\texttt{not allowed by the security policy} & Operación prohibida por \arch{policy.xml} (URL, \arch{@archivo}\ldots). & Use otra vía (descargue con \cmd{wget}, pase los archivos directamente).\\
		\texttt{no images defined} & No se leyó ninguna entrada (ruta incorrecta o comodín sin coincidencias). & Revise la ruta con \cmd{ls}.\\
		\texttt{no decode delegate for this image format} & Formato no soportado por la compilación, o extensión incorrecta. & \cmd{magick -list format}; compruebe el tipo real con \cmd{file archivo}.\\
		\texttt{FailedToExecuteCommand `gs'} & Falta Ghostscript para leer PDF/PS. & \cmd{sudo apt install ghostscript}.\\
		JPEG con fondo negro & Transparencia sin aplanar. & \opt{-background white -alpha remove -alpha off}.\\
		PDF rasterizado borroso & Faltó \opt{-density} antes del PDF. & \opt{-density 300} antes del nombre del PDF.\\*
		Originales modificados & \cmd{mogrify} sin \opt{-path}. & Recupere de una copia; use siempre \opt{-path}.\\
	\end{errores}
	
	Ayuda: \cmd{magick -help}, \cmd{magick -list option}, \cmd{man magick} y la referencia de opciones en \url{https://imagemagick.org/script/command-line-options.php}. Versión: \cmd{magick -version}.
	% =====================================================================
	\clearpage
	\section{Potrace}\label{sec:potrace}
	% =====================================================================
	
	\necesitaespacio{8\baselineskip}
	\subsection{Descripción general}
	
	\textbf{Qué es.} Un programa de \emph{vectorización}: convierte un mapa de bits en blanco y negro en contornos formados por curvas suaves (curvas de Bézier).
	
	\textbf{Raster frente a vector.} Una imagen \emph{raster} (PNG, JPEG) es una cuadrícula de píxeles: al ampliarla aparecen escalones. Una imagen \emph{vectorial} (SVG, PDF, EPS) describe formas con ecuaciones: se amplía sin pérdida. Potrace calcula las curvas que mejor siguen el borde entre las zonas negras y blancas.
	
	\textbf{Cuándo usarlo:} logotipos, firmas, sellos, iconos, dibujos a línea, tipografías escaneadas, planos sencillos: imágenes de dos colores con bordes nítidos.
	
	\textbf{Cuándo no usarlo:} fotografías o degradados (el resultado será una mancha bicolor); imágenes de muchos colores (Potrace traza un solo color por pasada); reconocimiento de texto: \textbf{Potrace no es un OCR} (convierte las letras en dibujos, no en texto editable).
	
	\necesitaespacio{8\baselineskip}
	\subsection{Instalación y comprobación}
	
	\begin{lstlisting}
	sudo apt install potrace    # incluye también mkbitmap
	potrace --version
	mkbitmap --version
\end{lstlisting}
	
	\necesitaespacio{8\baselineskip}
	\subsection{Sintaxis}
	
	\begin{lstlisting}[style=sintaxis]
  potrace [opciones] [<archivo.pbm|pgm|ppm|bmp> ...]
\end{lstlisting}
	
	Sin \opt{-o}, cada salida toma el nombre de la entrada con la extensión del formato elegido. Sin archivo de entrada (o con \val{-}), lee de la entrada estándar.
	
	\flujoh{logo.pbm (píxeles),potrace -s,logo.svg (curvas)}
	
	\necesitaespacio{25\baselineskip}
	\subsection{Formatos de salida (\emph{backends})}
	
	\necesitaespacio{21\baselineskip}
	\begin{opciones}[Opción]
		\opt{-s}, \opt{-b svg} & SVG: para la web, Inkscape, LaTeX (con \texttt{svg}) o edición posterior.\\*
		\opt{-b pdf} & PDF del tamaño del dibujo (ideal para \texttt{\bs{}\allowbreak{}includegraphics} en LaTeX).\\
		\opt{-b pdfpage} & PDF de página completa, con el tamaño de \opt{-P}.\\
		\opt{-e}, \opt{-b eps} & EPS (predeterminado).\\
		\opt{-p}, \opt{-b ps} & PostScript de página completa.\\
		\opt{-g}, \opt{-b pgm} & PGM: rasteriza de nuevo las curvas suavizadas (útil para escalar un mapa de bits sin escalones).\\
		\opt{-b dxf} & DXF para programas de CAD.\\
		\opt{-b geojson} & GeoJSON.\\
		\opt{-b gimppath} & Trazado para GIMP.\\*
		\opt{-b xfig} & Formato de XFig.\\
	\end{opciones}
	
	\necesitaespacio{14\baselineskip}
	\subsection{Parámetros principales}
	
	\necesitaespacio{10\baselineskip}
	\begin{opciones}
		\opt{-o ARCHIVO} & Archivo de salida (\val{-} para la salida estándar).\\*
		\opt{-t N}, \opt{-{}-turdsize N} & Elimina manchas de hasta \val{N} píxeles (predeterminado 2). Súbalo para limpiar el polvo de un escaneo.\\
		\opt{-a N}, \opt{-{}-alphamax N} & Umbral de esquinas (predeterminado 1): valores bajos producen más esquinas angulosas; altos, curvas más suaves (hasta 1,3334); \val{0} da un polígono.\\
		\opt{-n}, \opt{-{}-longcurve} & Desactiva la optimización de curvas (más nodos, más fidelidad).\\
		\opt{-O N}, \opt{-{}-opttolerance N} & Tolerancia de esa optimización (predeterminada 0,2): más alta, menos nodos.\\
		\opt{-z P}, \opt{-{}-turnpolicy P} & Cómo resolver ambigüedades en el contorno: \val{black}, \val{white}, \val{left}, \val{right}, \val{minority}, \val{majority}, \val{random}.\\
		\opt{-k N}, \opt{-{}-blacklevel N} & Umbral de negro para entradas en gris o color (0--1, predeterminado 0,5).\\
		\opt{-i}, \opt{-{}-invert} & Invierte el mapa de bits (traza el blanco en vez del negro).\\
		\opt{-C \#rrggbb} & Color de las formas trazadas (predeterminado negro).\\
		\opt{-{}-fillcolor \#rrggbb} & Color de relleno de los huecos interiores (predeterminado transparente).\\
		\opt{-{}-opaque} & Hace opacas las formas blancas.\\
		\opt{-{}-group}, \opt{-{}-flat} & Agrupa trazados relacionados, o todo en un único trazado (SVG).\\
		\opt{-W}, \opt{-H} & Ancho o alto del resultado con unidades: \val{5cm}, \val{2in}, \val{100pt}.\\
		\opt{-r N} & Resolución en ppp para formatos con dimensiones físicas.\\
		\opt{-x N} & Escala para formatos de píxeles (SVG, PGM, GeoJSON).\\
		\opt{-A ÁNGULO} & Gira el resultado en sentido antihorario.\\
		\opt{-M D}, \opt{-L}, \opt{-R}, \opt{-T}, \opt{-B} & Márgenes (todos o izquierdo, derecho, superior, inferior).\\
		\opt{-P T}, \opt{-{}-pagesize T} & Tamaño de página para \val{ps} y \val{pdfpage} (predeterminado \val{letter}; por ejemplo \val{a4}).\\
		\opt{-{}-tight} & Elimina el espacio en blanco alrededor del dibujo.\\
		\opt{-u N} & Cuantiza las coordenadas a 1/\val{N} píxel (predeterminado 10).\\*
		\opt{-c}, \opt{-2}, \opt{-3} & Compresión del PostScript/EPS: sin compresión, nivel 2 (predeterminado) o nivel 3.\\
	\end{opciones}
	
	\necesitaespacio{8\baselineskip}
	\subsection{Preparar la imagen: por qué blanco y negro}
	
	Potrace solo entiende mapas de bits PBM, PGM, PPM y BMP; un PNG o JPEG produce \texttt{file format not recognized}. Además, traza la frontera entre dos colores: si la entrada tiene grises, sombras o artefactos de JPEG, el umbral \opt{-k} decide de forma ciega qué es negro. Por eso conviene preparar la imagen con ImageMagick:
	
	\flujo{logo.png / foto.jpg,magick: gris + umbral,logo.pbm (blanco y negro puro),potrace,logo.svg}
	
	\begin{lstlisting}
	# 1) Convertir a gris y aplicar un umbral -> PBM (1 bit)
	magick logo.png -colorspace Gray -threshold 50% logo.pbm
	
	# 2) Vectorizar
	potrace logo.pbm -s -o logo.svg
	
	# En una sola línea, sin archivo intermedio
	magick logo.png -colorspace Gray -threshold 50% pbm:- | potrace -s -o logo.svg
\end{lstlisting}
	
	Ajuste el umbral (40\,\%, 60\,\%\ldots) hasta que la vista previa en blanco y negro contenga todos los detalles sin manchas. Para dibujos escaneados con iluminación desigual, \cmd{mkbitmap} da mejores resultados que un umbral fijo.
	
	\necesitaespacio{18\baselineskip}
	\subsubsection{\texttt{mkbitmap}: preprocesador de Potrace}
	
	\cmd{mkbitmap} aplica un filtro de paso alto (compensa fondos irregulares), amplía la imagen con interpolación (bordes más suaves) y aplica el umbral.
	
	\necesitaespacio{14\baselineskip}
	\begin{opciones}
		\opt{-f N} & Filtro de paso alto de radio \val{N} (predeterminado 4). \opt{-n} lo desactiva.\\*
		\opt{-s N} & Escala por un factor entero (predeterminado 2).\\
		\opt{-3}, \opt{-1} & Interpolación cúbica (predeterminada) o lineal.\\
		\opt{-t N} & Umbral (0--1, predeterminado 0,45).\\
		\opt{-b N} & Desenfoque previo de radio \val{N} (reduce ruido).\\
		\opt{-i} & Invierte (para dibujos con tiza sobre pizarra).\\*
		\opt{-g} & Salida en gris, sin umbral.\\
	\end{opciones}
	
	\begin{lstlisting}
	magick dibujo_escaneado.jpg -colorspace Gray dibujo.pgm
	mkbitmap -f 2 -s 2 -t 0.48 dibujo.pgm -o dibujo.pbm
	potrace -s -t 5 dibujo.pbm -o dibujo.svg
\end{lstlisting}
	
	\necesitaespacio{8\baselineskip}
	\subsection{Ejemplos}
	
	\begin{lstlisting}
	potrace logo.pbm -s -o logo.svg                     # SVG
	potrace -b pdf logo.pbm -o logo.pdf                 # PDF vectorial para LaTeX
	potrace -b pdfpage -P a4 logo.pbm -o logo_a4.pdf    # en una página A4
	potrace logo.pbm -s -t 5 -a 1.0 -O 0.2 -o limpio.svg   # menos manchas
	potrace logo.pbm -s -C "#1f4e79" -o azul.svg        # color corporativo
	potrace logo.pbm -s -W 5cm -o 5cm.svg               # ancho de 5 cm
	potrace firma.bmp -s -o firma.svg                   # BMP también se admite
\end{lstlisting}
	
	\paragraph{Imágenes de varios colores} Potrace no separa colores. La técnica consiste en crear una máscara en blanco y negro por cada color con ImageMagick (\opt{-fuzz} y \opt{-fill}/\opt{-opaque}, o umbrales por canal), vectorizar cada máscara con su \opt{-C} y combinar los SVG en un editor como Inkscape. Es un proceso manual: para logotipos de muchos colores, conviene un programa de diseño.
	
	\necesitaespacio{21\baselineskip}
	\subsection{Formatos de entrada y salida}
	
	\necesitaespacio{17\baselineskip}
	\begin{formatos}
		PBM & \compat & \nocompat & Formato ideal: 1 bit, sin ambigüedad.\\*
		PGM, PPM & \compat & \nocompat & Se umbralizan con \opt{-k} (0,5 por defecto).\\
		BMP & \compat & \nocompat & Admitido directamente.\\
		PNG, JPEG, TIFF & \nocompat & \nocompat & Convertir antes con ImageMagick.\\
		SVG, PDF, EPS, PS & \nocompat & \compat & Salidas vectoriales.\\
		DXF, GeoJSON, GIMP, XFig & \nocompat & \compat & Salidas para otros programas.\\*
		PGM & \compat & \compat & También como salida rasterizada suavizada.\\
	\end{formatos}
	
	\necesitaespacio{8\baselineskip}
	\subsection{Limitaciones}
	
	\begin{limitacion}
		\begin{itemize}
			\item Un solo color por pasada; no vectoriza degradados, sombras ni fotografías.
			\item El resultado depende por completo del umbral: un mal preprocesado da un mal vector.
			\item No es OCR: el texto se convierte en formas, no en caracteres.
			\item Las imágenes muy pequeñas o de baja resolución producen contornos pobres: amplíe con \cmd{mkbitmap -s} o parta de un escaneo de 300--600 ppp.
		\end{itemize}
	\end{limitacion}
	
	\necesitaespacio{21\baselineskip}
	\subsection{Errores frecuentes}
	
	\necesitaespacio{17\baselineskip}
	\begin{errores}
		\texttt{file format not recognized} & Entrada PNG/JPEG. & Convierta a PBM/BMP con \cmd{magick}.\\*
		SVG negro o vacío & Umbral inadecuado o imagen invertida. & Ajuste \opt{-threshold} o \opt{-k}; pruebe \opt{-i}.\\
		Muchos puntos sueltos & Ruido del escaneo. & Suba \opt{-t}; limpie antes con \cmd{mkbitmap} o \opt{-despeckle}.\\
		Curvas «blandas» en esquinas & \opt{-a} alto. & Baje \opt{-a} (por ejemplo 0,5).\\*
		Tamaño físico incorrecto & Resolución no indicada. & Use \opt{-W}/\opt{-H} o \opt{-r}.\\
	\end{errores}
	
	Ayuda: \cmd{potrace -{}-help}, \cmd{mkbitmap -{}-help}, \cmd{man potrace}. Versión: \cmd{potrace -{}-version}.
	
	% =====================================================================
	\clearpage
	\section{Pandoc}\label{sec:pandoc}
	% =====================================================================
	
	\necesitaespacio{8\baselineskip}
	\subsection{Descripción general}
	
	\textbf{Qué es.} Un conversor universal de documentos. Lee un formato, lo transforma en una representación interna (un árbol de bloques: títulos, párrafos, listas, tablas, enlaces\ldots) y la escribe en otro formato.
	
	\textbf{Qué problema resuelve.} Escribir una vez (normalmente en Markdown) y publicar en DOCX, PDF, HTML, EPUB, ODT o LaTeX; o convertir documentos existentes (DOCX, HTML) a Markdown.
	
	\textbf{Cuándo usarlo:} documentos de texto estructurado: apuntes, informes, artículos, libros, documentación técnica.
	
	\textbf{Cuándo no usarlo:} \textbf{Pandoc no lee PDF} (para PDF $\rightarrow$ Markdown, use Marker); tampoco conserva maquetaciones complejas (cuadros de texto flotantes, columnas manuales de Word, diseños gráficos).
	
	\necesitaespacio{8\baselineskip}
	\subsection{Instalación y comprobación}
	
	\begin{lstlisting}
	sudo apt install pandoc
	pandoc --version
	pandoc --list-input-formats      # formatos que puede leer
	pandoc --list-output-formats     # formatos que puede escribir
	# variantes de Markdown (+ activada, - desactivada)
	pandoc --list-extensions=markdown
\end{lstlisting}
	
	Para generar PDF, Pandoc necesita un \textbf{motor} externo; el predeterminado es \cmd{pdflatex} (paquetes LaTeX de la sección~\ref{sec:instalacion}).
	
	\necesitaespacio{8\baselineskip}
	\subsection{Sintaxis}
	
	\begin{lstlisting}[style=sintaxis]
  pandoc [opciones] [<entrada> ...] -o <salida>
  pandoc -f <formato_entrada> -t <formato_salida> [opciones] [<entrada> ...]
\end{lstlisting}
	
	Pandoc deduce los formatos por las extensiones (\arch{.md}, \arch{.docx}, \arch{.html}\ldots). \opt{-f} y \opt{-t} los fijan explícitamente (imprescindible al leer de la entrada estándar o con extensiones ambiguas). Si se dan varias entradas, se concatenan. Sin \opt{-o}, el resultado se escribe en pantalla.
	
	\flujoh{informe.md,pandoc,informe.docx / .pdf / .html}
	
	\necesitaespacio{14\baselineskip}
	\subsection{Opciones principales}
	
	\necesitaespacio{10\baselineskip}
	\begin{opciones}
		\opt{-f F}, \opt{-{}-from=F} & Formato de entrada: \val{markdown}, \val{gfm} (Markdown de GitHub), \val{docx}, \val{html}, \val{latex}, \val{odt}, \val{epub}, \val{rst}\ldots{} Admite extensiones: \val{markdown+hard\_\allowbreak{}line\_\allowbreak{}breaks}.\\*
		\opt{-t F}, \opt{-{}-to=F} & Formato de salida: \val{docx}, \val{html}, \val{latex}, \val{pdf}, \val{gfm}, \val{plain}, \val{odt}, \val{epub}, \val{beamer}, \val{pptx}\ldots\\
		\opt{-o ARCHIVO} & Archivo de salida.\\
		\opt{-s}, \opt{-{}-standalone} & Documento completo (cabecera HTML, preámbulo LaTeX). Automático para DOCX, ODT, EPUB y PDF; necesario para HTML o LaTeX autónomos.\\
		\opt{-{}-toc}, \opt{-{}-toc-depth=N} & Índice de contenidos y su profundidad.\\
		\opt{-N}, \opt{-{}-number-\allowbreak{}sections} & Numera las secciones.\\
		\opt{-M CLAVE=VALOR}, \opt{-{}-metadata} & Metadatos del documento (\val{title}, \val{author}, \val{date}, \val{lang}\ldots).\\
		\opt{-{}-metadata-\allowbreak{}file=\allowbreak{}ARCHIVO} & Metadatos desde un archivo YAML.\\
		\opt{-V CLAVE=VALOR}, \opt{-{}-variable} & Variables de la plantilla (márgenes, fuentes, tamaño de letra\ldots).\\
		\opt{-{}-template=\allowbreak{}ARCHIVO} & Plantilla propia (para HTML, LaTeX\ldots).\\
		\opt{-D F}, \opt{-{}-print-\allowbreak{}default-\allowbreak{}template=\allowbreak{}F} & Muestra la plantilla predeterminada (punto de partida para personalizar).\\
		\opt{-{}-pdf-\allowbreak{}engine=\allowbreak{}MOTOR} & Motor para PDF: \val{pdflatex}, \val{xelatex}, \val{lualatex}, \val{latexmk}, \val{tectonic}, \val{typst}, \val{wkhtmltopdf}, \val{weasyprint}, \val{pagedjs-cli}, \val{prince}, \val{context}, \val{groff}, \val{pdfroff}.\\
		\opt{-{}-pdf-\allowbreak{}engine-\allowbreak{}opt=\allowbreak{}OPC} & Opción adicional para el motor.\\
		\opt{-{}-reference-\allowbreak{}doc=\allowbreak{}ARCHIVO} & DOCX/ODT/PPTX de referencia del que se copian los estilos.\\
		\opt{-{}-print-\allowbreak{}default-\allowbreak{}data-\allowbreak{}file=\allowbreak{}ARCH} & Extrae archivos internos, como \val{reference.docx}, para personalizarlos.\\
		\opt{-{}-extract-\allowbreak{}media=\allowbreak{}CARPETA} & Extrae las imágenes al convertir desde DOCX, ODT o EPUB.\\
		\opt{-{}-resource-\allowbreak{}path=\allowbreak{}RUTAS} & Carpetas donde buscar imágenes y recursos (separadas por \val{:}).\\
		\opt{-{}-embed-\allowbreak{}resources} & Con \opt{-s}, incrusta imágenes, CSS y scripts en un único HTML (sustituye a la antigua \opt{-{}-self-\allowbreak{}contained}).\\
		\opt{-{}-highlight-\allowbreak{}style=\allowbreak{}ESTILO} & Estilo del resaltado de código: \val{pygments}, \val{tango}, \val{kate}, \val{monochrome}, \val{espresso}, \val{zenburn}, \val{haddock}, \val{breezedark}. \opt{-{}-no-highlight} lo desactiva.\\
		\opt{-{}-mathjax}, \opt{-{}-mathml}, \opt{-{}-katex} & Cómo se muestran las fórmulas en HTML.\\
		\opt{-{}-lof}, \opt{-{}-lot} & Índices de figuras y de tablas (salidas LaTeX/PDF).\\
		\opt{-C}, \opt{-{}-citeproc} & Procesa citas con \opt{-{}-bibliography=\allowbreak{}refs.bib} y \opt{-{}-csl=\allowbreak{}estilo.csl}.\\
		\opt{-L F}, \opt{-{}-lua-filter=F} & Filtro Lua para transformar el documento.\\
		\opt{-{}-wrap=\allowbreak{}auto|\allowbreak{}none|\allowbreak{}preserve} & Ajuste de líneas en salidas de texto; \val{none} evita cortes artificiales.\\
		\opt{-{}-shift-\allowbreak{}heading-\allowbreak{}level-\allowbreak{}by=\allowbreak{}N} & Sube o baja el nivel de todos los títulos.\\
		\opt{-{}-top-\allowbreak{}level-\allowbreak{}division=\allowbreak{}D} & En LaTeX/PDF, qué es un título de nivel 1: \val{section}, \val{chapter} o \val{part}.\\
		\opt{-{}-track-\allowbreak{}changes=\allowbreak{}M} & Al leer DOCX con control de cambios: \val{accept}, \val{reject} o \val{all}.\\
		\opt{-d ARCHIVO}, \opt{-{}-defaults} & Archivo YAML con todas las opciones (configuración reutilizable).\\*
		\opt{-{}-verbose}, \opt{-{}-fail-\allowbreak{}if-\allowbreak{}warnings} & Diagnóstico detallado; tratar los avisos como errores.\\
	\end{opciones}
	
	\begin{nota}[Opciones de versiones posteriores]
		Pandoc 3.7 no tiene \opt{-{}-syntax-\allowbreak{}highlighting} (aparece en 3.8) ni \opt{-{}-math-method} (3.11). Si las ve en Internet, use sus equivalentes en 3.7: \opt{-{}-highlight-\allowbreak{}style} y \opt{-{}-mathjax}/\opt{-{}-mathml}/\opt{-{}-katex}.
	\end{nota}
	
	\necesitaespacio{8\baselineskip}
	\subsection{Conceptos clave}
	
	\paragraph{Conversión directa} DOCX, HTML, ODT o EPUB se generan directamente: no hace falta nada más que Pandoc.
	
	\paragraph{Generar LaTeX} \cmd{pandoc doc.md -s -o doc.tex} produce un archivo \arch{.tex} completo que puede revisar, retocar y compilar usted mismo. Es la vía para controlar el resultado al detalle.
	
	\paragraph{Generar PDF} \cmd{pandoc doc.md -o doc.pdf} genera LaTeX internamente, llama al motor y borra los intermedios. Si el motor falla, Pandoc muestra el error de LaTeX.
	
	\paragraph{Motores} \cmd{pdflatex} es el más extendido, pero solo admite los caracteres Unicode que tiene configurados: símbolos como $\approx$, $\leq$ o subíndices Unicode escritos directamente provocan \texttt{Unicode character ... not set up for use with LaTeX}. \cmd{xelatex} y \cmd{lualatex} admiten Unicode y fuentes del sistema (\opt{-V mainfont="DejaVu Serif"}), pero si la fuente no tiene un carácter, lo \emph{omiten} con un aviso \texttt{Missing character} en lugar de detenerse: revise los avisos y elija una fuente con buena cobertura.
	
	\paragraph{Metadatos y variables} Los metadatos (\opt{-M} o el bloque YAML inicial) describen el documento: título, autor, fecha, idioma. Las variables (\opt{-V}) solo configuran la plantilla: márgenes, fuente, tamaño. Con \val{lang: es}, LaTeX usa la separación silábica y los rótulos en español («Tabla», «Figura»).
	
	\begin{lstlisting}[style=salida]
	---
	title: "Informe de laboratorio"
	author: "Nombre Apellido"
	date: "Octubre de 2026"
	lang: es
	geometry: margin=2.5cm
	fontsize: 11pt
	---
	# Introducción
	Texto del documento...
\end{lstlisting}
	
	\paragraph{Plantillas y documentos de referencia} Para DOCX, el aspecto se controla con un documento de referencia: extraiga el predeterminado, modifique sus \emph{estilos} en Word (no el texto) y úselo con \opt{-{}-reference-doc}.
	
	\necesitaespacio{8\baselineskip}
	\subsection{Ejemplos}
	
	\begin{lstlisting}
	# Markdown -> DOCX, con estilos propios
	pandoc informe.md -o informe.docx
	# editar estilos en Word
	pandoc -o plantilla.docx --print-default-data-file reference.docx
	pandoc informe.md --reference-doc=plantilla.docx -o informe.docx
	
	# Markdown -> PDF (pdflatex) con índice, secciones numeradas y márgenes
	pandoc informe.md -o informe.pdf --toc -N -M lang=es -V geometry:margin=2.5cm
	
	# Markdown -> PDF con Unicode completo y fuente del sistema
	pandoc informe.md -o informe.pdf --pdf-engine=xelatex -M lang=es -V \
	mainfont="DejaVu Serif"
	
	# Markdown -> LaTeX completo para revisar y compilar a mano
	pandoc informe.md -s -o informe.tex
	
	# Markdown -> HTML autónomo (imágenes y estilos incrustados)
	pandoc informe.md -s --embed-resources --toc -o informe.html
	
	# DOCX -> Markdown (GitHub), extrayendo las imágenes
	pandoc documento.docx -t gfm --extract-media=imagenes --wrap=none -o documento.md
	
	# HTML -> Markdown
	pandoc pagina.html -f html -t gfm -o pagina.md
	
	# Markdown -> EPUB y ODT
	pandoc libro.md -o libro.epub -M title="Mi libro" -M author="Autor"
	pandoc informe.md -o informe.odt
	
	# Markdown -> texto plano
	pandoc informe.md -t plain -o informe.txt
\end{lstlisting}
	
	\begin{lstlisting}
	# Configuración reutilizable con un archivo de valores por defecto
	cat > pdf-es.yaml <<'EOF'
	from: markdown
	to: pdf
	pdf-engine: xelatex
	metadata:
	lang: es
	variables:
	geometry: margin=2.5cm
	mainfont: DejaVu Serif
	EOF
	pandoc -d pdf-es.yaml informe.md -o informe.pdf
	
	# Lote: todos los Markdown de la carpeta a DOCX
	for f in *.md; do pandoc "$f" -o "${f%.md}.docx"; done
\end{lstlisting}
	
	\necesitaespacio{32\baselineskip}
	\subsection{Formatos de entrada y salida}
	
	\necesitaespacio{28\baselineskip}
	\begin{formatos}
		Markdown (\val{markdown}, \val{gfm}, \val{commonmark}) & \compat & \compat & Formato de trabajo recomendado.\\*
		DOCX & \compat & \compat & Estilos vía \opt{-{}-reference-\allowbreak{}doc}; imágenes con \opt{-{}-extract-\allowbreak{}media}.\\
		ODT & \compat & \compat & Equivalente para LibreOffice.\\
		HTML & \compat & \compat & Al leer se pierde el CSS; al escribir, \opt{-s} para documento completo.\\
		LaTeX & \compat & \compat & Al leer, solo entiende macros estándar.\\
		PDF & \nocompat & \condic & Solo salida, con un motor externo (LaTeX, Typst, HTML\ldots).\\
		EPUB & \compat & \compat & Libros electrónicos.\\
		TXT (\val{plain}) & \nocompat & \compat & Texto sin formato.\\
		PPTX & \nocompat & \compat & Presentaciones sencillas (también \val{beamer}, \val{revealjs}).\\*
		RST, Org, Typst, Jupyter (\val{ipynb}), CSV/TSV\ldots & \condic & \condic & Ver \cmd{-{}-list-\allowbreak{}input-\allowbreak{}formats} y \cmd{-{}-list-\allowbreak{}output-\allowbreak{}formats}: CSV y TSV solo de entrada.\\
	\end{formatos}
	
	\necesitaespacio{8\baselineskip}
	\subsection{Limitaciones}
	
	\begin{limitacion}
		\begin{itemize}
			\item \textbf{No lee PDF} ni imágenes.
			\item \textbf{Fidelidad}: Pandoc conserva la \emph{estructura}, no el \emph{diseño}. Cuadros de texto, columnas manuales, imágenes flotantes, encabezados de Word o CSS se pierden o se simplifican.
			\item \textbf{Tablas complejas}: celdas combinadas y tablas anidadas pueden simplificarse según el formato de destino.
			\item \textbf{LaTeX de entrada}: macros propias o paquetes poco comunes no se interpretan.
			\item \textbf{PDF}: depende del motor instalado; los errores suelen ser de LaTeX (paquete ausente, carácter Unicode).
			\item \textbf{Ida y vuelta}: DOCX $\rightarrow$ Markdown $\rightarrow$ DOCX no devuelve el documento original idéntico.
		\end{itemize}
	\end{limitacion}
	
	\necesitaespacio{14\baselineskip}
	\subsection{Errores frecuentes}
	
	\necesitaespacio{10\baselineskip}
	\begin{errores}
		\texttt{pdflatex not found. Please select a different -{}-pdf-\allowbreak{}engine or install pdflatex} & No hay LaTeX instalado. & Instale los paquetes \texttt{texlive-*} de la sección~\ref{sec:instalacion}.\\*
		\texttt{LaTeX Error: File `xxx.sty' not found} & Falta un paquete LaTeX. & \cmd{sudo apt install texlive-\allowbreak{}latex-\allowbreak{}extra} (o el paquete que lo contenga).\\
		\texttt{Unicode character ... not set up for use with LaTeX} & Carácter no admitido por \cmd{pdflatex}. & \opt{-{}-pdf-engine=xelatex} con una fuente adecuada, o escriba el símbolo como fórmula (\verb|$\approx$|).\\
		Símbolos ausentes en el PDF sin error & Fuente sin esos glifos (xelatex/lualatex). & Revise los avisos \texttt{Missing character}; cambie \opt{mainfont}.\\
		\texttt{Unknown option -{}-syntax-\allowbreak{}highlighting} & Opción de una versión posterior. & \opt{-{}-highlight-\allowbreak{}style}.\\
		\texttt{Could not fetch resource} & Imagen no encontrada. & Rutas relativas al lugar donde ejecuta Pandoc, o \opt{-{}-resource-\allowbreak{}path}.\\
		DOCX $\rightarrow$ Markdown sin imágenes & Falta \opt{-{}-extract-\allowbreak{}media}. & Añada \opt{-{}-extract-\allowbreak{}media=\allowbreak{}carpeta}.\\
		\texttt{YAML parse exception} & Bloque YAML mal formado (dos puntos o comillas). & Ponga los valores con \texttt{:} entre comillas.\\*
		Títulos o listas no reconocidos & Falta una línea en blanco antes del bloque. & Deje una línea en blanco antes de títulos y listas.\\
	\end{errores}
	
	Ayuda: \cmd{pandoc -{}-help}, \cmd{man pandoc} y el manual en \url{https://pandoc.org/MANUAL.html} (describe la última versión: compare con \cmd{pandoc -{}-version}).
	% =====================================================================
	\clearpage
	\section{QPDF}\label{sec:qpdf}
	% =====================================================================
	
	\necesitaespacio{8\baselineskip}
	\subsection{Descripción general}
	
	\textbf{Qué es.} Un programa que transforma archivos PDF \emph{a nivel estructural}: trabaja con los objetos internos del PDF (páginas, diccionarios, flujos, tablas de referencias) sin interpretar ni volver a dibujar su contenido. Por eso sus operaciones conservan intactos el texto, las fuentes y las imágenes.
	
	\textbf{Qué problema resuelve.} Seleccionar, reordenar, unir, dividir y rotar páginas; cifrar y descifrar; linealizar para la web; reparar tablas de referencias dañadas; comprimir sin pérdida; inspeccionar la estructura (incluso en JSON).
	
	\textbf{Diferencias con las otras herramientas de PDF.}
	\begin{itemize}
		\item \textbf{Poppler} \emph{lee y renderiza}: extrae texto, convierte páginas en imágenes. QPDF no dibuja nada.
		\item \textbf{Ghostscript} \emph{reinterpreta} el documento y escribe uno nuevo: puede recomprimir imágenes y cambiar fuentes. QPDF copia los objetos tal cual.
		\item \textbf{OCRmyPDF} añade una capa de texto reconocido. \textbf{QPDF no es una herramienta de OCR} ni modifica el contenido de las páginas.
	\end{itemize}
	
	\necesitaespacio{8\baselineskip}
	\subsection{Instalación y comprobación}
	
	\begin{lstlisting}
	sudo apt install qpdf
	qpdf --version
	qpdf --help            # temas de ayuda; qpdf --help=all muestra todo
	qpdf --help=--pages    # ayuda de una opción concreta
\end{lstlisting}
	
	\necesitaespacio{8\baselineskip}
	\subsection{Sintaxis general y códigos de salida}
	
	\begin{lstlisting}[style=sintaxis]
  qpdf [opciones] <entrada.pdf> [opciones] <salida.pdf>
  qpdf [opciones] --empty --pages <archivo> [rango] [<archivo> [rango] ...] -- \
  <salida.pdf>
  qpdf [opciones] <entrada.pdf> --replace-input
\end{lstlisting}
	
	QPDF siempre escribe un archivo \emph{nuevo}; no acepta la misma ruta como entrada y salida. Para modificar el archivo en su sitio existe \opt{-{}-replace-\allowbreak{}input}.
	
	\begingroup\small
	\necesitaespacio{9\baselineskip}
	\begin{longtable}{C{1.6cm}L{13.46cm}}
		\toprule\rowcolor{tablacab}\cab{Código} & \cab{Significado (\texttt{echo \$?} tras la orden)}\\\midrule
		\endfirsthead
		\bottomrule
		\endfoot
		0 & Sin errores ni avisos.\\*
		2 & Errores: el archivo de salida puede no haberse creado.\\*
		3 & Avisos: el resultado existe, pero QPDF encontró y corrigió problemas (salvo con \opt{-{}-warning-\allowbreak{}exit-\allowbreak{}0}, que devuelve 0).\\
	\end{longtable}
	\endgroup
	
	\necesitaespacio{22\baselineskip}
	\subsubsection{Sintaxis de rangos de páginas}\label{sec:qpdf-rangos}
	
	\necesitaespacio{18\baselineskip}
	\begin{opciones}[Rango]
		\val{1,3,5} & Páginas 1, 3 y 5.\\*
		\val{2-8} & De la 2 a la 8.\\
		\val{z} & Última página (equivale a \val{r1}).\\
		\val{r3} & Tercera página contando desde el final.\\
		\val{8-2} & Orden inverso (de la 8 a la 2); \val{z-1} invierte todo el documento.\\
		\val{1-z,x3-4} & Todas excepto la 3 y la 4 (\val{x} excluye del grupo anterior).\\
		\val{1-z:odd}, \val{1-z:even} & Impares o pares \emph{del conjunto resultante}.\\*
		\val{.} (como archivo) & En \opt{-{}-pages}, el punto representa el archivo de entrada principal.\\
	\end{opciones}
	
	\necesitaespacio{32\baselineskip}
	\subsection{Inspección y validación}\label{sec:qpdf-inspeccion}
	
	\necesitaespacio{28\baselineskip}
	\begin{opciones}
		\opt{-{}-check} & Comprueba la estructura y el contenido (decodifica todos los flujos). Informa de errores y avisos.\\*
		\opt{-{}-show-npages} & Número de páginas (útil en scripts).\\
		\opt{-{}-show-pages} & Objeto de cada página, su contenido y sus imágenes.\\
		\opt{-{}-show-\allowbreak{}encryption} & Tipo de cifrado y permisos.\\
		\opt{-{}-is-encrypted} & Solo código de salida: 0 si está cifrado, 2 si no.\\
		\opt{-{}-requires-\allowbreak{}password} & Solo código de salida: 0 si hace falta contraseña para abrirlo, 3 si está cifrado pero se abre sin ella, 2 si no está cifrado.\\
		\opt{-{}-json} & Descripción completa en JSON. Con \opt{-{}-json-\allowbreak{}key=\allowbreak{}pages} (u \val{outlines}, \val{encrypt}, \val{attachments}\ldots) limita las secciones.\\
		\opt{-{}-json-output} & Escribe el PDF completo en formato JSON de QPDF (editable y reconvertible con \opt{-{}-json-input}).\\
		\opt{-{}-show-\allowbreak{}object=\allowbreak{}N} & Muestra el objeto número \val{N} (\opt{-{}-raw-\allowbreak{}stream-\allowbreak{}data} o \opt{-{}-filtered-\allowbreak{}stream-\allowbreak{}data} para su flujo).\\
		\opt{-{}-show-xref} & Tabla de referencias cruzadas.\\
		\opt{-{}-check-\allowbreak{}linearization} & Verifica la linealización.\\*
		\opt{-{}-list-\allowbreak{}attachments} & Archivos adjuntos.\\
	\end{opciones}
	
	\begin{lstlisting}
	qpdf --check doc.pdf; echo "codigo: $?"
	qpdf --show-npages doc.pdf
	qpdf --show-encryption doc.pdf
	qpdf --requires-password doc.pdf; echo $?     # 0 = pide contraseña al abrir
	qpdf --json --json-key=pages doc.pdf | less
\end{lstlisting}
	
	\necesitaespacio{8\baselineskip}
	\subsection{Reparación}
	
	Muchos PDF dañados (descargas incompletas, tablas de referencias corruptas) se reparan simplemente reescribiéndolos: QPDF reconstruye la estructura al leer y escribe un archivo coherente.
	
	\begin{lstlisting}
		qpdf danado.pdf reparado.pdf; echo $?      # 3 = reparado con avisos
		qpdf --check reparado.pdf
\end{lstlisting}
	
	Para estudiar la estructura con un editor de texto, el modo QDF escribe un PDF «legible»:
	
	\begin{lstlisting}
		qpdf --qdf --object-streams=disable doc.pdf legible.pdf
\end{lstlisting}
	
	\begin{limitacion}
		QPDF repara la \emph{estructura}, no el contenido: si faltan datos de una imagen o de una página porque el archivo está truncado, no puede recuperarlos. Si \cmd{qpdf} no lo arregla, pruebe a reescribir con Ghostscript (que reinterpreta el contenido) o con \cmd{pdftocairo -pdf}.
	\end{limitacion}
	
	\necesitaespacio{30\baselineskip}
	\subsection{Páginas: seleccionar, unir, dividir y rotar}\label{sec:qpdf-paginas}
	
	\necesitaespacio{26\baselineskip}
	\begin{opciones}
		\opt{-{}-pages A [R] B [R] ... -{}-} & Construye el documento con las páginas indicadas de cada archivo. Termina con \opt{-{}-}. Si un archivo está cifrado: \val{archivo -{}-password=\allowbreak{}CLAVE rango}.\\*
		\opt{-{}-empty} & Parte de un PDF vacío en lugar de un archivo de entrada (para uniones).\\
		\opt{-{}-split-\allowbreak{}pages[=\allowbreak{}N]} & Divide en archivos de una página (o de \val{N} páginas). \val{\%d} en el nombre de salida se sustituye por el número; sin él, el número se añade al final.\\
		\opt{-{}-rotate=\allowbreak{}[+|\allowbreak{}-]ÁNGULO[:RANGO]} & Gira 90, 180 o 270 grados; \val{+} y \val{-} son relativos a la rotación actual; sin rango, todas las páginas.\\
		\opt{-{}-collate[=N]} & Intercala las páginas de los archivos de \opt{-{}-pages} (1.ª de A, 1.ª de B, 2.ª de A\ldots).\\
		\opt{-{}-overlay F [opc] -{}-} & Superpone las páginas de \val{F} encima (sellos, marcas de agua). Opciones \opt{-{}-to}, \opt{-{}-from}, \opt{-{}-repeat}.\\
		\opt{-{}-underlay F [opc] -{}-} & Igual, pero debajo del contenido (membretes, fondos).\\*
		\opt{-{}-flatten-\allowbreak{}rotation} & Incorpora la rotación al contenido en lugar de dejarla como atributo.\\
	\end{opciones}
	
	\begin{lstlisting}
	# Unir varios PDF completos
	qpdf --empty --pages portada.pdf cap1.pdf cap2.pdf -- libro.pdf
	
	# Unir con selección: páginas 1-3 de a.pdf y la última de b.pdf
	qpdf --empty --pages a.pdf 1-3 b.pdf z -- seleccion.pdf
	
	# Extraer páginas del propio archivo (el punto es la entrada)
	qpdf doc.pdf --pages . 1,3,5-7 -- extracto.pdf
	
	# Quitar la página 2
	qpdf doc.pdf --pages . 1-z,x2 -- sin_pag2.pdf
	
	# Invertir el orden
	qpdf doc.pdf --pages . z-1 -- invertido.pdf
	
	# Una página por archivo: pag-1.pdf, pag-2.pdf, ...
	qpdf --split-pages doc.pdf pag-%d.pdf
	
	# Bloques de 10 páginas: bloque-01-10.pdf, bloque-11-20.pdf, ...
	qpdf --split-pages=10 doc.pdf bloque.pdf
	
	# Girar 90 grados la página 1 y 180 todas las demás
	qpdf doc.pdf girado.pdf --rotate=+90:1 --rotate=180:2-z
	
	# Escáner sin dúplex: intercalar caras anteriores y posteriores (invertidas)
	qpdf --collate --empty --pages anversos.pdf reversos.pdf z-1 -- completo.pdf
	
	# Marca de agua: superponer la página 1 de sello.pdf en todas las páginas
	qpdf doc.pdf --overlay sello.pdf --from=1 --repeat=1 -- doc_sellado.pdf
\end{lstlisting}
	
	\necesitaespacio{8\baselineskip}
	\subsection{Cifrado y descifrado}\label{sec:qpdf-cifrado}
	
	Un PDF puede tener dos contraseñas: la de \textbf{usuario} (necesaria para abrirlo) y la de \textbf{propietario} (permite cambiar los permisos: imprimir, copiar, modificar).
	
	\begin{lstlisting}[style=sintaxis]
  qpdf --encrypt --user-password=U --owner-password=P --bits=256 [restricciones] \
  -- <entrada> <salida>
\end{lstlisting}
	
	\necesitaespacio{20\baselineskip}
	\begin{opciones}
		\opt{-{}-user-\allowbreak{}password=\allowbreak{}U} & Contraseña de apertura (vacía: se abre sin contraseña pero conserva las restricciones).\\*
		\opt{-{}-owner-\allowbreak{}password=\allowbreak{}P} & Contraseña de propietario.\\
		\opt{-{}-bits=256} & AES de 256 bits (recomendado). Los cifrados de 40 y 128 bits son débiles.\\
		\opt{-{}-print=\allowbreak{}full|\allowbreak{}low|\allowbreak{}none} & Impresión permitida en alta calidad, baja o ninguna.\\
		\opt{-{}-modify=\allowbreak{}all|\allowbreak{}annotate|\allowbreak{}form|\allowbreak{}assembly|\allowbreak{}none} & Grado de modificación permitido.\\
		\opt{-{}-extract=y|n} & Permitir copiar texto e imágenes.\\
		\opt{-{}-decrypt} & Elimina el cifrado (requiere la contraseña con \opt{-{}-password}).\\*
		\opt{-{}-password=\allowbreak{}CLAVE} & Contraseña para abrir la entrada.\\
	\end{opciones}
	
	\begin{lstlisting}
	# Cifrar con contraseña de apertura
	qpdf --encrypt --user-password='lectura' --owner-password='propietario' \
	--bits=256 -- \
	doc.pdf doc_cifrado.pdf
	
	# Sin contraseña de apertura, pero sin copiar ni modificar
	qpdf --encrypt --user-password= --owner-password='propietario' --bits=256 \
	--print=low --modify=none --extract=n -- doc.pdf doc_restringido.pdf
	
	# Descifrar un PDF del que conoce la contraseña
	qpdf --password='propietario' --decrypt doc_cifrado.pdf doc_abierto.pdf
\end{lstlisting}
	
	\begin{importante}[Uso legítimo]
		Descifre solo documentos propios o para los que tenga autorización y la contraseña. Las contraseñas escritas en la orden quedan en el historial de Bash (\cmd{history}); para no dejarlas, escriba un espacio antes del comando (Ubuntu no guarda esas líneas con la configuración predeterminada de \texttt{HISTCONTROL}) o use un archivo: \opt{-{}-password-\allowbreak{}file=\allowbreak{}archivo}.
	\end{importante}
	
	\necesitaespacio{26\baselineskip}
	\subsection{Optimización, linealización y compresión sin pérdida}
	
	\necesitaespacio{22\baselineskip}
	\begin{opciones}
		\opt{-{}-linearize} & «Vista web rápida»: reorganiza el archivo para mostrar la primera página antes de descargarlo entero.\\*
		\opt{-{}-object-\allowbreak{}streams=\allowbreak{}generate} & Agrupa objetos en flujos comprimidos (reduce el tamaño). \val{preserve} (predeterminado) o \val{disable}.\\
		\opt{-{}-compress-\allowbreak{}streams=\allowbreak{}y} & Comprime los flujos sin comprimir.\\
		\opt{-{}-recompress-\allowbreak{}flate} & Vuelve a comprimir los flujos ya comprimidos con Flate.\\
		\opt{-{}-compression-\allowbreak{}level=\allowbreak{}N} & Nivel de Flate de 1 a 9 (9 = máxima compresión, más lento).\\
		\opt{-{}-optimize-\allowbreak{}images} & \textbf{Con pérdida}: recomprime como JPEG las imágenes que no lo están cuando el resultado es menor.\\
		\opt{-{}-remove-\allowbreak{}unreferenced-\allowbreak{}resources=\allowbreak{}yes} & Elimina recursos no usados por las páginas.\\*
		\opt{-{}-deterministic-\allowbreak{}id} & Identificador del archivo calculado a partir de su contenido (resultados reproducibles).\\
	\end{opciones}
	
	\begin{lstlisting}
  qpdf --linearize doc.pdf doc_web.pdf
  qpdf --object-streams=generate --compress-streams=y --recompress-flate \
  --compression-level=9 doc.pdf doc_compacto.pdf      # sin pérdida
  qpdf --optimize-images doc.pdf doc_jpeg.pdf              # con pérdida en imágenes
\end{lstlisting}
	
	\begin{nota}
		La compresión sin pérdida de QPDF suele lograr reducciones modestas; cuando el tamaño se debe a imágenes de alta resolución, la reducción importante viene de Ghostscript (sección~\ref{sec:gs-reduccion}), que sí reduce la resolución de las imágenes.
	\end{nota}
	
	\necesitaespacio{30\baselineskip}
	\subsection{Metadatos, adjuntos y otras transformaciones}
	
	\necesitaespacio{26\baselineskip}
	\begin{opciones}
		\opt{-{}-remove-info} & Elimina el diccionario de información (título, autor\ldots), salvo la fecha de modificación.\\*
		\opt{-{}-remove-\allowbreak{}metadata} & Elimina los metadatos XMP.\\
		\opt{-{}-remove-\allowbreak{}structure} & Elimina el árbol de estructura lógica (PDF etiquetado).\\
		\opt{-{}-remove-\allowbreak{}page-\allowbreak{}labels} & Elimina la numeración lógica de páginas (por ejemplo, i, ii, iii).\\
		\opt{-{}-set-\allowbreak{}page-\allowbreak{}labels E... -{}-} & Asigna la numeración lógica. Cada especificación es \val{página:tipo[/\allowbreak{}inicio[/\allowbreak{}prefijo]]}, con tipos \val{D} (arábigos), \val{r}/\val{R} (romanos) y \val{a}/\val{A} (letras). \val{1:r 3:D} numera i, ii y después 1, 2, 3\ldots\\
		\opt{-{}-add-\allowbreak{}attachment F -{}-} & Adjunta un archivo.\\
		\opt{-{}-remove-\allowbreak{}attachment=\allowbreak{}CLAVE} & Elimina un adjunto (clave mostrada por \opt{-{}-list-\allowbreak{}attachments}).\\
		\opt{-{}-flatten-\allowbreak{}annotations=\allowbreak{}all|\allowbreak{}print|\allowbreak{}screen} & Integra las anotaciones (sellos, comentarios, campos rellenados) en el contenido.\\*
		\opt{-{}-generate-\allowbreak{}appearances} & Genera la apariencia de los campos de formulario (útil antes de aplanar).\\
	\end{opciones}
	
	\begin{lstlisting}
	qpdf --remove-info --remove-metadata doc.pdf doc_anonimo.pdf
	qpdf doc.pdf etiquetado.pdf --set-page-labels 1:r 3:D --     # i, ii, 1, 2, 3...
	qpdf doc.pdf --add-attachment datos.csv -- doc_con_datos.pdf
	qpdf --generate-appearances --flatten-annotations=all formulario.pdf \
	formulario_plano.pdf
\end{lstlisting}
	
	\begin{nota}[Escribir título o autor]
		QPDF no tiene una opción directa para \emph{asignar} título o autor. Puede hacerse editando su representación JSON (\opt{-{}-json-output} y \opt{-{}-update-\allowbreak{}from-\allowbreak{}json}, uso avanzado), indicándolo al crear el PDF (\cmd{img2pdf -{}-title}, \cmd{ocrmypdf -{}-title}, metadatos de Pandoc) o con la herramienta externa \cmd{exiftool} (paquete \texttt{libimage-\allowbreak{}exiftool-\allowbreak{}perl}): \cmd{exiftool -Title=\allowbreak{}"..." doc.pdf}.
	\end{nota}
	
	\necesitaespacio{12\baselineskip}
	\subsection{Sobrescribir el archivo de entrada}
	
	\begin{advertencia}[\texttt{-{}-replace-\allowbreak{}input}]
		\cmd{qpdf doc.pdf -{}-replace-\allowbreak{}input -{}-linearize} sustituye el original por el resultado. QPDF escribe primero un archivo temporal y solo reemplaza si no hubo errores, pero no queda copia del original. Úselo solo sobre copias.
	\end{advertencia}
	
	\necesitaespacio{17\baselineskip}
	\subsection{Formatos de entrada y salida}
	
	\necesitaespacio{13\baselineskip}
	\begin{formatos}
		PDF & \compat & \compat & Entrada y salida principales (cifrados, dañados, linealizados\ldots).\\*
		JSON (QPDF) & \compat & \compat & Representación completa del PDF (\opt{-{}-json-output}, \opt{-{}-json-input}).\\
		QDF & \compat & \compat & PDF «legible» para estudio (\opt{-{}-qdf}).\\*
		Imágenes, texto & \nocompat & \nocompat & QPDF no renderiza ni extrae texto (use Poppler).\\
	\end{formatos}
	
	\necesitaespacio{8\baselineskip}
	\subsection{Limitaciones y seguridad}
	
	\begin{limitacion}
		\begin{itemize}
			\item No renderiza, no extrae texto, no hace OCR, no edita el texto de las páginas.
			\item La reducción de tamaño sin pérdida es limitada; \opt{-{}-optimize-\allowbreak{}images} es con pérdida.
			\item Las restricciones de propietario (imprimir, copiar) dependen de que el lector las respete: no son una protección fuerte. La contraseña de \emph{usuario} con AES-256 sí cifra el contenido.
			\item Al unir documentos, la estructura de nivel de documento (marcadores, formularios) de los archivos secundarios puede no conservarse completa: revise el resultado.
		\end{itemize}
	\end{limitacion}
	
	\necesitaespacio{30\baselineskip}
	\subsection{Errores frecuentes}
	
	\necesitaespacio{26\baselineskip}
	\begin{errores}
		\texttt{invalid password} (código 2) & PDF cifrado sin contraseña o con contraseña incorrecta. & \opt{-{}-password=\allowbreak{}'...'}; comillas simples si contiene \texttt{\$} o \texttt{!}.\\*
		Código de salida 3 & Avisos: el PDF tenía defectos que QPDF corrigió. & Revise el resultado; \opt{-{}-warning-\allowbreak{}exit-\allowbreak{}0} si el script debe continuar.\\
		\texttt{input file and output file are the same} & Misma ruta como entrada y salida. & Otro nombre, o \opt{-{}-replace-\allowbreak{}input} (sobre una copia).\\
		Las páginas no aparecen en \opt{-{}-pages} & Falta el \opt{-{}-} final. & Termine la lista con \opt{-{}-} antes del nombre de salida.\\
		Rango no válido & Página inexistente o sintaxis incorrecta. & \cmd{qpdf -{}-show-\allowbreak{}npages}; sintaxis en la sección~\ref{sec:qpdf-rangos}.\\
		\texttt{operation for ... attempted on object of type null} & Archivo dañado. & Reescriba con \cmd{qpdf} o con Ghostscript.\\*
		Contraseña visible en el historial & Se escribió en la orden. & Espacio inicial o \opt{-{}-password-\allowbreak{}file}.\\
	\end{errores}
	
	Ayuda: \cmd{qpdf -{}-help}, \cmd{qpdf -{}-help=all}, \cmd{qpdf -{}-help=-{}-opción}, \cmd{man qpdf} y \url{https://qpdf.readthedocs.io}. Versión: \cmd{qpdf -{}-version}.
	
	% =====================================================================
	\clearpage
	\section{OCRmyPDF}\label{sec:ocrmypdf}
	% =====================================================================
	
	\necesitaespacio{8\baselineskip}
	\subsection{Descripción general}
	
	\textbf{Qué es.} Un programa que añade a un PDF escaneado una \textbf{capa de texto invisible} obtenida por OCR (reconocimiento óptico de caracteres), alineada con las palabras de la imagen. El aspecto de las páginas no cambia, pero el texto pasa a poder seleccionarse, copiarse y buscarse.
	
	\textbf{Qué problema resuelve.} Convertir escaneos en PDF buscables y archivables (PDF/A), sin rehacer el documento.
	
	\textbf{Cuándo usarlo:} PDF de escáner o de aplicaciones de escaneo del móvil; documentos a los que hay que buscar texto o pasar después por Marker o \cmd{pdftotext}.
	
	\textbf{Cuándo no usarlo:} si el PDF ya es digital (no hace falta OCR); si lo que necesita es Markdown estructurado (Marker); si necesita reconocer escritura a mano (Tesseract está diseñado para texto impreso).
	
	\necesitaespacio{22\baselineskip}
	\subsubsection{Tres tipos de PDF}
	
	\begingroup\small
	\necesitaespacio{18\baselineskip}
	\begin{longtable}{L{3.3cm}L{5.96cm}L{5.8cm}}
		\toprule\rowcolor{tablacab}\cab{Tipo} & \cab{Qué contiene} & \cab{Cómo reconocerlo}\\\midrule
		\endfirsthead
		\bottomrule
		\endfoot
		\textbf{PDF digital} & Texto real con fuentes, generado por un programa (Word, LaTeX\ldots). Texto seleccionable y nítido a cualquier ampliación. & \cmd{pdffonts} lista fuentes normales; \cmd{pdftotext} devuelve el texto.\\*
		\textbf{PDF escaneado} & Solo imágenes de las páginas. No se puede seleccionar ni buscar texto. & \cmd{pdffonts} sin fuentes; \cmd{pdftotext} vacío; \cmd{pdfimages -list} muestra una imagen por página.\\*
		\textbf{PDF escaneado con OCR} & Imágenes de las páginas \emph{más} una capa de texto invisible reconocida. & \cmd{pdffonts} muestra una fuente como \texttt{GlyphLessFont}; \cmd{pdftotext} devuelve texto (con posibles errores de OCR).\\
	\end{longtable}
	\endgroup
	
	\necesitaespacio{8\baselineskip}
	\subsection{OCRmyPDF y Tesseract}
	
	\textbf{Tesseract} es el \emph{motor} de OCR: recibe la imagen de una página y devuelve los caracteres reconocidos con su posición. No sabe manejar PDF. \textbf{OCRmyPDF} es el \emph{orquestador}: abre el PDF, rasteriza cada página, la prepara, llama a Tesseract, construye la capa de texto invisible, la coloca bajo la imagen original y optimiza el resultado.
	
	\flujo{escaneado.pdf,OCRmyPDF: rasterizar (Ghostscript),imagen de cada página,preprocesado opcional (rotar{,} enderezar{,} limpiar),Tesseract (reconocimiento),texto + posiciones,OCRmyPDF: capa invisible + PDF/A + optimización,escaneado\_ocr.pdf}
	
	\begin{itemize}
		\item La calidad del reconocimiento depende de Tesseract, de los \textbf{datos de idioma} instalados y de la resolución de la imagen.
		\item OCRmyPDF conserva las imágenes originales salvo que pida operaciones que las cambian (\opt{-{}-deskew}, \opt{-{}-clean-final}, \opt{-{}-force-ocr} u optimizaciones con pérdida).
	\end{itemize}
	
	\necesitaespacio{8\baselineskip}
	\subsection{Instalación, idiomas y comprobación}
	
	\begin{lstlisting}
	sudo apt install ocrmypdf tesseract-ocr-spa unpaper pngquant jbig2
	ocrmypdf --version
	tesseract --version
	tesseract --list-langs              # idiomas instalados (eng, osd, spa...)
	apt-cache search tesseract-ocr- | less   # paquetes de idiomas disponibles
	# ejemplo: francés y portugués
	sudo apt install tesseract-ocr-fra tesseract-ocr-por
\end{lstlisting}
	
	Los datos de idioma se guardan en \arch{/\allowbreak{}usr/\allowbreak{}share/\allowbreak{}tesseract-\allowbreak{}ocr/\allowbreak{}5/\allowbreak{}tessdata/}. El paquete \texttt{osd} (detección de orientación y escritura) es el que usa \opt{-{}-rotate-pages}.
	
	\begingroup\small
	\necesitaespacio{16\baselineskip}
	\begin{longtable}{L{3.4cm}L{11.66cm}}
		\toprule\rowcolor{tablacab}\cab{Dependencia} & \cab{Función}\\\midrule
		\endfirsthead
		\bottomrule
		\endfoot
		Tesseract (\texttt{tesseract-ocr}) & Motor de OCR (obligatorio).\\*
		\texttt{tesseract-\allowbreak{}ocr-\allowbreak{}XXX} & Datos de cada idioma (\texttt{spa}, \texttt{eng}\ldots).\\
		Ghostscript & Rasterizado de páginas y conversión a PDF/A (obligatorio en esta versión).\\
		\texttt{unpaper} & Limpieza (\opt{-{}-clean}, \opt{-{}-clean-final}). Opcional.\\
		\texttt{pngquant} & Compresión con pérdida de imágenes PNG en \opt{-O 2} y \opt{-O 3}. Opcional.\\*
		\texttt{jbig2} (jbig2enc) & Compresión JBIG2 de páginas en blanco y negro. Opcional.\\
	\end{longtable}
	\endgroup
	
	\necesitaespacio{8\baselineskip}
	\subsection{Sintaxis}
	
	\begin{lstlisting}[style=sintaxis]
  ocrmypdf [opciones] <entrada.pdf | imagen | -> <salida.pdf | ->
\end{lstlisting}
	
	\begin{lstlisting}
	ocrmypdf entrada.pdf salida.pdf          # idioma predeterminado: inglés
	ocrmypdf -l spa entrada.pdf salida.pdf   # español
\end{lstlisting}
	
	El archivo de salida se sobrescribe si existe. Si la salida es el mismo archivo que la entrada, la entrada solo se actualiza si el proceso termina bien.
	
	\necesitaespacio{27\baselineskip}
	\subsection{Opciones}
	
	\necesitaespacio{24\baselineskip}
	\subsubsection{Idioma, salida y metadatos}
	
	\necesitaespacio{20\baselineskip}
	\begin{opciones}
		\opt{-l IDIOMAS} & Idiomas de Tesseract separados por \val{+}: \val{spa}, \val{spa+eng}. Indíquelos siempre: con el idioma correcto mejoran mucho las tildes y la eñe.\\*
		\opt{-{}-output-\allowbreak{}type T} & \val{pdfa} (predeterminado: PDF/A-2b, para archivo a largo plazo), \val{pdf} (cambia lo menos posible el original), \val{pdfa-1}, \val{pdfa-2}, \val{pdfa-3}, \val{none} (sin PDF, útil con \opt{-{}-sidecar}).\\
		\opt{-{}-sidecar [ARCHIVO]} & Guarda además el texto reconocido en un \arch{.txt} (\val{-} para la pantalla).\\
		\opt{-{}-image-dpi N} & Resolución que debe suponerse si la entrada es una imagen sin resolución válida.\\*
		\opt{-{}-title}, \opt{-{}-author}, \opt{-{}-subject}, \opt{-{}-keywords} & Metadatos del PDF de salida.\\
	\end{opciones}
	
	\necesitaespacio{29\baselineskip}
	\subsubsection{Preprocesado de imagen}
	
	\necesitaespacio{25\baselineskip}
	\begin{opciones}
		\opt{-r}, \opt{-{}-rotate-pages} & Detecta la orientación del texto y gira las páginas (90, 180, 270 grados).\\*
		\opt{-{}-rotate-\allowbreak{}pages-\allowbreak{}threshold N} & Confianza mínima para girar (súbala si gira páginas que no debía).\\
		\opt{-d}, \opt{-{}-deskew} & Endereza páginas ligeramente inclinadas (modifica la imagen).\\
		\opt{-c}, \opt{-{}-clean} & Limpia la imagen con \cmd{unpaper} \emph{solo para el OCR}; el PDF conserva la imagen original.\\
		\opt{-i}, \opt{-{}-clean-final} & Limpia y además usa la imagen limpia en el resultado (puede borrar contenido deseado).\\
		\opt{-{}-unpaper-\allowbreak{}args "ARGS"} & Argumentos adicionales para \cmd{unpaper} (requiere \opt{-{}-clean}).\\
		\opt{-{}-remove-\allowbreak{}background} & Intenta blanquear el fondo de páginas en gris o color.\\
		\opt{-{}-oversample N} & Remuestrea las imágenes a al menos \val{N} ppp antes del OCR (puede mejorar escaneos de baja resolución).\\*
		\opt{-{}-remove-\allowbreak{}vectors} & Experimental: ignora objetos vectoriales en el OCR (evita caracteres falsos).\\
	\end{opciones}
	
	\necesitaespacio{17\baselineskip}
	\subsubsection{Páginas que ya tienen texto}
	
	Por defecto, si una página ya contiene texto, OCRmyPDF \textbf{se detiene sin generar el archivo} (error \texttt{PriorOcrFoundError}, código 6), para no estropear PDF digitales o ya procesados. Hay tres alternativas, incompatibles entre sí:
	
	\necesitaespacio{13\baselineskip}
	\begin{opciones}
		\opt{-s}, \opt{-{}-skip-text} & Omite las páginas con texto y hace OCR del resto. Ideal para documentos mixtos (páginas digitales y escaneadas).\\*
		\opt{-{}-redo-ocr} & Detecta y elimina una capa OCR invisible anterior y la rehace. Conserva el texto visible y real.\\*
		\opt{-f}, \opt{-{}-force-ocr} & Rasteriza \emph{todo} (también el texto digital) y hace OCR. El resultado es un PDF de imágenes con capa OCR: se pierde la calidad vectorial del texto. Último recurso.\\
	\end{opciones}
	
	\necesitaespacio{26\baselineskip}
	\subsubsection{Optimización y calidad}
	
	\necesitaespacio{22\baselineskip}
	\begin{opciones}
		\opt{-O 0} & Sin optimización.\\*
		\opt{-O 1} & Predeterminado: optimizaciones sin pérdida (incluye JBIG2 sin pérdida si \texttt{jbig2} está instalado).\\
		\opt{-O 2}, \opt{-O 3} & Optimizaciones con pérdida de JPEG (y PNG con \cmd{pngquant}); \val{3} es más agresivo.\\
		\opt{-{}-jpeg-\allowbreak{}quality Q}, \opt{-{}-png-\allowbreak{}quality Q} & Calidad (1--100) usada por las optimizaciones con pérdida.\\
		\opt{-{}-jbig2-lossy} & JBIG2 con pérdida: archivos muy pequeños, pero puede \emph{confundir caracteres parecidos} (un 6 por un 8). No recomendable en documentos con cifras importantes.\\
		\opt{-{}-pdfa-\allowbreak{}image-\allowbreak{}compression M} & En PDF/A: \val{auto}, \val{jpeg} o \val{lossless}.\\*
		\opt{-{}-fast-\allowbreak{}web-\allowbreak{}view MB} & Linealiza los archivos mayores que el umbral indicado (\val{0}: siempre).\\
	\end{opciones}
	
	\necesitaespacio{34\baselineskip}
	\subsubsection{Rendimiento y control avanzado}
	
	\necesitaespacio{30\baselineskip}
	\begin{opciones}
		\opt{-j N}, \opt{-{}-jobs N} & Núcleos de CPU (predeterminado: todos). Redúzcalo si WSL2 se queda sin memoria.\\*
		\opt{-{}-pages RANGO} & Solo esas páginas (desde 1): \val{1-3,7}.\\
		\opt{-{}-skip-big MPX} & No hace OCR de páginas de más de \val{MPX} megapíxeles.\\
		\opt{-{}-tesseract-\allowbreak{}timeout S} & Abandona el OCR de una página tras \val{S} segundos.\\
		\opt{-{}-tesseract-\allowbreak{}pagesegmode N} & Modo de segmentación de Tesseract (\cmd{tesseract -{}-help-\allowbreak{}psm}): \val{6} para un bloque uniforme de texto, \val{4} para columnas de tamaño variable.\\
		\opt{-{}-tesseract-\allowbreak{}oem N} & Motor de Tesseract: \val{1} redes neuronales (LSTM), \val{3} predeterminado.\\
		\opt{-{}-tesseract-\allowbreak{}thresholding M} & Umbralizado de Tesseract (por ejemplo \val{adaptive-otsu} o \val{sauvola}) para escaneos de mala calidad.\\
		\opt{-{}-user-\allowbreak{}words ARCHIVO} & Lista de palabras especializadas que Tesseract debe considerar.\\
		\opt{-{}-invalidate-\allowbreak{}digital-\allowbreak{}signatures} & Permite procesar un PDF firmado (la firma dejará de ser válida).\\
		\opt{-{}-pdf-\allowbreak{}renderer R} & Cómo se escribe la capa de texto: \val{auto}, \val{hocr}, \val{sandwich}.\\*
		\opt{-k}, \opt{-q}, \opt{-v N} & Conservar temporales, silencioso, detallado (diagnóstico).\\
	\end{opciones}
	
	\necesitaespacio{8\baselineskip}
	\subsection{Ejemplos}
	
	\begin{lstlisting}
	cd /mnt/c/Users/NOMBRE/Downloads
	
	# OCR en español
	ocrmypdf -l spa "acta escaneada.pdf" "acta escaneada_ocr.pdf"
	
	# Documento bilingüe, páginas giradas o inclinadas, limpieza para el OCR
	ocrmypdf -l spa+eng --rotate-pages --deskew --clean escaneo.pdf escaneo_ocr.pdf
	
	# Guardar también el texto reconocido
	ocrmypdf -l spa --sidecar escaneo.txt escaneo.pdf escaneo_ocr.pdf
	
	# Solo el texto, sin generar PDF
	ocrmypdf -l spa --sidecar escaneo.txt --output-type none escaneo.pdf -
	
	# PDF mixto: páginas digitales y escaneadas
	ocrmypdf -l spa --skip-text mixto.pdf mixto_ocr.pdf
	
	# Rehacer un OCR antiguo de mala calidad
	ocrmypdf -l spa --redo-ocr ocr_antiguo.pdf ocr_nuevo.pdf
	
	# Cambiar lo mínimo el original (sin conversión a PDF/A)
	ocrmypdf -l spa --output-type pdf escaneo.pdf escaneo_ocr.pdf
	
	# Archivo más pequeño (con pérdida en imágenes)
	ocrmypdf -l spa -O 3 --jpeg-quality 70 escaneo.pdf escaneo_ligero.pdf
	
	# Desde imágenes: una imagen directamente, o varias a través de img2pdf
	ocrmypdf -l spa --image-dpi 300 pagina.png pagina_ocr.pdf
	img2pdf --pagesize A4 pag_*.jpg | ocrmypdf -l spa - documento_ocr.pdf
\end{lstlisting}
	
	\necesitaespacio{6\baselineskip}
	\subsubsection{Conservar la calidad de las imágenes}
	
	\begin{itemize}
		\item Use \opt{-{}-output-\allowbreak{}type pdf} y \opt{-O 1} (o \opt{-O 0}) para no recomprimir.
		\item Evite \opt{-{}-clean-final} y \opt{-{}-force-ocr} salvo necesidad.
		\item Para PDF/A sin pérdida: \opt{-{}-pdfa-\allowbreak{}image-\allowbreak{}compression lossless}.
		\item \opt{-{}-deskew} modifica la imagen: úselo solo si la inclinación perjudica el OCR o la lectura.
	\end{itemize}
	
	\necesitaespacio{8\baselineskip}
	\subsection{Procesamiento por lotes}
	
	\begin{lstlisting}
	mkdir -p ~/trabajo/ocr
	for f in ~/trabajo/escaneos/*.pdf; do
	out="$HOME/trabajo/ocr/$(basename "$f")"
	[ -e "$out" ] && continue                     # reanudar sin repetir
	ocrmypdf -l spa --rotate-pages --deskew --skip-text --jobs 4 "$f" "$out" \
	|| echo "Fallo ($?) en: $f" >> ~/trabajo/ocr/errores.log
	done
\end{lstlisting}
	
	OCRmyPDF ya usa varios núcleos por documento: procese los archivos de uno en uno y regule los núcleos con \opt{-{}-jobs}.
	
	\necesitaespacio{26\baselineskip}
	\subsection{Códigos de salida}
	
	\begingroup\small
	\necesitaespacio{22\baselineskip}
	\begin{longtable}{C{1.4cm}L{13.66cm}}
		\toprule\rowcolor{tablacab}\cab{Código} & \cab{Significado}\\\midrule
		\endfirsthead
		\toprule\rowcolor{tablacab}\cab{Código} & \cab{Significado}\\\midrule
		\endhead
		\bottomrule
		\endfoot
		0 & Correcto.\\*
		1 & Argumentos no válidos.\\
		2 & Archivo de entrada no válido o ilegible.\\
		3 & Falta una dependencia (por ejemplo, Tesseract o un idioma).\\
		4 & El PDF de salida no es válido.\\
		5 & Error de acceso a archivos (permisos, ruta).\\
		6 & La página ya tiene texto (use \opt{-{}-skip-text}, \opt{-{}-redo-ocr} o \opt{-{}-force-ocr}).\\
		7 & Error de un programa auxiliar (Tesseract, Ghostscript\ldots).\\
		8 & PDF cifrado: descífrelo antes con QPDF.\\
		9 & Configuración no válida.\\
		10 & Falló la conversión a PDF/A (el resultado puede ser un PDF normal).\\
		15 & Otro error.\\*
		130 & Interrumpido con \tecla{Ctrl}+\tecla{C}.\\
	\end{longtable}
	\endgroup
	
	\necesitaespacio{12\baselineskip}
	\subsection{OCRmyPDF y Marker}
	
	OCRmyPDF y Marker no hacen lo mismo: OCRmyPDF produce un \emph{PDF} con texto buscable; Marker produce \emph{Markdown} estructurado. Encadenarlos tiene sentido cuando el equipo es modesto y el documento es un escaneo de texto corrido:
	
	\flujo{PDF escaneado,OCRmyPDF,PDF con capa de texto,Marker (-\kern0pt-disable\_ocr),Markdown}
	
	Así, el reconocimiento lo hace Tesseract en CPU (relativamente ligero) y Marker solo estructura el texto, sin arrancar su modelo de visión. A cambio, se pierden las ecuaciones y la calidad depende del OCR de Tesseract. Con un equipo potente o documentos con fórmulas, es mejor dejar que Marker reconozca directamente el escaneo.
	
	\necesitaespacio{20\baselineskip}
	\subsection{Formatos de entrada y salida}
	
	\necesitaespacio{16\baselineskip}
	\begin{formatos}
		PDF & \compat & \compat & Entrada escaneada o mixta; salida PDF o PDF/A.\\*
		PDF/A (1b, 2b, 3b) & --- & \compat & Predeterminado: PDF/A-2b.\\
		Imagen (PNG, JPEG, TIFF\ldots) & \condic & \nocompat & Una sola imagen; para varias, img2pdf primero. Puede requerir \opt{-{}-image-dpi}.\\
		TXT (\opt{-{}-sidecar}) & \nocompat & \compat & Texto reconocido.\\*
		PDF cifrado & \nocompat & \nocompat & Descifrar antes (código 8).\\
	\end{formatos}
	
	\necesitaespacio{8\baselineskip}
	\subsection{Limitaciones}
	
	\begin{limitacion}
		\begin{itemize}
			\item \textbf{Errores de reconocimiento}: siempre hay errores; aumentan con baja resolución (menos de 200 ppp), papel sucio, fotocopias de fotocopias o tipografías inusuales. Lo recomendable es escanear a 300 ppp.
			\item \textbf{Tablas}: el texto queda en su sitio, pero la capa OCR no sabe qué es una tabla.
			\item \textbf{Fórmulas}: Tesseract no reconoce matemáticas.
			\item \textbf{Escritura a mano}: resultados pobres.
			\item \textbf{Inclinación y orientación}: requieren \opt{-{}-deskew} y \opt{-{}-rotate-pages}; inclinaciones grandes o páginas curvadas (libros fotografiados) se reconocen mal.
			\item \textbf{Idiomas}: sin el paquete del idioma correcto, el reconocimiento empeora; hay que indicarlo con \opt{-l}.
			\item \textbf{Diseños complejos}: columnas, recuadros y texto sobre imágenes pueden reconocerse en orden incorrecto.
			\item \textbf{PDF/A}: la conversión puede modificar el archivo (fuentes, colores, formularios, JavaScript).
			\item \textbf{Firmas digitales}: OCR invalida las firmas.
		\end{itemize}
	\end{limitacion}
	
	\necesitaespacio{14\baselineskip}
	\subsection{Errores frecuentes}
	
	\necesitaespacio{10\baselineskip}
	\begin{errores}
		\texttt{PriorOcrFoundError: page already has text!} (código 6) & El PDF ya tiene texto. & \opt{-{}-skip-text}, \opt{-{}-redo-ocr} o \opt{-{}-force-ocr}.\\*
		\texttt{OCR engine does not have language data for ...} & Idioma no instalado. & \cmd{sudo apt install tesseract-\allowbreak{}ocr-\allowbreak{}spa} (o el que falte).\\
		Tildes y eñes mal reconocidas & Idioma no indicado (inglés por defecto). & \opt{-l spa}.\\
		Código 8 & PDF cifrado. & \cmd{qpdf -{}-password=\allowbreak{}... -{}-decrypt}.\\
		\texttt{EncryptedPdfError} o firmas & PDF cifrado o firmado. & Descifrar o \opt{-{}-invalidate-\allowbreak{}digital-\allowbreak{}signatures}.\\
		\texttt{unpaper} no encontrado & Falta la dependencia de \opt{-{}-clean}. & \cmd{sudo apt install unpaper}.\\
		Muy lento o \texttt{Killed} & Muchas páginas en paralelo, memoria limitada de WSL2. & Baje \opt{-{}-jobs}; aumente memoria en \arch{.wslconfig}.\\
		Archivo resultante mayor que el original & Conversión a PDF/A o imágenes recomprimidas sin pérdida. & \opt{-{}-output-\allowbreak{}type pdf}, \opt{-O 2}/\opt{-O 3}, o Ghostscript después.\\
		Páginas giradas por error & Detección de orientación con poca confianza. & Suba \opt{-{}-rotate-\allowbreak{}pages-\allowbreak{}threshold} o quite \opt{-r}.\\*
		\texttt{ImportError} al arrancar & Paquetes de Python mezclados (\cmd{sudo pip}). & Ver sección~\ref{sec:pep668}; reinstale con \cmd{apt}.\\
	\end{errores}
	
	Ayuda: \cmd{ocrmypdf -{}-help}, \cmd{man ocrmypdf}, \cmd{tesseract -{}-help-extra} y \url{https://ocrmypdf.readthedocs.io} (documenta la última versión: compare con \cmd{ocrmypdf -{}-version}).
	% =====================================================================
	\clearpage
	\section{Ghostscript}\label{sec:ghostscript}
	% =====================================================================
	
	\necesitaespacio{8\baselineskip}
	\subsection{Descripción general}
	
	\textbf{Qué es.} Un intérprete de los lenguajes \textbf{PostScript} y \textbf{PDF}. PostScript es un lenguaje de programación que describe páginas (fue el estándar de las impresoras profesionales); PDF deriva de él, pero es un formato de documento, no un programa. Ghostscript lee el documento, lo \emph{interpreta} (calcula qué hay que dibujar en cada página) y envía el resultado a un \textbf{dispositivo de salida}: un nuevo PDF, un archivo PostScript, imágenes PNG, JPEG o TIFF, texto o el formato de una impresora.
	
	\textbf{Qué problema resuelve.} Convertir entre PostScript y PDF, rasterizar páginas, \emph{reescribir} un PDF completo (lo que permite reducir su tamaño, convertirlo a escala de grises, incrustar fuentes o ajustar la página) y producir formatos para imprenta e impresión.
	
	\textbf{Cuándo usarlo:} para reducir el tamaño de un PDF con imágenes pesadas; para convertir PS/EPS a PDF o viceversa; para rasterizar cuando se necesita un dispositivo concreto (TIFF de fax, PNG con transparencia); para normalizar PDF problemáticos.
	
	\textbf{Cuándo no usarlo:} para operaciones de páginas que deben conservar el archivo intacto (unir, dividir, rotar, cifrar: QPDF lo hace sin reinterpretar); para extraer texto con calidad (Poppler); para OCR (OCRmyPDF, que usa Ghostscript internamente).
	
	\begin{importante}[Ghostscript reescribe]
		Con el dispositivo \texttt{pdfwrite}, Ghostscript no «copia» el PDF: lo interpreta y genera uno nuevo. Las imágenes pueden recomprimirse, las fuentes se incrustan en subconjuntos y algunas estructuras (formularios, etiquetas de accesibilidad, contenido interactivo) pueden cambiar o perderse. Compare siempre el resultado con el original.
	\end{importante}
	
	\necesitaespacio{8\baselineskip}
	\subsection{Instalación y comprobación}
	
	\begin{lstlisting}
  sudo apt install ghostscript
  gs --version                                 # número de versión
  gs -h | sed -n '/Available devices/,/Search path/p'   # dispositivos disponibles
\end{lstlisting}
	
	El paquete instala además guiones que llaman a \cmd{gs} con opciones predefinidas: \cmd{ps2pdf}, \cmd{pdf2ps}, \cmd{eps2eps}, \cmd{ps2ps}, \cmd{ps2ascii}, \cmd{ps2epsi}\ldots
	
	\necesitaespacio{8\baselineskip}
	\subsection{Sintaxis general}
	
	\begin{lstlisting}[style=sintaxis]
  gs [opciones] -sDEVICE=<dispositivo> -o <salida> <entrada> [<entrada2> ...]
\end{lstlisting}
	
	\necesitaespacio{10\baselineskip}
	\begin{opciones}
		\opt{-sDEVICE=\allowbreak{}NOMBRE} & Dispositivo de salida (tabla siguiente). Las opciones \opt{-s} definen valores de texto.\\*
		\opt{-o ARCHIVO} & Archivo de salida. Equivale a \opt{-sOutputFile=\allowbreak{}ARCHIVO -dNOPAUSE -dBATCH}: es la forma recomendada y más corta.\\
		\opt{-sOutputFile=\allowbreak{}ARCHIVO} & Forma tradicional del archivo de salida; necesita además \opt{-dNOPAUSE} y \opt{-dBATCH}.\\
		\opt{-dNOPAUSE} & No se detiene tras cada página esperando una tecla.\\
		\opt{-dBATCH} & Sale al terminar en lugar de quedarse en el intérprete interactivo (\texttt{GS>}).\\
		\opt{-q}, \opt{-dQUIET} & Suprime los mensajes informativos.\\
		\opt{-r PPP} & Resolución en puntos por pulgada para dispositivos de imagen (\opt{-r300}, o \opt{-r300x150} por eje).\\
		\opt{-dFirstPage=N}, \opt{-dLastPage=N} & Primera y última página (desde 1).\\
		\opt{-sPageList=\allowbreak{}LISTA} & Páginas concretas: \val{1,3,5-7}; \val{5-} hasta el final; \val{odd}/\val{even}.\\
		\opt{-sPAPERSIZE=a4} & Tamaño de papel (\val{a4}, \val{letter}\ldots); con \opt{-dFIXEDMEDIA} se impone a todas las páginas.\\
		\opt{-dPDFFitPage} & Escala cada página para que quepa en el papel fijado.\\
		\opt{-dEPSCrop} & Al convertir EPS, recorta la página a la caja del dibujo.\\
		\opt{-dPDFSETTINGS=\allowbreak{}/\allowbreak{}P} & Preajuste de \texttt{pdfwrite} (sección~\ref{sec:gs-reduccion}).\\
		\opt{-dCompatibilityLevel=\allowbreak{}N} & Versión de PDF de salida (\val{1.4}, \val{1.7}\ldots).\\
		\opt{-dSAFER} & Modo seguro: restringe el acceso a archivos. \textbf{Activado por defecto}; \opt{-dNOSAFER} lo desactiva (no lo haga con archivos ajenos).\\*
		\opt{-dTextAlphaBits=\allowbreak{}4}, \opt{-dGraphicsAlphaBits=\allowbreak{}4} & Suavizado de texto y gráficos en dispositivos de imagen.\\
	\end{opciones}
	
	\paragraph{Una imagen por página} En dispositivos de imagen, el nombre de salida debe llevar un marcador de número de página: \arch{pag-\%03d.png} produce \arch{pag-001.png}, \arch{pag-002.png}\ldots{} Sin marcador, cada página sobrescribe a la anterior y solo queda la última.
	
	\necesitaespacio{34\baselineskip}
	\subsection{Dispositivos principales}
	
	\necesitaespacio{30\baselineskip}
	\begin{opciones}[Dispositivo]
		\val{pdfwrite} & Escribe un PDF nuevo (reescribir, comprimir, convertir PS $\rightarrow$ PDF).\\*
		\val{ps2write} & PostScript de nivel 2.\\
		\val{eps2write} & EPS (una página).\\
		\val{png16m} & PNG en color de 24 bits.\\
		\val{pngalpha} & PNG con transparencia.\\
		\val{pnggray}, \val{pngmono}, \val{png256} & PNG en grises, en blanco y negro o con paleta de 256 colores.\\
		\val{jpeg}, \val{jpeggray} & JPEG en color o en grises; calidad con \opt{-dJPEGQ=N} (0--100).\\
		\val{tiff24nc}, \val{tiffgray}, \val{tifflzw} & TIFF en color, en grises o comprimido con LZW.\\
		\val{tiffg4} & TIFF en blanco y negro con compresión de fax CCITT G4 (archivos muy pequeños para OCR o archivo).\\
		\val{pbmraw}, \val{pgmraw}, \val{ppmraw} & Formatos PNM (por ejemplo, para Potrace).\\
		\val{txtwrite} & Extrae el texto.\\
		\val{inkcov}, \val{ink\_cov} & Cobertura de tinta CMYK de cada página (estimar consumo de impresión).\\*
		\val{bbox} & Caja que encierra el dibujo de cada página.\\
	\end{opciones}
	
	\necesitaespacio{10\baselineskip}
	\subsection{Ejemplos}
	
	\necesitaespacio{6\baselineskip}
	\subsubsection{PDF a PDF}
	
	\begin{lstlisting}
	gs -sDEVICE=pdfwrite -o reescrito.pdf original.pdf                  # reescribir
	gs -sDEVICE=pdfwrite -dPDFSETTINGS=/ebook -o reducido.pdf original.pdf
	gs -sDEVICE=pdfwrite -dFirstPage=2 -dLastPage=5 -o pags_2a5.pdf original.pdf
	gs -sDEVICE=pdfwrite -sPageList=1,3,5-7 -o seleccion.pdf original.pdf
	gs -sDEVICE=pdfwrite -sColorConversionStrategy=Gray \
	-dProcessColorModel=/DeviceGray \
	-o gris.pdf original.pdf  # escala de grises
	gs -sDEVICE=pdfwrite -sPAPERSIZE=a4 -dFIXEDMEDIA -dPDFFitPage -o a4.pdf carta.pdf
	# texto -> contornos
	gs -sDEVICE=pdfwrite -dNoOutputFonts -o contornos.pdf original.pdf
	gs -sDEVICE=pdfwrite -o unido.pdf a.pdf b.pdf c.pdf  # unir (reinterpretando)
\end{lstlisting}
	
	\begin{nota}[Unir con Ghostscript o con QPDF]
		\cmd{gs} une PDF reinterpretándolos (útil si además quiere reducir tamaño). Si solo quiere unir sin alterar nada, \cmd{qpdf -{}-empty -{}-pages ...} es más fiel y más rápido.
	\end{nota}
	
	\necesitaespacio{6\baselineskip}
	\subsubsection{PDF a imágenes}
	
	\begin{lstlisting}
	gs -sDEVICE=png16m -r150 -dTextAlphaBits=4 -dGraphicsAlphaBits=4 -o \
	pag-%03d.png doc.pdf
	gs -sDEVICE=jpeg -dJPEGQ=90 -r150 -o pag-%03d.jpg doc.pdf
	gs -sDEVICE=pngalpha -r300 -dFirstPage=1 -dLastPage=1 -o portada.png doc.pdf
	gs -sDEVICE=tiffg4 -r300 -o fax.tif escaneado.pdf  # TIFF multipágina 1 bit
	gs -sDEVICE=pnggray -r300 -sPageList=4 -o pag4_gris.png doc.pdf
\end{lstlisting}
	
	\necesitaespacio{6\baselineskip}
	\subsubsection{PostScript y EPS}
	
	\begin{lstlisting}
	gs -sDEVICE=ps2write -o salida.ps doc.pdf          # PDF -> PostScript
	gs -sDEVICE=pdfwrite -o salida.pdf entrada.ps      # PostScript -> PDF
	gs -sDEVICE=pdfwrite -dEPSCrop -o figura.pdf figura.eps   # EPS -> PDF ajustado
	gs -sDEVICE=eps2write -dFirstPage=3 -dLastPage=3 -o pag3.eps doc.pdf
	ps2pdf entrada.ps salida.pdf                       # guiones equivalentes
	pdf2ps doc.pdf salida.ps
\end{lstlisting}
	
	\necesitaespacio{6\baselineskip}
	\subsubsection{Otros usos}
	
	\begin{lstlisting}
	# texto (Poppler suele dar mejor resultado)
	gs -sDEVICE=txtwrite -o texto.txt doc.pdf
	gs -q -o - -sDEVICE=inkcov doc.pdf            # cobertura C M Y K por página
	# cajas de dibujo (se muestran en pantalla)
	gs -q -dBATCH -dNOPAUSE -sDEVICE=bbox doc.pdf
\end{lstlisting}
	
	\necesitaespacio{13\baselineskip}
	\subsection{Reducción del tamaño de un PDF}\label{sec:gs-reduccion}
	
	\necesitaespacio{10\baselineskip}
	\subsubsection{Por qué ocupa tanto un PDF}
	
	El tamaño suele deberse a tres causas: \textbf{imágenes} de alta resolución o poco comprimidas (la más frecuente), \textbf{fuentes} incrustadas completas, y \textbf{estructura} poco eficiente (objetos duplicados, flujos sin comprimir). \texttt{pdfwrite} actúa sobre las tres: reduce la resolución de las imágenes que superan un umbral, recomprime, incrusta subconjuntos de fuentes y escribe una estructura nueva.
	
	\necesitaespacio{21\baselineskip}
	\subsubsection{Preajustes \texttt{-dPDFSETTINGS}}
	
	\begingroup\small
	\necesitaespacio{17\baselineskip}
	\begin{longtable}{L{2.0cm}L{2.6cm}L{1.5cm}L{2.96cm}L{1.0cm}L{3.1cm}}
		\toprule\rowcolor{tablacab}\cab{Preajuste} & \cab{Imágenes color/gris} & \cab{1 bit} & \cab{Color} & \cab{PDF} & \cab{Uso}\\\midrule
		\endfirsthead
		\toprule\rowcolor{tablacab}\cab{Preajuste} & \cab{Imágenes color/gris} & \cab{1 bit} & \cab{Color} & \cab{PDF} & \cab{Uso}\\\midrule
		\endhead
		\bottomrule
		\endfoot
		\val{/default} & Sin reducción & Sin reducción & Sin cambios & --- & Máxima fidelidad.\\*
		\val{/screen} & 72 ppp & 300 ppp & Convierte a sRGB & 1.5 & Solo pantalla; calidad baja.\\
		\val{/ebook} & 150 ppp & 300 ppp & Convierte a sRGB & 1.5 & Lectura en pantalla y correo.\\
		\val{/printer} & 300 ppp & 1200 ppp & Independiente del dispositivo & 1.7 & Impresión de oficina.\\*
		\val{/prepress} & 300 ppp & 1200 ppp & Sin cambios & 1.7 & Imprenta.\\
	\end{longtable}
	\endgroup
	
	Las imágenes solo se reducen si su resolución supera 1,5 veces la del preajuste (por ejemplo, con \val{/ebook}, las de más de 225 ppp). La documentación de Ghostscript advierte que \val{/prepress} no significa «máxima calidad»: la salida más fiel al original se obtiene sin \opt{-dPDFSETTINGS} (o con \val{/default}), y los preajustes alteran el documento de varias formas.
	
	\begin{lstlisting}
		# Probar varios preajustes y comparar tamaños
		for p in screen ebook printer prepress; do
		gs -q -sDEVICE=pdfwrite -dPDFSETTINGS=/$p -o "doc_$p.pdf" doc.pdf
		done
		ls -lh doc.pdf doc_*.pdf
\end{lstlisting}
	
	\necesitaespacio{6\baselineskip}
	\subsubsection{Control fino de la resolución}
	
	\begin{lstlisting}
		gs -sDEVICE=pdfwrite -dCompatibilityLevel=1.7 \
		-dDownsampleColorImages=true -dColorImageResolution=200 \
		-dColorImageDownsampleType=/Bicubic \
		-dDownsampleGrayImages=true  -dGrayImageResolution=200 \
		-dGrayImageDownsampleType=/Bicubic \
		-dDownsampleMonoImages=true  -dMonoImageResolution=400 \
		-o medida.pdf doc.pdf
\end{lstlisting}
	
	\necesitaespacio{24\baselineskip}
	\subsubsection{Efectos sobre el contenido}
	
	\begingroup\small
	\necesitaespacio{20\baselineskip}
	\begin{longtable}{L{2.8cm}L{12.26cm}}
		\toprule\rowcolor{tablacab}\cab{Elemento} & \cab{Efecto de la reducción}\\\midrule
		\endfirsthead
		\bottomrule
		\endfoot
		Texto & Sigue siendo vectorial y nítido; no se rasteriza (salvo \opt{-dNoOutputFonts}, que lo convierte en contornos).\\*
		Fuentes & Se incrustan en subconjuntos (solo los caracteres usados): reduce tamaño sin cambiar el aspecto.\\
		Imágenes & Pierden resolución y pueden recomprimirse en JPEG: es la principal fuente de ahorro y de pérdida de calidad.\\
		Gráficos vectoriales & Se conservan; los muy complejos apenas cambian de tamaño.\\
		Color & \val{/screen} y \val{/ebook} convierten a sRGB; puede alterar colores de imprenta.\\*
		Interactividad & Formularios, etiquetas de accesibilidad o JavaScript pueden cambiar o perderse; revise enlaces y marcadores.\\
	\end{longtable}
	\endgroup
	
	\begin{advertencia}[No hay garantía de reducción]
		Si el PDF ya estaba optimizado, contiene sobre todo vectores o usa compresiones muy eficientes (JBIG2 en escaneos), el resultado puede ser igual o incluso mayor. En una prueba con un PDF de 118\,KB generado por LaTeX, todos los preajustes produjeron unos 19\,KB (el ahorro vino de las fuentes, no de las imágenes) y \val{/screen} no fue el más pequeño. Compare siempre tamaños y aspecto, y conserve el original.
	\end{advertencia}
	
	Si el PDF es un escaneo, OCRmyPDF con \opt{-O 2} o \opt{-O 3} (y \texttt{jbig2} instalado) suele comprimir mejor que \texttt{pdfwrite}; si el problema es solo estructural, QPDF lo resuelve sin pérdida (sección~\ref{sec:flujo6}).
	
	\necesitaespacio{8\baselineskip}
	\subsection{PDF/A}
	
	\begin{lstlisting}
	gs -dPDFA=2 -dPDFACompatibilityPolicy=1 -sColorConversionStrategy=RGB \
	-sDEVICE=pdfwrite -o archivo_pdfa.pdf doc.pdf
\end{lstlisting}
	
	Genera un PDF orientado a PDF/A-2, pero la conformidad estricta puede requerir más ajustes (perfil de color de salida). Para archivo, OCRmyPDF genera PDF/A de forma más automática; en ambos casos conviene comprobar el resultado con un validador como veraPDF.
	
	\necesitaespacio{8\baselineskip}
	\subsection{Seguridad}\label{sec:gs-seguridad}
	
	PostScript es un lenguaje de programación completo, y los intérpretes de PDF y PostScript han tenido vulnerabilidades graves. Ghostscript activa por defecto el modo \opt{-dSAFER}, que impide que un documento lea o escriba archivos arbitrarios: en una prueba, un PostScript que intentaba leer \arch{/etc/passwd} se detuvo con el error \texttt{/\allowbreak{}invalidfileaccess}.
	
	\begin{itemize}
		\item No use \opt{-dNOSAFER} con documentos que no haya creado usted.
		\item Si un trabajo legítimo necesita leer un archivo concreto, autorícelo de forma puntual: \opt{-{}-permit-\allowbreak{}file-\allowbreak{}read=\allowbreak{}/\allowbreak{}ruta/}.
		\item Mantenga Ghostscript actualizado (\cmd{sudo apt update \&\& sudo apt upgrade}): las correcciones de seguridad llegan como actualizaciones.
		\item ImageMagick y OCRmyPDF llaman internamente a Ghostscript: procesar un PDF con ellos también ejecuta el intérprete.
	\end{itemize}
	
	\necesitaespacio{24\baselineskip}
	\subsection{Diferencias con QPDF, Poppler e ImageMagick}
	
	\begingroup\small
	\necesitaespacio{20\baselineskip}
	\begin{longtable}{L{2.5cm}L{3.66cm}L{4.1cm}L{4.0cm}}
		\toprule\rowcolor{tablacab}\cab{Herramienta} & \cab{Cómo trata el PDF} & \cab{Puntos fuertes} & \cab{No usar para}\\\midrule
		\endfirsthead
		\bottomrule
		\endfoot
		Ghostscript & Lo interpreta y escribe uno nuevo. & Reducir tamaño, PS$\leftrightarrow$PDF, rasterizar, normalizar. & Operaciones que deben conservar el archivo intacto.\\*
		QPDF & Copia y reorganiza sus objetos sin interpretarlos. & Páginas, cifrado, reparación estructural, linealización. & Reducir imágenes, rasterizar, extraer texto.\\
		Poppler & Lo lee y lo renderiza. & Texto, imágenes, rasterizado de calidad, diagnóstico. & Modificar o comprimir.\\*
		ImageMagick & Lo rasteriza a través de Ghostscript. & Procesar después las imágenes. & PDF largos o conservar el texto.\\
	\end{longtable}
	\endgroup
	
	\necesitaespacio{22\baselineskip}
	\subsection{Formatos de entrada y salida}
	
	\necesitaespacio{18\baselineskip}
	\begin{formatos}
		PDF & \compat & \compat & Entrada interpretada; salida con \texttt{pdfwrite}.\\*
		PostScript (PS) & \compat & \compat & Entrada nativa; salida con \texttt{ps2write}.\\
		EPS & \compat & \compat & Entrada con \opt{-dEPSCrop}; salida con \texttt{eps2write} (una página).\\
		PNG, JPEG, TIFF & \nocompat & \compat & Dispositivos de imagen (no lee imágenes como entrada).\\
		PBM, PGM, PPM & \nocompat & \compat & \texttt{pbmraw}, \texttt{pgmraw}, \texttt{ppmraw}.\\
		TXT & \nocompat & \compat & \texttt{txtwrite}.\\*
		PDF/A & --- & \condic & Con \opt{-dPDFA}; conformidad sujeta a validación.\\
	\end{formatos}
	
	\necesitaespacio{8\baselineskip}
	\subsection{Limitaciones}
	
	\begin{limitacion}
		\begin{itemize}
			\item Reescribe el documento: posibles cambios en imágenes, color, formularios, accesibilidad y metadatos.
			\item La reducción de tamaño no está garantizada.
			\item Las fuentes no incrustadas se sustituyen por otras disponibles (avisos de fuente no encontrada).
			\item Rasterizar a alta resolución consume mucha memoria y disco.
			\item Opciones numerosas y poco intuitivas; los preajustes cambian muchos parámetros a la vez.
		\end{itemize}
	\end{limitacion}
	
	\necesitaespacio{14\baselineskip}
	\subsection{Errores frecuentes}
	
	\necesitaespacio{10\baselineskip}
	\begin{errores}
		\texttt{Unknown device} & Nombre de dispositivo mal escrito o no compilado. & \cmd{gs -h} lista los disponibles.\\*
		\texttt{/\allowbreak{}undefinedfilename} & Archivo inexistente o ruta con espacios sin comillas. & Compruebe con \cmd{ls -l "ruta"}.\\
		\texttt{/\allowbreak{}invalidfileaccess} & Modo SAFER impidió un acceso a archivo. & Si es legítimo, \opt{-{}-permit-\allowbreak{}file-\allowbreak{}read=}; nunca \opt{-dNOSAFER} con archivos ajenos.\\
		Solo aparece la última página como imagen & Falta \val{\%d} en el nombre de salida. & Use \arch{pag-\%03d.png}.\\
		Se queda en el indicador \texttt{GS>} & Faltan \opt{-dBATCH} y \opt{-dNOPAUSE}. & Use \opt{-o}, o escriba \texttt{quit}.\\
		PDF mayor que el original & Contenido ya optimizado o recompresión desfavorable. & Otro preajuste, QPDF sin pérdida, o conserve el original.\\
		Imágenes borrosas tras comprimir & Preajuste demasiado agresivo (\val{/screen}). & \val{/ebook} o \val{/printer}, o resolución a medida.\\
		Avisos de fuentes sustituidas & Fuentes no incrustadas en el original. & Revise el aspecto; incruste fuentes al generar el PDF de origen.\\*
		\texttt{Unrecoverable error} & Error genérico (suele acompañar a otro más arriba). & Lea el primer mensaje de error; quite \opt{-q} para ver detalles.\\
	\end{errores}
	
	Ayuda: \cmd{gs -h}, \cmd{man gs} y la documentación en \url{https://ghostscript.readthedocs.io}. Versión: \cmd{gs -{}-version}.
	% =====================================================================
	\clearpage
	\section{Comparación entre herramientas}\label{sec:comparacion}
	% =====================================================================
	
	\necesitaespacio{14\baselineskip}
	\subsection{Visión de conjunto}
	
	\begingroup\footnotesize
	\setlength{\tabcolsep}{4pt}
	\necesitaespacio{10\baselineskip}
	\begin{longtable}{L{2.1cm}L{3.5cm}L{2.9cm}L{2.9cm}L{3.30cm}}
		\toprule\rowcolor{tablacab}\cab{Herramienta} & \cab{Función principal} & \cab{Entrada} & \cab{Salida} & \cab{Mejor uso}\\\midrule
		\endfirsthead
		\toprule\rowcolor{tablacab}\cab{Herramienta} & \cab{Función principal} & \cab{Entrada} & \cab{Salida} & \cab{Mejor uso}\\\midrule
		\endhead
		\bottomrule
		\endfoot
		\textbf{Marker} & Extraer y estructurar el contenido de documentos & PDF, imágenes; DOCX, PPTX, XLSX, HTML, EPUB con extras & Markdown, JSON, HTML, chunks, imágenes & PDF $\rightarrow$ Markdown con títulos, tablas y fórmulas\\*
		\textbf{Poppler} & Leer, extraer y renderizar PDF & PDF & TXT, PNG, JPEG, TIFF, SVG, PS, EPS, HTML, PDF & Diagnóstico, texto de PDF digitales, rasterizado\\
		\textbf{img2pdf} & Empaquetar imágenes en PDF sin pérdida & JPEG, PNG, TIFF, GIF\ldots & PDF & Escaneos y fotos $\rightarrow$ PDF\\
		\textbf{ImageMagick} & Procesar y convertir imágenes & Cientos de formatos raster; PDF/PS vía Ghostscript & Imágenes; PDF raster & Redimensionar, recortar, ajustar, lotes de imágenes\\
		\textbf{Potrace} & Vectorizar mapas de bits & PBM, PGM, PPM, BMP & SVG, PDF, EPS, PS, DXF\ldots & Logotipos y dibujos $\rightarrow$ vector\\
		\textbf{Pandoc} & Convertir entre formatos de documento & Markdown, DOCX, HTML, LaTeX, ODT, EPUB\ldots & Los anteriores, PDF (con motor), PPTX\ldots & Markdown $\rightarrow$ DOCX/PDF/HTML\\
		\textbf{QPDF} & Transformar la estructura del PDF & PDF, JSON de QPDF & PDF, JSON & Unir, dividir, rotar, cifrar, reparar sin alterar contenido\\
		\textbf{OCRmyPDF} & Añadir capa de texto OCR & PDF escaneado, una imagen & PDF/PDF/A buscable, TXT & Escaneos $\rightarrow$ PDF buscable\\*
		\textbf{Ghostscript} & Interpretar PostScript/PDF y reescribir & PDF, PS, EPS & PDF, PS, EPS, PNG, JPEG, TIFF, TXT & Reducir tamaño, PS$\leftrightarrow$PDF\\
	\end{longtable}
	\endgroup
	
	\necesitaespacio{14\baselineskip}
	\subsection{¿Qué herramienta para cada necesidad?}\label{sec:necesidades}
	
	\begingroup\small
	\necesitaespacio{10\baselineskip}
	\begin{longtable}{L{4.6cm}L{4.2cm}L{6.26cm}}
		\toprule\rowcolor{tablacab}\cab{Necesidad} & \cab{Herramienta recomendada} & \cab{Alternativa o comentario}\\\midrule
		\endfirsthead
		\toprule\rowcolor{tablacab}\cab{Necesidad} & \cab{Herramienta recomendada} & \cab{Alternativa o comentario}\\\midrule
		\endhead
		\bottomrule
		\endfoot
		PDF $\rightarrow$ texto & Poppler (\cmd{pdftotext}) & Si sale vacío, OCRmyPDF antes.\\*
		PDF $\rightarrow$ imágenes & Poppler (\cmd{pdftoppm}) & Ghostscript (dispositivos PNG/TIFF); ImageMagick para pocas páginas.\\
		Imágenes $\rightarrow$ PDF & img2pdf & ImageMagick si hay que transformar antes.\\
		PDF $\rightarrow$ Markdown & Marker & \cmd{pdftotext} si basta el texto sin estructura.\\
		Markdown $\rightarrow$ PDF & Pandoc + LaTeX & Pandoc con \cmd{xelatex} para Unicode completo.\\
		Imagen $\rightarrow$ SVG & ImageMagick + Potrace & Solo imágenes de dos colores.\\
		Redimensionar imágenes & ImageMagick & \cmd{magick mogrify -path} para lotes.\\
		Extraer imágenes de un PDF & Poppler (\cmd{pdfimages}) & Gráficos vectoriales: \cmd{pdftocairo -svg}.\\
		Unir PDF & QPDF & \cmd{pdfunite} para casos simples.\\
		Separar PDF & QPDF (\opt{-{}-split-pages}) & \cmd{pdfseparate}.\\
		Reparar o reestructurar un PDF & QPDF & Ghostscript si QPDF no basta.\\
		Inspeccionar la estructura & QPDF (\opt{-{}-check}, \opt{-{}-json}) & \cmd{pdfinfo}, \cmd{pdffonts}, \cmd{pdfimages -list}.\\
		OCR de PDF escaneado & OCRmyPDF & Marker si el objetivo es Markdown.\\
		Optimizar PDF (web, estructura) & QPDF (\opt{-{}-linearize}, flujos de objetos) & OCRmyPDF \opt{-{}-fast-\allowbreak{}web-\allowbreak{}view} en escaneos.\\
		Reducir tamaño de PDF & Ghostscript (\opt{-dPDFSETTINGS}) & OCRmyPDF \opt{-O 3} en escaneos.\\
		PDF $\rightarrow$ PostScript & Ghostscript (\texttt{ps2write}) & Poppler (\cmd{pdftops}).\\
		PostScript $\rightarrow$ PDF & Ghostscript (\cmd{ps2pdf}) & ---\\
		PDF $\rightarrow$ PDF transformado & QPDF (estructura) o Ghostscript (contenido) & Ver sección~\ref{sec:comparacionpdf}.\\
		Cifrar o descifrar con contraseña & QPDF & ---\\
		DOCX $\rightarrow$ Markdown o PDF & Pandoc & Fidelidad limitada en maquetas complejas.\\
		Extraer tablas de un PDF & Marker (\texttt{TableConverter}) & \cmd{pdftotext -layout} para tablas simples.\\
		PDF $\rightarrow$ HTML & Poppler (\cmd{pdftohtml}) & Marker (\opt{-{}-output\_\allowbreak{}format html}) para HTML estructurado.\\
		Marca de agua o membrete & QPDF (\opt{-{}-overlay}/\opt{-{}-underlay}) & ImageMagick para imágenes sueltas.\\*
		Procesar cientos de archivos & Bash (bucles, \cmd{find}) + la herramienta & Sección~\ref{sec:flujo10}.\\
	\end{longtable}
	\endgroup
	
	% =====================================================================
	\clearpage
	\section{Comparación específica de herramientas PDF}\label{sec:comparacionpdf}
	% =====================================================================
	
	Siete de las nueve herramientas trabajan con PDF, pero \textbf{no son intercambiables}: algunas operaciones aparecen en varias, aunque cada una las realiza de manera distinta y con consecuencias diferentes.
	
	\paragraph{Poppler} Orientado a la \emph{lectura}, extracción, renderizado y conversión de PDF. No modifica documentos (salvo separar y unir de forma simple).
	
	\paragraph{QPDF} Orientado a la \emph{transformación estructural}: reorganiza los objetos del PDF sin interpretarlos. Lo que entra es lo que sale: el contenido de las páginas no cambia.
	
	\paragraph{Ghostscript} Orientado a la \emph{interpretación}: entiende el documento como un programa de dibujo y genera uno nuevo. Por eso puede comprimir, convertir y rasterizar, y por eso también puede alterar el resultado.
	
	\paragraph{OCRmyPDF} Orientado a \emph{incorporar OCR} a PDF escaneados. Usa Ghostscript, Tesseract y otras bibliotecas por debajo.
	
	\paragraph{ImageMagick} Orientado a \emph{imágenes}. Ante un PDF, solo puede rasterizarlo (vía Ghostscript) o crear un PDF de imágenes.
	
	\paragraph{Marker} Orientado a \emph{extraer y estructurar} el contenido en Markdown, JSON o HTML. No produce PDF.
	
	\paragraph{img2pdf} Orientado a \emph{crear} PDF a partir de imágenes sin transformar sus datos.
	
	\necesitaespacio{14\baselineskip}
	\subsection{Matriz de operaciones}
	
	\compat{} = función principal \quad \condic{} = posible, con limitaciones o pérdidas \quad (vacío) = no aplicable
	
	\begingroup\footnotesize
	\setlength{\tabcolsep}{3pt}
	\necesitaespacio{10\baselineskip}
	\begin{longtable}{L{3.3cm}C{1.32cm}C{1.1cm}C{1.95cm}C{2.15cm}C{2.22cm}C{1.24cm}C{1.37cm}}
		\toprule\rowcolor{tablacab}\cabml{3.3cm}{Operación} & \cabmc{1.32cm}{Poppler} & \cabmc{1.1cm}{QPDF} & \cabmc{1.95cm}{Ghostscript} & \cabmc{2.15cm}{OCRmyPDF} & \cabmc{2.22cm}{ImageMagick} & \cabmc{1.24cm}{Marker} & \cabmc{1.37cm}{img2pdf}\\\midrule
		\endfirsthead
		\toprule\rowcolor{tablacab}\cabml{3.3cm}{Operación} & \cabmc{1.32cm}{Poppler} & \cabmc{1.1cm}{QPDF} & \cabmc{1.95cm}{Ghostscript} & \cabmc{2.15cm}{OCRmyPDF} & \cabmc{2.22cm}{ImageMagick} & \cabmc{1.24cm}{Marker} & \cabmc{1.37cm}{img2pdf}\\\midrule
		\endhead
		\bottomrule
		\endfoot
		Información y metadatos & \compat & \compat & & & \condic & & \\*
		Validar y reparar estructura & \condic & \compat & \condic & & & & \\
		Extraer texto & \compat & & \condic & \condic & & \compat & \\
		Extraer imágenes incrustadas & \compat & & & & & \condic & \\
		Rasterizar páginas & \compat & & \compat & & \condic & & \\
		Unir documentos & \compat & \compat & \condic & & \condic & & \\
		Separar páginas & \compat & \compat & \condic & & \condic & & \\
		Seleccionar y reordenar páginas & & \compat & \condic & & & & \\
		Rotar páginas & & \compat & & \condic & & & \\
		Cifrar / descifrar & & \compat & & & & & \\
		Linealizar (vista web rápida) & & \compat & & \condic & & & \\
		Reducir tamaño (imágenes) & & \condic & \compat & \condic & & & \\
		Compresión sin pérdida & & \compat & \condic & \condic & & & \\
		PostScript $\leftrightarrow$ PDF & \condic & & \compat & & & & \\
		OCR & & & & \compat & & \condic & \\
		PDF $\rightarrow$ Markdown & & & & & & \compat & \\*
		Imágenes $\rightarrow$ PDF & & & & \condic & \condic & & \compat\\
	\end{longtable}
	\endgroup
	
	Notas a la matriz: Poppler solo convierte PDF a PostScript (no al revés) y «repara» únicamente reescribiendo con \cmd{pdftocairo -pdf}; Ghostscript une y separa reinterpretando el contenido; OCRmyPDF solo rota según la orientación detectada del texto y solo acepta una imagen; ImageMagick une, separa y crea PDF rasterizando; Marker hace OCR, pero produce Markdown, no PDF.
	
	\begin{ejemplo}[La misma operación, tres resultados]
		Unir \arch{a.pdf} y \arch{b.pdf}: con \cmd{qpdf -{}-empty -{}-pages a.pdf b.pdf -{}- u.pdf} el resultado contiene exactamente los mismos objetos (texto, fuentes, imágenes); con \cmd{gs -sDEVICE=\allowbreak{}pdfwrite -o u.pdf a.pdf b.pdf}, un documento reescrito, quizá más pequeño y con imágenes recomprimidas; con \cmd{magick a.pdf b.pdf u.pdf}, un PDF de imágenes sin texto seleccionable. Las tres órdenes «unen PDF», pero solo la primera conserva el documento intacto.
	\end{ejemplo}
	
	% =====================================================================
	\clearpage
	\section{Flujos de trabajo combinando varias herramientas}\label{sec:flujos}
	% =====================================================================
	
	Cada flujo indica el problema, la cadena de herramientas, los comandos y por qué se eligió cada herramienta. Los comandos suponen que está en la carpeta de trabajo y que las rutas no contienen espacios; si los tienen, use comillas.
	
	\necesitaespacio{8\baselineskip}
	\subsection{Flujo 1: PDF digital a Markdown}\label{sec:flujo1}
	
	\flujoh{PDF digital,Marker,Markdown + imágenes}
	
	\begin{lstlisting}
  pdffonts articulo.pdf | head        # 1) confirmar que tiene texto (hay fuentes)
  source ~/marker-env/bin/activate
  # 2)
  marker_single articulo.pdf --output_dir md --mode fast --ocr_inline_math true
  ls md/articulo/  # 3) articulo.md, articulo_meta.json, imágenes
\end{lstlisting}
	
	\textbf{Por qué así:} Poppler confirma en un segundo que el PDF es digital; Marker usa entonces el texto incrustado (rápido) y solo recurre al modelo de visión para fórmulas. Sin fórmulas, \opt{-{}-disable\_ocr} es aún más rápido.
	
	\necesitaespacio{8\baselineskip}
	\subsection{Flujo 2: PDF escaneado, OCR y Markdown}\label{sec:flujo2}
	
	\flujo{PDF escaneado,OCRmyPDF,PDF con capa de texto,Marker (-\kern0pt-disable\_ocr),Markdown}
	
	\begin{lstlisting}
	ocrmypdf -l spa --rotate-pages --deskew escaneo.pdf escaneo_ocr.pdf
	pdftotext -l 1 escaneo_ocr.pdf - | head          # comprobar que ya hay texto
	source ~/marker-env/bin/activate
	marker_single escaneo_ocr.pdf --output_dir md --disable_ocr
\end{lstlisting}
	
	\textbf{Por qué así:} en un equipo sin GPU, Tesseract (dentro de OCRmyPDF) reconoce el texto con menos recursos que el modelo de visión de Marker, y Marker se limita a estructurarlo. \textbf{Cuándo no:} si el escaneo contiene fórmulas o tablas complejas y el equipo lo permite, es mejor \cmd{marker\_\allowbreak{}single escaneo.pdf -{}-force\_\allowbreak{}ocr}, que reconoce directamente con mayor calidad.
	
	\necesitaespacio{8\baselineskip}
	\subsection{Flujo 3: PDF a imágenes, procesamiento y de nuevo a PDF}\label{sec:flujo3}
	
	\flujo{doc.pdf,Poppler (pdftoppm),PNG por página,ImageMagick (mogrify),PNG procesados,img2pdf,doc\_procesado.pdf}
	
	\begin{lstlisting}
	mkdir -p pags proc
	pdftoppm -png -r 300 doc.pdf pags/p                        # 1) rasterizar
	# 2) blanquear fondo
	magick mogrify -path proc -colorspace Gray -level 15%,90% pags/*.png
	# 3) reempaquetar a tamaño real
	img2pdf --imgsize 300dpi proc/*.png -o doc_procesado.pdf
\end{lstlisting}
	
	\textbf{Por qué así:} Poppler rasteriza mejor y con menos memoria que ImageMagick; ImageMagick hace el retoque; img2pdf empaqueta sin recomprimir. \opt{-{}-imgsize 300dpi} asegura que las páginas recuperen su tamaño físico. \textbf{Importante:} el resultado es un PDF de imágenes; si el original tenía texto seleccionable, se pierde (aplique OCR después si lo necesita).
	
	\necesitaespacio{8\baselineskip}
	\subsection{Flujo 4: imagen a vector}\label{sec:flujo4}
	
	\flujo{logo.png,ImageMagick,logo.pbm (blanco y negro),Potrace,logo.svg / logo.pdf}
	
	\begin{lstlisting}
	magick logo.png -colorspace Gray -threshold 50% logo.pbm   # 1) blanco y negro
	potrace logo.pbm -s -t 4 -o logo.svg                       # 2) vector SVG
	potrace logo.pbm -b pdf -o logo.pdf                        # o PDF para LaTeX
\end{lstlisting}
	
	\textbf{Por qué así:} Potrace solo admite PBM/PGM/PPM/BMP y traza una frontera entre dos colores; ImageMagick controla con precisión esa frontera (el umbral).
	
	\necesitaespacio{8\baselineskip}
	\subsection{Flujo 5: Markdown a PDF}\label{sec:flujo5}
	
	\flujo{informe.md,Pandoc,informe.tex,latexmk / xelatex,informe.pdf}
	
	\begin{lstlisting}
	# Vía directa (Pandoc llama al motor y borra los intermedios)
	pandoc informe.md -o informe.pdf --toc -N -M lang=es -V geometry:margin=2.5cm
	
	# Vía controlada: generar LaTeX, revisarlo y compilarlo
	pandoc informe.md -s -o informe.tex -M lang=es -V geometry:margin=2.5cm
	latexmk -pdf informe.tex        # o latexmk -xelatex si usa fuentes del sistema
\end{lstlisting}
	
	\textbf{Por qué así:} la vía directa es cómoda; la controlada permite corregir el \arch{.tex} cuando el PDF no queda como se espera y ver los mensajes completos de LaTeX.
	
	\necesitaespacio{8\baselineskip}
	\subsection{Flujo 6: optimización de un PDF}\label{sec:flujo6}
	
	Antes de optimizar, averigüe \emph{por qué} pesa el PDF:
	
	\begin{lstlisting}
	pdfinfo doc.pdf | grep -E 'Pages|File size|Optimized'
	# ¿muchas imágenes, de alta resolución (x-ppi)?
	pdfimages -list doc.pdf | head -20
	pdffonts doc.pdf  # ¿fuentes completas sin subconjunto (sub = no)?
\end{lstlisting}
	
	\begingroup\small
	\necesitaespacio{13\baselineskip}
	\begin{longtable}{L{4.4cm}L{2.6cm}L{8.06cm}}
		\toprule\rowcolor{tablacab}\cab{Diagnóstico} & \cab{Herramienta} & \cab{Comando}\\\midrule
		\endfirsthead
		\bottomrule
		\endfoot
		Estructura poco eficiente, sin linealizar & QPDF & \cmd{qpdf -{}-object-\allowbreak{}streams=\allowbreak{}generate -{}-recompress-\allowbreak{}flate -{}-compression-\allowbreak{}level=\allowbreak{}9 -{}-linearize doc.pdf o.pdf}\\*
		Imágenes de alta resolución & Ghostscript & \cmd{gs -sDEVICE=\allowbreak{}pdfwrite -dPDFSETTINGS=\allowbreak{}/\allowbreak{}ebook -o o.pdf doc.pdf}\\*
		Escaneo (páginas = imágenes) & OCRmyPDF & \cmd{ocrmypdf -l spa -O 3 doc.pdf o.pdf} (o \opt{-{}-skip-text} si ya tiene OCR)\\
	\end{longtable}
	\endgroup
	
	\begin{lstlisting}
	# Probar las tres estrategias y quedarse con la más pequeña (revisando la calidad)
	qpdf --object-streams=generate --recompress-flate --compression-level=9 doc.pdf \
	op_qpdf.pdf
	gs -q -sDEVICE=pdfwrite -dPDFSETTINGS=/ebook -o op_gs.pdf doc.pdf
	qpdf --linearize op_gs.pdf op_gs_web.pdf       # Ghostscript + linealización
	ls -lS doc.pdf op_*.pdf                         # ordenados por tamaño
\end{lstlisting}
	
	\textbf{Por qué así:} QPDF no pierde calidad, así que se prueba primero; Ghostscript reduce imágenes (con pérdida); OCRmyPDF optimiza escaneos con técnicas específicas (JBIG2, cuantización). Combinar Ghostscript y después QPDF permite reducir y linealizar.
	
	\necesitaespacio{8\baselineskip}
	\subsection{Flujo 7: inspección y reparación}\label{sec:flujo7}
	
	\begin{lstlisting}
	# 1) Diagnóstico
	pdfinfo doc.pdf; pdffonts doc.pdf
	qpdf --check doc.pdf; echo "codigo: $?"         # 0 bien, 2 errores, 3 avisos
	# 2) Reparación estructural (sin tocar el contenido)
	qpdf doc.pdf reparado.pdf
	# 3) Si persiste: reinterpretar el contenido
	gs -sDEVICE=pdfwrite -o reparado_gs.pdf doc.pdf
	# 4) Último recurso: reescribir con Cairo
	pdftocairo -pdf doc.pdf reparado_cairo.pdf
	# 5) Verificar el resultado
	qpdf --check reparado.pdf && pdftotext -l 1 reparado.pdf - | head
	pdftoppm -png -r 72 -f 1 -l 1 reparado.pdf control && explorer.exe control-1.png
\end{lstlisting}
	
	\textbf{Por qué así:} Poppler informa sin modificar nada; QPDF repara estructura sin riesgo para el contenido; Ghostscript y Cairo reinterpretan (más agresivos, más capaces con daños de contenido). Se avanza de lo menos a lo más invasivo.
	
	\necesitaespacio{8\baselineskip}
	\subsection{Flujo 8: OCR y optimización de escaneos}\label{sec:flujo8}
	
	\flujo{fotos o escaneos JPG/PNG,img2pdf,escaneo.pdf,OCRmyPDF (OCR + limpieza + -O 2),escaneo\_ocr.pdf (buscable y ligero)}
	
	\begin{lstlisting}
	img2pdf --pagesize A4 --rotation=ifvalid escaneos/*.jpg -o escaneo.pdf
	ocrmypdf -l spa --rotate-pages --deskew --clean -O 2 \
	--sidecar escaneo.txt escaneo.pdf escaneo_ocr.pdf
	pdffonts escaneo_ocr.pdf | grep -i glyph        # capa OCR presente
	ls -lh escaneo.pdf escaneo_ocr.pdf
\end{lstlisting}
	
	\textbf{Por qué así:} img2pdf crea el PDF sin degradar las fotos; OCRmyPDF endereza, reconoce y comprime (JBIG2 sin pérdida para páginas en blanco y negro si \texttt{jbig2} está instalado). Para archivo a largo plazo, deje la salida PDF/A predeterminada; para cambiar lo mínimo, \opt{-{}-output-\allowbreak{}type pdf}.
	
	\necesitaespacio{8\baselineskip}
	\subsection{Flujo 9: PDF a imágenes y OCR}\label{sec:flujo9}
	
	\textbf{Cuándo NO conviene rasterizar antes:} casi nunca hace falta, porque OCRmyPDF ya rasteriza internamente cada página. Y nunca rasterice un PDF digital para «hacerle OCR»: perdería el texto real a cambio de un texto reconocido con errores.
	
	\textbf{Cuándo SÍ tiene sentido:} cuando necesita un preprocesado que OCRmyPDF no ofrece (eliminar un fondo de color, recortar márgenes con anotaciones, aumentar mucho el contraste de una fotocopia gastada).
	
	\flujo{escaneado.pdf,pdftoppm (gris{,} 300 ppp),PNG,ImageMagick (limpieza),PNG limpios,img2pdf + OCRmyPDF,limpio\_ocr.pdf}
	
	\begin{lstlisting}
	mkdir -p r l
	pdftoppm -gray -r 300 -png escaneado.pdf r/p
	magick mogrify -path l -level 25%,85% -fuzz 10% -trim +repage r/*.png
	img2pdf --imgsize 300dpi l/*.png | ocrmypdf -l spa - limpio_ocr.pdf
\end{lstlisting}
	
	\necesitaespacio{8\baselineskip}
	\subsection{Flujo 10: procesamiento masivo}\label{sec:flujo10}
	
	Este script recorre una carpeta de PDF, detecta si cada uno es escaneado (sin fuentes), aplica OCR solo cuando hace falta, extrae el texto y, opcionalmente, convierte a Markdown con Marker. Registra los errores y puede reanudarse.
	
	\begin{lstlisting}
	#!/usr/bin/env bash
	# masivo.sh - OCR selectivo + texto + Markdown para todos los PDF de ENTRADA
	set -u
	shopt -s nullglob
	ENTRADA="$HOME/trabajo/entrada"
	SALIDA="$HOME/trabajo/salida"
	# ejecutable del entorno, sin activarlo
	MARKER="$HOME/marker-env/bin/marker_single"
	mkdir -p "$SALIDA"/{pdf,txt,md}
	LOG="$SALIDA/registro.log"
	
	for f in "$ENTRADA"/*.pdf; do
	b="$(basename "$f" .pdf)"
	pdf="$SALIDA/pdf/$b.pdf"
	[ -e "$SALIDA/txt/$b.txt" ] && { echo "Omitido (ya hecho): $b"; continue; }
	
	nfuentes=$(pdffonts "$f" 2>/dev/null | tail -n +3 | wc -l)
	if [ "$nfuentes" -eq 0 ]; then
	echo "OCR: $b"
	ocrmypdf -q -l spa --rotate-pages --deskew "$f" "$pdf" \
	|| { echo "ERROR ocrmypdf ($?): $b" >> "$LOG"; continue; }
	else
	cp "$f" "$pdf"                             # digital: se copia tal cual
	fi
	
	pdftotext -layout "$pdf" "$SALIDA/txt/$b.txt" || echo "ERROR pdftotext: $b" \
	>> "$LOG"
	"$MARKER" "$pdf" --output_dir "$SALIDA/md" --disable_ocr \
	|| echo "ERROR marker: $b" >> "$LOG"
	done
	echo "Terminado. Errores registrados en: $LOG"
\end{lstlisting}
	
	\begin{lstlisting}
	chmod +x masivo.sh
	# ver y guardar la salida a la vez
	./masivo.sh 2>&1 | tee ~/trabajo/salida/consola.log
\end{lstlisting}
	
	\textbf{Por qué así:} \cmd{pdffonts} decide en milisegundos si hace falta OCR, lo que evita procesar de nuevo PDF digitales; las comprobaciones \texttt{[ -e ... ]} permiten reanudar; el ejecutable de Marker se llama con su ruta dentro del entorno, sin necesidad de activarlo en el script. Si no necesita Markdown, elimine las dos líneas de Marker; si los documentos tienen fórmulas, quite \opt{-{}-disable\_ocr}.
	% =====================================================================
	\clearpage
	\section{Tabla de formatos}\label{sec:formatos}
	% =====================================================================
	
	Relación entre formatos y herramientas según la documentación y las pruebas de cada programa. \textbf{L} = lectura (entrada); \textbf{E} = escritura (salida); \textbf{L/E} = ambas; \condic{} = con condiciones (ver notas); --- = no compatible o no aplicable.
	
	\begingroup\footnotesize
	\setlength{\tabcolsep}{2pt}
	\necesitaespacio{10\baselineskip}
	\begin{longtable}{L{2.15cm}C{1.15cm}C{1.22cm}C{1.27cm}C{2.04cm}C{1.19cm}C{1.15cm}C{1.01cm}C{1.98cm}C{1.8cm}}
		\toprule\rowcolor{tablacab}\cabml{2.15cm}{\scriptsize Formato} & \cabmc{1.15cm}{\scriptsize Marker} & \cabmc{1.22cm}{\scriptsize Poppler} & \cabmc{1.27cm}{\scriptsize img2pdf} & \cabmc{2.04cm}{\scriptsize ImageMagick} & \cabmc{1.19cm}{\scriptsize Potrace} & \cabmc{1.15cm}{\scriptsize Pandoc} & \cabmc{1.01cm}{\scriptsize QPDF} & \cabmc{1.98cm}{\scriptsize OCRmyPDF} & \cabmc{1.8cm}{\scriptsize Ghostscript}\\\midrule
		\endfirsthead
		\toprule\rowcolor{tablacab}\cabml{2.15cm}{\scriptsize Formato} & \cabmc{1.15cm}{\scriptsize Marker} & \cabmc{1.22cm}{\scriptsize Poppler} & \cabmc{1.27cm}{\scriptsize img2pdf} & \cabmc{2.04cm}{\scriptsize ImageMagick} & \cabmc{1.19cm}{\scriptsize Potrace} & \cabmc{1.15cm}{\scriptsize Pandoc} & \cabmc{1.01cm}{\scriptsize QPDF} & \cabmc{1.98cm}{\scriptsize OCRmyPDF} & \cabmc{1.8cm}{\scriptsize Ghostscript}\\\midrule
		\endhead
		\bottomrule
		\endfoot
		PDF & L & L/E\textsuperscript{1} & E & \condic L/E\textsuperscript{2} & E & \condic E\textsuperscript{3} & L/E & L/E & L/E\\*
		PostScript, EPS & --- & E & --- & \condic L/E\textsuperscript{2} & E & --- & --- & --- & L/E\\
		PNG & L & E & L & L/E & --- & --- & --- & \condic L\textsuperscript{4} & E\\
		JPEG & L & E & L & L/E & --- & --- & --- & \condic L\textsuperscript{4} & E\\
		TIFF & --- & E & L & L/E & --- & --- & --- & \condic L\textsuperscript{4} & E\\
		GIF, WebP & --- & --- & \condic L & L/E & --- & --- & --- & --- & ---\\
		BMP & --- & --- & \condic L & L/E & L & --- & --- & --- & ---\\
		PBM, PGM, PPM & --- & E & L & L/E & L (E: PGM) & --- & --- & \condic L\textsuperscript{4} & E\\
		SVG & --- & E & --- & \condic L/E\textsuperscript{5} & E & --- & --- & --- & ---\\
		DXF, GeoJSON & --- & --- & --- & --- & E & --- & --- & --- & ---\\
		Markdown & E & --- & --- & --- & --- & L/E & --- & --- & ---\\
		DOCX & \condic L\textsuperscript{6} & --- & --- & --- & --- & L/E & --- & --- & ---\\
		ODT & --- & --- & --- & --- & --- & L/E & --- & --- & ---\\
		PPTX, XLSX & \condic L\textsuperscript{6} & --- & --- & --- & --- & E (PPTX) & --- & --- & ---\\
		HTML & \condic L\textsuperscript{6}/E & E & --- & --- & --- & L/E & --- & --- & ---\\
		EPUB & \condic L\textsuperscript{6} & --- & --- & --- & --- & L/E & --- & --- & ---\\
		LaTeX & --- & --- & --- & --- & --- & L/E & --- & --- & ---\\
		TXT & --- & E & --- & --- & --- & E & --- & E & E\\*
		JSON & E & --- & --- & --- & --- & L/E\textsuperscript{7} & L/E\textsuperscript{8} & --- & ---\\
	\end{longtable}
	\endgroup
	
	{\footnotesize
		\textsuperscript{1} Escritura de PDF solo con \cmd{pdfseparate}, \cmd{pdfunite}, \cmd{pdfattach} y \cmd{pdftocairo -pdf}.
		\textsuperscript{2} A través de Ghostscript, rasterizando (sin texto seleccionable).
		\textsuperscript{3} Mediante un motor externo (LaTeX, Typst, motores HTML).
		\textsuperscript{4} Una sola imagen por ejecución; para varias, img2pdf primero.
		\textsuperscript{5} Lectura rasterizando; la escritura no vectoriza (para eso, Potrace).
		\textsuperscript{6} Requiere \texttt{marker-\allowbreak{}pdf[full]}.
		\textsuperscript{7} Árbol interno de Pandoc.
		\textsuperscript{8} Representación JSON propia de QPDF.\par}
	
	% =====================================================================
	\clearpage
	\section{Tabla de dependencias}\label{sec:dependencias}
	% =====================================================================
	
	Dependencias según los paquetes de Ubuntu 26.04 y la documentación de cada proyecto. \cmd{apt} instala automáticamente las dependencias obligatorias y, por defecto, las recomendadas.
	
	\begingroup\small
	\necesitaespacio{10\baselineskip}
	\begin{longtable}{L{2.2cm}L{3.96cm}L{4.3cm}L{4.0cm}}
		\toprule\rowcolor{tablacab}\cab{Herramienta} & \cab{Dependencias importantes} & \cab{Dependencias opcionales} & \cab{Para qué sirven}\\\midrule
		\endfirsthead
		\toprule\rowcolor{tablacab}\cab{Herramienta} & \cab{Dependencias importantes} & \cab{Dependencias opcionales} & \cab{Para qué sirven}\\\midrule
		\endhead
		\bottomrule
		\endfoot
		Marker & Python $\geq$ 3.10 (entorno con 3.12), PyTorch $\geq$ 2.7, \texttt{surya-ocr}, \texttt{pdftext}, \texttt{transformers}, Pillow; \cmd{llama-server} (CPU) o Docker + NVIDIA Container Toolkit (GPU) & \texttt{marker-\allowbreak{}pdf[full]} (mammoth, python-pptx, openpyxl, ebooklib, weasyprint); servicios LLM; \texttt{streamlit}; \texttt{fastapi} y \texttt{uvicorn} & PyTorch: modelos pequeños; surya: OCR, maquetación y fórmulas; extras: formatos distintos de PDF, interfaz y API\\*
		Poppler & \texttt{libpoppler}, Cairo, FreeType, LittleCMS & Fuentes del sistema & Renderizado, fuentes y gestión del color\\
		img2pdf & Python 3, \texttt{python3-\allowbreak{}img2pdf}, Pillow, pikepdf & --- & Leer imágenes y escribir el PDF\\
		ImageMagick & Bibliotecas MagickCore/MagickWand 7 & Ghostscript, netpbm, \texttt{libmagickcore-\allowbreak{}7.q16-\allowbreak{}10-\allowbreak{}extra}, fuentes & Ghostscript: leer PDF/PS/EPS; extra: formatos adicionales\\
		Potrace & \texttt{libpotrace} & ImageMagick (preparar la entrada) & Vectorización; \cmd{mkbitmap} se incluye\\
		Pandoc & \texttt{pandoc-data} & TeX Live (\cmd{pdflatex}, \cmd{xelatex}, \cmd{lualatex}), \cmd{latexmk}, Typst, WeasyPrint & Motores para generar PDF\\
		QPDF & \texttt{libqpdf} & --- & Lectura y escritura de la estructura del PDF\\
		OCRmyPDF & Tesseract, Ghostscript, pikepdf, img2pdf, pdfminer.six, Pillow, reportlab & \texttt{unpaper}, \texttt{pngquant}, \texttt{jbig2} (recomendados) y datos de idiomas & Tesseract: OCR; Ghostscript: rasterizado y PDF/A; opcionales: limpieza y compresión\\
		Tesseract & \texttt{libtesseract}, Leptonica, datos \texttt{eng} y \texttt{osd} & \texttt{tesseract-\allowbreak{}ocr-\allowbreak{}XXX} (spa, fra\ldots) & Motor OCR, idiomas y detección de orientación\\*
		Ghostscript & \texttt{libgs} & Fuentes del sistema & Intérprete PostScript/PDF y dispositivos\\
	\end{longtable}
	\endgroup
	
	% =====================================================================
	\clearpage
	\section{Tabla de comandos esenciales}\label{sec:comandos}
	% =====================================================================
	
	\begingroup\small
	\necesitaespacio{10\baselineskip}
	\begin{longtable}{L{5.0cm}L{10.06cm}}
		\toprule\rowcolor{tablacab}\cab{Tarea} & \cab{Comando}\\\midrule
		\endfirsthead
		\toprule\rowcolor{tablacab}\cab{Tarea} & \cab{Comando}\\\midrule
		\endhead
		\bottomrule
		\endfoot
		\multicolumn{2}{l}{\textbf{\color{acento}Inspección}}\\*
		Información de un PDF & \cmd{pdfinfo doc.pdf}\\
		¿Tiene texto o es escaneado? & \cmd{pdffonts doc.pdf}\\
		Número de páginas & \cmd{qpdf -{}-show-\allowbreak{}npages doc.pdf}\\
		Validar estructura & \cmd{qpdf -{}-check doc.pdf}\\
		Inventario de imágenes & \cmd{pdfimages -list doc.pdf}\\
		Información de una imagen & \cmd{magick identify foto.jpg}\\
		\multicolumn{2}{l}{\textbf{\color{acento}Extracción y conversión de PDF}}\\
		PDF $\rightarrow$ texto & \cmd{pdftotext -layout doc.pdf doc.txt}\\
		PDF $\rightarrow$ PNG (todas) & \cmd{pdftoppm -png -r 150 doc.pdf pag}\\
		PDF $\rightarrow$ JPEG (una página) & \cmd{pdftoppm -jpeg -r 150 -f 3 -l 3 doc.pdf p3}\\
		PDF $\rightarrow$ SVG (una página) & \cmd{pdftocairo -svg -f 1 -l 1 doc.pdf p1.svg}\\
		Extraer imágenes incrustadas & \cmd{pdfimages -all doc.pdf img}\\
		PDF $\rightarrow$ HTML & \cmd{pdftohtml -s -noframes doc.pdf doc.html}\\
		PDF $\rightarrow$ PostScript & \cmd{gs -sDEVICE=\allowbreak{}ps2write -o doc.ps doc.pdf}\\
		PostScript $\rightarrow$ PDF & \cmd{ps2pdf doc.ps doc.pdf}\\
		PDF $\rightarrow$ Markdown & \cmd{marker\_\allowbreak{}single doc.pdf -{}-output\_\allowbreak{}dir md}\\
		\multicolumn{2}{l}{\textbf{\color{acento}Páginas y estructura}}\\
		Unir PDF & \cmd{qpdf -{}-empty -{}-pages a.pdf b.pdf -{}- unido.pdf}\\
		Extraer páginas & \cmd{qpdf doc.pdf -{}-pages . 1-\allowbreak{}3,\allowbreak{}7 -{}- extracto.pdf}\\
		Una página por archivo & \cmd{qpdf -{}-split-\allowbreak{}pages doc.pdf pag-\%d.pdf}\\
		Rotar la página 2 & \cmd{qpdf doc.pdf girado.pdf -{}-rotate=\allowbreak{}+90:2}\\
		Reparar & \cmd{qpdf danado.pdf reparado.pdf}\\
		Cifrar (AES-256) & \cmd{qpdf -{}-encrypt -{}-user-\allowbreak{}password=\allowbreak{}U -{}-owner-\allowbreak{}password=\allowbreak{}P -{}-bits=\allowbreak{}256 -{}- e.pdf s.pdf}\\
		Descifrar & \cmd{qpdf -{}-password=\allowbreak{}P -{}-decrypt e.pdf s.pdf}\\
		Linealizar para la web & \cmd{qpdf -{}-linearize doc.pdf web.pdf}\\
		\multicolumn{2}{l}{\textbf{\color{acento}Tamaño y OCR}}\\
		Reducir tamaño & \cmd{gs -sDEVICE=\allowbreak{}pdfwrite -dPDFSETTINGS=\allowbreak{}/\allowbreak{}ebook -o s.pdf doc.pdf}\\
		Comprimir sin pérdida & \cmd{qpdf -{}-object-\allowbreak{}streams=\allowbreak{}generate -{}-recompress-\allowbreak{}flate doc.pdf s.pdf}\\
		OCR en español & \cmd{ocrmypdf -l spa escaneo.pdf escaneo\_\allowbreak{}ocr.pdf}\\
		OCR de documento mixto & \cmd{ocrmypdf -l spa -{}-skip-\allowbreak{}text mixto.pdf s.pdf}\\
		OCR con enderezado y rotación & \cmd{ocrmypdf -l spa -{}-rotate-\allowbreak{}pages -{}-deskew e.pdf s.pdf}\\
		\multicolumn{2}{l}{\textbf{\color{acento}Imágenes}}\\
		Imágenes $\rightarrow$ PDF & \cmd{img2pdf *.jpg -o album.pdf}\\
		Imágenes $\rightarrow$ PDF A4 & \cmd{img2pdf -{}-pagesize A4 *.jpg -o album.pdf}\\
		Convertir formato & \cmd{magick imagen.png imagen.jpg}\\
		Redimensionar al 50\,\% & \cmd{magick imagen.png -resize 50\% salida.png}\\
		Escala de grises & \cmd{magick foto.jpg -colorspace Gray gris.jpg}\\
		Lote a otra carpeta & \cmd{magick mogrify -path web -resize "1600x1600>" *.jpg}\\
		Imagen $\rightarrow$ SVG & \cmd{magick logo.png -threshold 50\% logo.pbm \&\& potrace logo.pbm -s -o logo.svg}\\
		\multicolumn{2}{l}{\textbf{\color{acento}Documentos}}\\
		Markdown $\rightarrow$ DOCX & \cmd{pandoc doc.md -o doc.docx}\\
		Markdown $\rightarrow$ PDF & \cmd{pandoc doc.md -o doc.pdf -M lang=\allowbreak{}es}\\
		Markdown $\rightarrow$ LaTeX & \cmd{pandoc doc.md -s -o doc.tex}\\
		DOCX $\rightarrow$ Markdown & \cmd{pandoc doc.docx -t gfm -{}-extract-\allowbreak{}media=\allowbreak{}img -o doc.md}\\*
		HTML $\rightarrow$ Markdown & \cmd{pandoc pagina.html -t gfm -o pagina.md}\\
	\end{longtable}
	\endgroup
	
	% =====================================================================
	\clearpage
	\section{Errores frecuentes del entorno y solución}\label{sec:errores}
	% =====================================================================
	
	Los errores propios de cada herramienta están al final de su capítulo. Esta tabla recoge los que aparecen con cualquiera de ellas en Ubuntu sobre WSL2.
	
	\necesitaespacio{10\baselineskip}
	\begin{errores}
		\texttt{command not found} & Programa no instalado, entorno virtual no activo o error de escritura. & \cmd{command -v programa}; instale con \cmd{apt}; active \arch{\tl/marker-env} para Marker.\\*
		\texttt{No such file or directory} & Ruta incorrecta, espacios sin comillas o mayúsculas distintas. & Comillas; \cmd{ls -l "ruta"}; tecla \tecla{Tab}.\\
		\texttt{Permission denied} al escribir & Carpeta sin permisos o archivo abierto en un programa de Windows. & Cierre el programa; escriba en \arch{\tl/trabajo}.\\
		\texttt{Permission denied} al ejecutar un script & El archivo no tiene permiso de ejecución. & \cmd{chmod +x script.sh} o \cmd{bash script.sh}.\\
		\texttt{\$'\bs{}\allowbreak{}r': command not found} & Script guardado con finales de línea de Windows (CRLF). & Guarde con LF o \cmd{sed -i 's/\allowbreak{}\bs{}\allowbreak{}r\$/\allowbreak{}/\allowbreak{}' script.sh}.\\
		\texttt{error: externally-\allowbreak{}managed-\allowbreak{}environment} & \cmd{pip} fuera de un entorno virtual. & Entorno virtual o \cmd{pipx} (sección~\ref{sec:pep668}).\\
		\texttt{ImportError} en una herramienta de \cmd{apt} & Paquetes de Python instalados con \cmd{sudo pip} que se superponen. & Desinstale lo instalado con \cmd{sudo pip}; reinstale con \cmd{apt}.\\
		\texttt{E: Unable to locate package} & Índice de paquetes desactualizado o nombre incorrecto. & \cmd{sudo apt update}; \cmd{apt search nombre}.\\
		\texttt{Could not get lock /\allowbreak{}var/\allowbreak{}lib/\allowbreak{}dpkg/\allowbreak{}lock-\allowbreak{}frontend} & Otra instalación en curso. & Espere a que termine; no borre el bloqueo.\\
		\texttt{Argument list too long} & Comodín con demasiados archivos. & Bucles, \cmd{find} o \cmd{xargs}.\\
		\texttt{Killed} & Memoria insuficiente en WSL2. & Menos procesos en paralelo; más memoria en \arch{.wslconfig}.\\
		\texttt{No space left on device} & Disco de Linux o temporal lleno. & \cmd{df -h}; borre intermedios o cachés (\cmd{du -sh \tl/\allowbreak{}.cache/\allowbreak{}*}).\\
		Procesos muy lentos sobre \arch{/mnt/c} & Acceso entre sistemas de archivos. & Copie a \arch{\tl/trabajo} (sección~\ref{sec:rendimiento}).\\
		Archivos \arch{*:Zone.Identifier} & Copia hecha con el Explorador de Windows. & \cmd{find . -name '*:Zone.Identifier' -delete} tras revisarlos.\\
		Orden incorrecto en lotes & Orden alfabético del comodín. & Ceros a la izquierda o \cmd{sort -V}.\\*
		Opción «desconocida» copiada de Internet & Opción de otra versión. & Compare con \cmd{-{}-version} y \cmd{-{}-help}.\\
	\end{errores}
	
	\necesitaespacio{22\baselineskip}
	\subsection{Obtener ayuda y versión desde la terminal}
	
	\begingroup\small
	\necesitaespacio{18\baselineskip}
	\begin{longtable}{L{2.6cm}L{6.06cm}L{6.4cm}}
		\toprule\rowcolor{tablacab}\cab{Herramienta} & \cab{Ayuda} & \cab{Versión}\\\midrule
		\endfirsthead
		\bottomrule
		\endfoot
		Marker & \cmd{marker\_\allowbreak{}single -{}-help}; \cmd{marker\_\allowbreak{}single config -{}-help} & \cmd{pip show marker-\allowbreak{}pdf}\\*
		Poppler & \cmd{pdftoppm -h}; \cmd{man pdftoppm} & \cmd{pdftoppm -v}\\
		img2pdf & \cmd{img2pdf -{}-help} & \cmd{img2pdf -{}-version}\\
		ImageMagick & \cmd{magick -help}; \cmd{man magick} & \cmd{magick -version}\\
		Potrace & \cmd{potrace -{}-help}; \cmd{man potrace} & \cmd{potrace -{}-version}\\
		Pandoc & \cmd{pandoc -{}-help}; \cmd{man pandoc} & \cmd{pandoc -{}-version}\\
		QPDF & \cmd{qpdf -{}-help}; \cmd{qpdf -{}-help=\allowbreak{}all} & \cmd{qpdf -{}-version}\\
		OCRmyPDF & \cmd{ocrmypdf -{}-help}; \cmd{man ocrmypdf} & \cmd{ocrmypdf -{}-version}\\
		Tesseract & \cmd{tesseract -{}-help-\allowbreak{}extra} & \cmd{tesseract -{}-version}\\*
		Ghostscript & \cmd{gs -h}; \cmd{man gs} & \cmd{gs -{}-version}\\
	\end{longtable}
	\endgroup
	
	% =====================================================================
	\clearpage
	\section{Seguridad}\label{sec:seguridad}
	% =====================================================================
	
	\necesitaespacio{12\baselineskip}
	\subsection{Riesgos de procesar documentos no confiables}
	
	Los PDF y, sobre todo, los PostScript pueden contener más que texto e imágenes: PostScript es un lenguaje de programación, y un PDF puede incluir JavaScript, archivos adjuntos y estructuras deliberadamente malformadas. Los programas que los interpretan (Ghostscript, Poppler, ImageMagick y, a través de ellos, OCRmyPDF) han tenido vulnerabilidades que permitían, con un archivo preparado, ejecutar código o leer archivos del sistema. Las imágenes también pueden estar diseñadas para agotar memoria («bombas de descompresión»).
	
	\necesitaespacio{8\baselineskip}
	\subsection{Protecciones que ya están activas}
	
	\begin{itemize}
		\item \textbf{Ghostscript} funciona por defecto en modo \opt{-dSAFER}, que bloquea el acceso a archivos desde el documento (sección~\ref{sec:gs-seguridad}).
		\item \textbf{ImageMagick} aplica \arch{/\allowbreak{}etc/\allowbreak{}ImageMagick-\allowbreak{}7/\allowbreak{}policy.xml}: límites de recursos, prohibición de leer listas \arch{@archivo} y de descargar URL (sección~\ref{sec:im-policy}).
		\item \textbf{img2pdf} (Pillow) rechaza imágenes con un número de píxeles sospechoso, salvo \opt{-{}-pillow-\allowbreak{}limit-\allowbreak{}break}.
	\end{itemize}
	
	\necesitaespacio{8\baselineskip}
	\subsection{Buenas prácticas de seguridad}
	
	\begin{enumerate}
		\item \textbf{Trabaje sobre copias}, nunca sobre el único ejemplar de un documento.
		\item \textbf{No ejecute archivos desconocidos} ni scripts descargados sin leerlos antes.
		\item \textbf{Mantenga el sistema actualizado}: \cmd{sudo apt update \&\& sudo apt upgrade}; dentro del entorno de Marker, \cmd{pip install -{}-upgrade marker-\allowbreak{}pdf}.
		\item \textbf{No desactive las protecciones} (\opt{-dNOSAFER}, relajar \arch{policy.xml}, \opt{-{}-pillow-\allowbreak{}limit-\allowbreak{}break}) para procesar archivos de terceros.
		\item \textbf{No use \cmd{sudo}} para procesar documentos: las herramientas no lo necesitan, y un fallo tendría privilegios de administrador.
		\item \textbf{Revise rutas y comodines} antes de operaciones masivas (\cmd{ls} o \cmd{echo} primero). Recuerde que desde WSL2 se puede borrar o sobrescribir también lo que hay en \arch{/mnt/c}.
		\item \textbf{Cuidado con las opciones que sobrescriben}: \opt{-{}-replace-\allowbreak{}input} (QPDF), \cmd{magick mogrify} sin \opt{-path}, \cmd{pdfunite} sin nombre de salida, y cualquier \opt{-o} que apunte a un archivo existente (OCRmyPDF, Ghostscript e img2pdf sobrescriben sin preguntar).
		\item \textbf{Contraseñas}: evite escribirlas en la orden (quedan en el historial); use \opt{-{}-password-\allowbreak{}file} en QPDF o un espacio inicial.
		\item \textbf{Datos confidenciales y servicios externos}: el modo \opt{-{}-use\_llm} de Marker envía contenido a terceros (salvo con Ollama local).
		\item \textbf{Metadatos}: antes de compartir, considere eliminar metadatos con \cmd{qpdf -{}-remove-\allowbreak{}info -{}-remove-\allowbreak{}metadata} o \cmd{magick ... -strip} en imágenes.
	\end{enumerate}
	
	% =====================================================================
	\clearpage
	\section{Buenas prácticas}\label{sec:buenas}
	% =====================================================================
	
	\begin{enumerate}
		\item \textbf{Organice el trabajo} en \arch{\tl/trabajo} con carpetas separadas de entrada, intermedios y salida. Es más rápido que \arch{/mnt/c} y evita mezclar originales y resultados.
		\item \textbf{Inspeccione antes de procesar}: \cmd{pdfinfo}, \cmd{pdffonts} y \cmd{pdfimages -list} dicen en segundos qué herramienta conviene.
		\item \textbf{Pruebe con pocas páginas} (\opt{-f}/\opt{-l}, \opt{-{}-pages}, \opt{-{}-page\_range}) antes de lanzar un documento o un lote completo.
		\item \textbf{Elija la herramienta menos invasiva}: QPDF antes que Ghostscript para operaciones de páginas; img2pdf antes que ImageMagick para crear PDF; \opt{-{}-skip-text} antes que \opt{-{}-force-ocr}.
		\item \textbf{No sobrescriba originales}: nombres de salida nuevos y descriptivos (\arch{informe\_\allowbreak{}ocr.pdf}, \arch{informe\_\allowbreak{}ebook.pdf}).
		\item \textbf{Verifique los resultados}: \cmd{qpdf -{}-check}, \cmd{pdftotext -l 1 ... -}, una página rasterizada para comprobar el aspecto y una comparación de tamaños con \cmd{ls -lh}.
		\item \textbf{Nombres ordenables}: use ceros a la izquierda en las secuencias (\arch{pag-001.png}) y evite espacios y caracteres especiales en los archivos que genere.
		\item \textbf{Comillas siempre} en variables y rutas: \texttt{"\$f"}.
		\item \textbf{Ensayo en seco} de los bucles con \cmd{echo} antes de ejecutarlos de verdad.
		\item \textbf{Registre los errores} de los lotes en un archivo y permita reanudarlos (comprobando si la salida ya existe).
		\item \textbf{Convierta lo repetitivo en scripts} y guárdelos con el resto del proyecto; anote las versiones usadas (\cmd{-{}-version}) para poder reproducir resultados.
		\item \textbf{Conserve la calidad en los pasos intermedios}: PNG o TIFF mientras procesa, JPEG solo en la entrega final.
	\end{enumerate}
	
	% =====================================================================
	\appendix
	\clearpage
	\section{Referencia rápida}\label{sec:refrapida}
	% =====================================================================
	
	\begingroup
	\lstset{basicstyle=\ttfamily\footnotesize, aboveskip=4pt, belowskip=6pt}
	
	\necesitaespacio{8\baselineskip}
	\subsection*{Entorno}
	\begin{lstlisting}
	cd /mnt/c/Users/NOMBRE/Downloads        # Descargas de Windows
	explorer.exe .                          # abrir la carpeta actual en Windows
	wslpath 'C:\Users\NOMBRE\Documents'     # traducir una ruta de Windows
	source ~/marker-env/bin/activate        # activar el entorno de Marker
\end{lstlisting}
	
	\necesitaespacio{8\baselineskip}
	\subsection*{Marker}
	\begin{lstlisting}
	marker_single doc.pdf --output_dir md                       # PDF -> Markdown
	# solo 2 primeras páginas (desde 0)
	marker_single doc.pdf --output_dir md --page_range "0-1"
	marker_single doc.pdf --output_dir md --mode fast --ocr_inline_math true
	marker_single doc.pdf --output_dir md --disable_ocr         # sin modelo de visión
	marker_single doc.pdf --output_dir md --force_ocr           # reconocer todo
	marker_single doc.pdf --output_dir js --output_format json
	marker carpeta_pdf --output_dir md --skip_existing          # lote
\end{lstlisting}
	
	\necesitaespacio{8\baselineskip}
	\subsection*{Poppler}
	\begin{lstlisting}
	pdfinfo archivo.pdf
	pdffonts archivo.pdf
	pdftotext archivo.pdf salida.txt
	pdftotext -layout -f 1 -l 3 archivo.pdf -
	pdftoppm -png archivo.pdf pagina
	pdftoppm -jpeg -jpegopt quality=85 -r 150 archivo.pdf pagina
	pdftoppm -png -r 300 -f 2 -l 2 -singlefile archivo.pdf pag2
	pdftocairo -svg -f 1 -l 1 archivo.pdf pag1.svg
	pdfimages -list archivo.pdf ; pdfimages -all archivo.pdf img
	pdftohtml -s -noframes archivo.pdf archivo.html
	pdftops archivo.pdf archivo.ps
	pdfseparate archivo.pdf pagina-%d.pdf
	pdfunite a.pdf b.pdf unido.pdf          # el último es la SALIDA
\end{lstlisting}
	
	\necesitaespacio{8\baselineskip}
	\subsection*{img2pdf}
	\begin{lstlisting}
	img2pdf *.png -o documento.pdf
	img2pdf --pagesize A4 --auto-orient *.jpg -o documento.pdf
	img2pdf --pagesize A4 --border 1cm --rotation=ifvalid IMG_*.jpg -o movil.pdf
	find . -maxdepth 1 -name '*.png' -print0 | sort -zV | img2pdf --from-file - -o natural.pdf
\end{lstlisting}
	
	\necesitaespacio{8\baselineskip}
	\subsection*{ImageMagick}
	\begin{lstlisting}
	magick imagen.png imagen.jpg
	magick imagen.png -resize 50% salida.png
	magick foto.jpg -resize "1920x1080>" -quality 82 -strip web.jpg
	magick foto.jpg -crop 800x600+100+50 +repage recorte.jpg
	magick foto.jpg -colorspace Gray gris.jpg
	magick a.jpg b.jpg +append juntas.jpg
	magick -density 300 "doc.pdf[0]" -background white -alpha remove -alpha off pag1.png
	mkdir -p web && magick mogrify -path web -resize 50% *.jpg
	magick identify -format "%f %wx%h\n" *.jpg
\end{lstlisting}
	
	\necesitaespacio{8\baselineskip}
	\subsection*{Potrace}
	\begin{lstlisting}
	magick logo.png -colorspace Gray -threshold 50% logo.pbm
	potrace logo.pbm -s -o logo.svg
	potrace imagen.bmp -s -o salida.svg
	potrace -b pdf logo.pbm -o logo.pdf
	mkbitmap -f 2 -s 2 -t 0.48 dibujo.pgm -o dibujo.pbm
\end{lstlisting}
	
	\necesitaespacio{8\baselineskip}
	\subsection*{Pandoc}
	\begin{lstlisting}
	pandoc documento.md -o documento.docx
	pandoc documento.md -o documento.pdf --toc -N -M lang=es -V geometry:margin=2.5cm
	pandoc documento.md -o documento.pdf --pdf-engine=xelatex -M lang=es
	pandoc documento.md -s -o documento.tex
	pandoc documento.md -s --embed-resources -o documento.html
	pandoc documento.docx -t gfm --extract-media=img -o documento.md
	pandoc pagina.html -t gfm -o pagina.md
\end{lstlisting}
	
	\necesitaespacio{8\baselineskip}
	\subsection*{QPDF}
	\begin{lstlisting}
	qpdf --check archivo.pdf
	qpdf --show-npages archivo.pdf
	qpdf --empty --pages a.pdf b.pdf -- unido.pdf
	qpdf archivo.pdf --pages . 1-3,7 -- extracto.pdf
	qpdf --split-pages archivo.pdf pag-%d.pdf
	qpdf archivo.pdf girado.pdf --rotate=+90:1
	qpdf archivo.pdf reparado.pdf
	qpdf --linearize archivo.pdf web.pdf
	qpdf --encrypt --user-password=U --owner-password=P --bits=256 -- archivo.pdf cifrado.pdf
	qpdf --password=P --decrypt cifrado.pdf descifrado.pdf
\end{lstlisting}
	
	\necesitaespacio{8\baselineskip}
	\subsection*{OCRmyPDF}
	\begin{lstlisting}
	ocrmypdf entrada.pdf salida.pdf
	ocrmypdf -l spa entrada.pdf salida.pdf
	ocrmypdf -l spa+eng --rotate-pages --deskew --clean entrada.pdf salida.pdf
	ocrmypdf -l spa --skip-text mixto.pdf salida.pdf
	ocrmypdf -l spa --redo-ocr ocr_viejo.pdf salida.pdf
	ocrmypdf -l spa --sidecar texto.txt entrada.pdf salida.pdf
	ocrmypdf -l spa -O 3 entrada.pdf ligero.pdf
	tesseract --list-langs
\end{lstlisting}
	
	\necesitaespacio{8\baselineskip}
	\subsection*{Ghostscript}
	\begin{lstlisting}
	gs -sDEVICE=pdfwrite -o salida.pdf entrada.pdf
	gs -sDEVICE=pdfwrite -dPDFSETTINGS=/ebook -o reducido.pdf entrada.pdf
	gs -sDEVICE=png16m -r150 -dTextAlphaBits=4 -dGraphicsAlphaBits=4 -o pag-%03d.png \
	entrada.pdf
	gs -sDEVICE=jpeg -dJPEGQ=90 -r150 -o pag-%03d.jpg entrada.pdf
	gs -sDEVICE=pdfwrite -sPageList=1,3,5-7 -o seleccion.pdf entrada.pdf
	gs -sDEVICE=ps2write -o salida.ps entrada.pdf
	ps2pdf entrada.ps salida.pdf
\end{lstlisting}
	\endgroup
	
	% =====================================================================
	\clearpage
	\section{Documentación oficial de referencia}\label{sec:fuentes}
	% =====================================================================
	
	\begingroup\small
	\necesitaespacio{30\baselineskip}
	\begin{longtable}{L{3.2cm}L{11.86cm}}
		\toprule\rowcolor{tablacab}\cab{Recurso} & \cab{Dirección}\\\midrule
		\endfirsthead
		\toprule\rowcolor{tablacab}\cab{Recurso} & \cab{Dirección}\\\midrule
		\endhead
		\bottomrule
		\endfoot
		Marker & \url{https://github.com/datalab-to/marker}\\*
		Surya (modelo de Marker) & \url{https://github.com/datalab-to/surya}\\
		llama.cpp & \url{https://github.com/ggml-org/llama.cpp}\\
		Poppler & \url{https://poppler.freedesktop.org/}\\
		img2pdf & \url{https://gitlab.mister-muffin.de/josch/img2pdf}\\
		ImageMagick & \url{https://imagemagick.org/script/command-line-options.php}\\
		ImageMagick (seguridad) & \url{https://imagemagick.org/script/security-policy.php}\\
		Potrace & \url{https://potrace.sourceforge.net/}\\
		Pandoc & \url{https://pandoc.org/MANUAL.html}\\
		QPDF & \url{https://qpdf.readthedocs.io/}\\
		OCRmyPDF & \url{https://ocrmypdf.readthedocs.io/}\\
		Tesseract & \url{https://tesseract-ocr.github.io/tessdoc/}\\
		Ghostscript & \url{https://ghostscript.readthedocs.io/}\\
		WSL (Microsoft) & \url{https://learn.microsoft.com/windows/wsl/}\\
		Ubuntu en WSL & \url{https://documentation.ubuntu.com/wsl/}\\
		Paquetes de Ubuntu & \url{https://packages.ubuntu.com/}\\*
		Páginas \texttt{man} & \url{https://manpages.ubuntu.com/}\\
	\end{longtable}
	\endgroup
	
	\begin{nota}
		La documentación en línea suele describir la última versión publicada de cada programa, que puede ser posterior a la de Ubuntu 26.04. Ante cualquier duda, la referencia definitiva es la ayuda del programa instalado (\opt{-{}-help}, \cmd{man}).
	\end{nota}
	
\end{document}